\chapter{The Deligne series II: characters, the Arakawa--Kawasetsu MLDE and \texorpdfstring{$q$}{q}-series identities}\label{ch:L25}
\begingroup\setlength{\emergencystretch}{3em}\sloppy

\providecommand{\ltwentyfiveg}{\mathfrak g}
\providecommand{\ltwentyfivevac}{|0\rangle}
\providecommand{\ltwentyfiveD}{\mathcal D}
\providecommand{\ltwentyfivehv}{h^{\vee}}
\providecommand{\ltwentyfiveEu}{\delta_q}

This chapter was written against an earlier state of the LeanVOA library. For this edition, \texttt{scripts/textbook/sync.py} checks that every Lean quotation below occurs verbatim (up to whitespace and comments) in the source files of that release. The machine
record (builds and \texttt{\#print axioms}) is older than that release; Section~\ref{l25:sec:machine} states exactly
what it covers. The list of files covered by the chapter is in Section~\ref{l25:sec:files}.

%======================================================================
\section{Overview: the chain, its status, and how to read this chapter}\label{l25:sec:overview}
%======================================================================

\subsection{What this chapter is about}

The \emph{Deligne exceptional series} is
$A_1\subset A_2\subset G_2\subset D_4\subset F_4\subset E_6\subset E_7\subset E_8$. For its seven higher-rank
members, the level $k=-1-\ltwentyfivehv/6$ gives a simple affine vertex algebra $V=L_k(\ltwentyfiveg)$ of central
charge $c=-2(\ltwentyfivehv+1)$. Physicists know these as the chiral algebras of the rank-one four-dimensional
$\mathcal N=2$ theories (for $A_2$, $D_4$, $E_{6,7,8}$). The project's main theorem (Chapter~\ref{ch:L27}) is
\emph{classical freeness}. Its statement uses three objects, glossed in Section~\ref{l25:sec:glossary}: the
$C_2$-algebra $R_V=V/C_2(V)$, the arc space $J_\infty$ of a scheme, and an explicit space $I_2^{\rm ns}$ of
quadrics (the $C_2$-symbols of the weight-two singular vectors). It says that the graded algebra of $V$ for Li's
filtration is the coordinate ring of the arc space of $\operatorname{Spec}\operatorname{Sym}(\ltwentyfiveg)/(I_2^{\rm ns})$.
For $A_2$ this is now a Lean theorem with no hypotheses (\lean{VOA.DeligneAudit.A2.classicallyFree},
Chapter~\ref{ch:L27}; aggregate build green on 2026-10-28, standard axioms); for the other six types it is proved conditionally on named inputs.

For counting purposes, classical freeness comes down to one equality of dimensions for each conformal weight $N$.
On one side is $\dim V_N$, the number of quantum states of weight $N$. On the other is the number of classical
states, $\sum_p\dim\mathrm{Carrier}_{p,N}$. Here the \emph{carrier} is a combinatorial model of the arc space
(Chapter~\ref{ch:L24}), bigraded by Li degree $p$ and weight $N$. One inequality,
$\dim V_N\le\sum_p\dim\mathrm{Carrier}_{p,N}$, comes from the presentation (Chapter~\ref{ch:L24}). The other,
$\sum_p\dim\mathrm{Carrier}_{p,N}\le\dim V_N$, is a \emph{lower bound on $\dim V_N$}, and this chapter supplies the
part of it that is about $V$. That part has three layers.
\begin{enumerate}
\item \textbf{The VOA side} (Sections~\ref{l25:sec:recurrence}--\ref{l25:sec:alltypes}). For every type,
  $\dim V_N\ge c_N$, where $c_N$ are the coefficients of the unique normalised solution of a second-order
  \emph{modular linear differential equation} (MLDE; primer in Section~\ref{l25:sec:recurrence}). The proof is
  formal. One relation among states of weight four, the \emph{vacuum quartic relation}, forces the MLDE, and the
  indicial equation then forces the equation's constant. These are the first MLDEs in the library proved for specific
  vertex algebras. They are identities of \emph{formal} $q$-series (equivalently, recurrences for the coefficients
  $\dim V_N$): nothing analytic is formalised, neither convergence nor modular invariance nor a second solution. The general Arakawa--Kawasetsu theorem of Chapter~\ref{ch:L16}
  (\lean{IsSqPrimary.exists_monicMLDEOn_all_analyticBlocks_of_etingofFinite_uncond}) is conditional only on
  Etingof-finiteness (\lean{EtingofFinite}, which holds for every $C_2$-cofinite VOA by \lean{EtingofFinite_of_C2Cofinite});
  its side inputs \lean{AKWeightGrading} and \lean{E46SerreClosed} are now theorems (\lean{VOA.akWeightGrading},
  \lean{VOA.e46SerreClosed}). It gives, for abstract analytic genus-one blocks, a monic MLDE of \emph{some} order with
  unspecified coefficients; it does not give the second-order equations proved here. The route here needs neither that
  theorem nor modularity (warning in Section~\ref{l25:sec:block}).
\item \textbf{The $q$-series side} (Sections~\ref{l25:sec:qseries}--\ref{l25:sec:closed}). A library of formal
  $q$-series identities, proved from scratch: the Jacobi triple product, Euler's pentagonal theorem, the
  sum side of the quintuple product, Legendre's four-triangular-numbers theorem, Ramanujan's three differential identities, and
  Lorenz's and Borwein's level-three identities. These give closed forms of $c_N$ for $A_2$, $G_2$, $D_4$, $E_6$,
  $E_7$, $E_8$.
\item \textbf{The Demazure side} (Sections~\ref{l25:sec:demazure}--\ref{l25:sec:star}). Level-one affine Demazure
  characters, computed by an algorithm. The \emph{star identity} ``Demazure sum $=c_N$'' is proved for $A_2$,
  proved for $D_4$ on the unclamped route (Chapter~\ref{ch:L27}), and is a named input for the other five types,
  where it is a conjecture (Conjecture~\ref{l27:conj:star}).
\end{enumerate}

\subsection{The chain of inequalities and its status}\label{l25:sec:chain}

Write $\mathcal C_N=\sum_{p=0}^N\dim\mathrm{Carrier}_{p,N}$ (the classical count), $D_N$ for the Demazure sum
(Definition~\ref{l25:def:dsum}), and $c_N=$\lean{DeligneCharacter.mldeCoeff h N} (Definition~\ref{l25:def:mlde}).
The whole argument, at each weight $N$, is the chain
\begin{equation}\label{l25:eq:chain}
\mathcal C_N\ \overset{\text{(a)}}{\le}\ D_N\ \overset{\text{(b)}}{\le}\ c_N\ \overset{\text{(c)}}{\le}\ \dim V_N
\ \overset{\text{(d)}}{\le}\ \mathcal C_N .
\end{equation}
The chain starts and ends at $\mathcal C_N$. Once every link is proved, every link is an equality, and
$\mathcal C_N=\dim V_N$ is the counting content of classical freeness. The links are as follows.
\begin{itemize}
\item[(a)] \emph{The carrier side}, $\mathcal C_N\le D_N$: Chapters~\ref{ch:L24} and~\ref{ch:L26}, conditional on
  the fusion inputs U5 and U6 (Section~\ref{l25:sec:glossary}); for $A_2$ these inputs are proved, so (a) holds
  unconditionally there (Chapter~\ref{ch:L27}). Not in this chapter.
\item[(b)] \emph{The star bound}, $D_N\le c_N$: Section~\ref{l25:sec:star}. This is the one-sided form of the star
  identity $D_N=c_N$.
\item[(c)] \emph{The target lower bound}, $c_N\le\dim V_N$: Sections~\ref{l25:sec:recurrence}--\ref{l25:sec:alltypes}.
  It is proved for all seven types, with no hypothesis.
\item[(d)] \emph{The upper bound}, $\dim V_N\le\mathcal C_N$: current-jet surjectivity plus the geometric carrier
  embedding (Chapter~\ref{ch:L24}; Theorem~\ref{l25:thm:ofle} below).
\end{itemize}
Links (b) and (c) together are the \emph{Demazure bound} $D_N\le\dim V_N$ (Definition~\ref{l25:def:inputs}), which
is what Chapter~\ref{ch:L27} consumes. The table below records, for each type and each claim, what is a Lean
theorem. The entries mean the following. \emph{Theorem}: a Lean theorem with no hypotheses. \emph{Derivable}: it
follows by composing exported Lean theorems, but no single Lean statement says it. \emph{Conditional}: a Lean
theorem with the named hypothesis. \emph{Open}: a named \lean{Prop} that is not proved.
\begin{center}\small
\begin{tabular}{p{0.06\textwidth}p{0.12\textwidth}p{0.15\textwidth}p{0.28\textwidth}p{0.23\textwidth}}
\toprule
type & (c) $\dim V_N\ge c_N$ & $\dim V_N=c_N$ & (b) star: $D_N\le c_N$ or $=$ & closed form of $c_N$\\
\midrule
$A_2$ & theorem & theorem & theorem ($=$; Thm.~\ref{l25:thm:a2star}) & theorem: $\prod_{r\ge1}(1+q^r)^8$\\
$G_2$ & theorem & derivable$^\dagger$ & \textbf{open} (\lean{G2StarBound}; Conj.~\ref{l27:conj:star}) & theorem: $\theta_{A_2}(q)\frac{(q^3;q^3)_\infty^6}{(q;q)_\infty^8}$\\
$D_4$ & theorem & derivable$^\dagger$ & $\mathtt{d4SchurQ}=c_N$ theorem; $D_N=\mathtt{d4SchurQ}$ conditional on \lean{D4DemazurePositivity} & theorem: $\frac{\delta_qE_4}{240\,(q)_\infty^{10}}$ (shifted)\\
$F_4$ & theorem & derivable$^\dagger$ & \textbf{open} (\lean{F4StarBound}; Conj.~\ref{l27:conj:star}) & none in the library\\
$E_6$ & theorem & derivable$^\ddagger$ & \textbf{open} (\lean{E6StarBound}; Conj.~\ref{l27:conj:star}) & theorem (Thm.~\ref{l25:thm:closed})\\
$E_7$ & theorem & derivable$^\ddagger$ & \textbf{open} (\lean{E7StarBound}; Conj.~\ref{l27:conj:star}) & theorem (Thm.~\ref{l25:thm:closed})\\
$E_8$ & theorem & derivable$^\ddagger$ & \textbf{open} (\lean{E8StarBound}; Conj.~\ref{l27:conj:star}) & theorem (Thm.~\ref{l25:thm:closed})\\
\bottomrule
\end{tabular}
\end{center}
\noindent $^\dagger$The composition was typechecked in a scratch file for this chapter. $^\ddagger$The composition
was not typechecked, because the $E_{6,7,8}$ modules are too heavy to import for a single check. For both marks see
Caveat~\ref{l25:cav:equality}. The $A_2$ product entry has a further gap: the equality
$\dim V_N=[q^N]\prod(1+q^r)^8$ is itself only derivable (Remark~\ref{l25:rem:a2cons}). Every ``theorem'' of the
star column is a statement about the computed Demazure tables. Reading those tables as characters of Demazure
modules is the input U7 (Caveat~\ref{l25:cav:U7}).

\paragraph{What Chapter~\ref{ch:L27} still needs from this chapter.} Its final theorems for $G_2$, $F_4$, $E_6$,
$E_7$, $E_8$ take the star bound \lean{XStarBound} as an explicit hypothesis, the open entry of column (b). The
$A_2$ final theorem takes no hypothesis at all: (b) and (c) are both proved here, and the carrier-side inputs are proved
in Chapters~\ref{ch:L26} and~\ref{ch:L27}. The $D_4$ final theorem
takes no star hypothesis either, because it works with the unclamped Demazure series (see
Chapter~\ref{ch:L27}). Link (c) is used in full for all seven types.

\subsection{Which sections are on the critical path}

The \textbf{critical path}, needed for links (b) and (c) in Chapter~\ref{ch:L27}, consists of these sections:
Section~\ref{l25:sec:interface} (the interface), Sections~\ref{l25:sec:recurrence}--\ref{l25:sec:alltypes} (link (c)),
Section~\ref{l25:sec:demazure} (the Demazure tables), Section~\ref{l25:sec:star} (link (b)), and those $q$-series
results that feed the $A_2$ and $D_4$ star identities: the triple product, Legendre, Ramanujan's identities, the
$A_2$ product, and the $D_4$ closed form.

The following are \textbf{side results}, off the critical path at this commit. They are marked as such where they
occur.
\begin{itemize}
\item The $G_2$ target side (Lorenz, Borwein, Theorem~\ref{l25:thm:g2target}). It becomes critical only if the open
  \lean{G2DemazureEqKW} is proved.
\item Euler's pentagonal theorem and the sum side of the quintuple product (Theorem~\ref{l25:thm:pent}), except as ingredients of
  Lorenz's formula.
\item The $E_6$, $E_7$, $E_8$ closed forms (Theorem~\ref{l25:thm:closed}). The star bounds target $c_N$ directly.
\item The $q=1$ counts and the energy bounds (Section~\ref{l25:sec:counts}). They serve Chapter~\ref{ch:L26} and the
  \lean{energy} field of the $D_4$ positivity input.
\end{itemize}

\subsection{Glossary of terms and named inputs}\label{l25:sec:glossary}

The terms below are used in this chapter but defined elsewhere. Each gets a one-line gloss here and a pointer to
where it is developed.
\begin{itemize}
\item \textbf{$C_2$-algebra} $R_V=V/C_2(V)$, where $C_2(V)$ is spanned by the states $a_{(-2)}b$. It is a
  commutative Poisson algebra, the ``classical limit'' of $V$ (Chapter~\ref{ch:L12}). Its spectrum is the
  \textbf{associated scheme}; the reduced variety underneath is the \textbf{associated variety}.
\item \textbf{Li filtration}: a decreasing filtration of $V$ whose degree-$0$ graded piece is $R_V$. Its associated
  graded $\mathrm{gr}V$ is a quotient of the \textbf{arc space} (jet ring) $J_\infty R_V$, the ring of functions on
  formal curves $t\mapsto x(t)$ in $\operatorname{Spec}R_V$ (Chapters~\ref{ch:L11} and~\ref{ch:L24}).
\item \textbf{Quasi-lisse}: the associated variety has finitely many symplectic leaves. Arakawa and Kawasetsu prove
  that the characters of such VOAs satisfy an MLDE (the general theorem is in Chapter~\ref{ch:L16}).
\item \textbf{Higgs branch}: in the 4d $\mathcal N=2$ theory, the moduli space on which the associated variety of the
  chiral algebra is realised. For the Deligne series it is the minimal nilpotent orbit closure.
\item \textbf{Joseph relations, $I_2$, $I_2^{\rm ns}$.} $\operatorname{Sym}^2\ltwentyfiveg=V(2\theta)\oplus
  I_2^{\rm ns}\oplus\C Q$, with $Q$ the Casimir quadric. For $\ltwentyfiveg$ not of type $A$, $I_2=I_2^{\rm ns}\oplus\C Q$
  generates the associated graded of the Joseph ideal (literature, not Lean); for $A_2$ there is no Joseph ideal, and
  in Lean $I_2^{\rm ns}$ is simply the ideal spanned by explicit lists of quadrics (Chapter~\ref{ch:L24}). The superscript ``ns'' means \emph{non-singlet}: $I_2^{\rm ns}$ is the part of $I_2$ without
  the $\ltwentyfiveg$-invariant $Q$. These are the $C_2$-symbols of the weight-two singular vectors (Chapter~\ref{ch:L24},
  ``The Deligne series'').
\item \textbf{U3}: $X_a\,Q\in(I_2^{\rm ns})$ in $\operatorname{Sym}^3\ltwentyfiveg$ for every current $X_a$, so $Q^2$ lies in
  the ideal too. It is certified per type, as the field \lean{u3} of \lean{DeligneCarrierData} (Chapter~\ref{ch:L24}).
\item \textbf{U5}: the fusion bound (Input~T of Chapter~\ref{ch:L27}). The fusion module $\mathcal F_m$ has dimension
  $\le(\dim\ltwentyfiveg)^m$, with graded character $F_m$. Its published source for the total dimension is
  Kus--Littelmann, Pacific J.\ Math.\ 278, Theorem~4.2, with one flagged step (p.~437); Feigin's earlier proof has a
  gap (Chapter~\ref{ch:L27}). It is a named hypothesis, reduced in Chapter~\ref{ch:L26}.
\item \textbf{U6}: Ravinder's mixed fusion recursion, which expresses $\mathrm{Hilb}_qD(1,n\theta)$ through the
  $F_m$. Named; Chapter~\ref{ch:L26}.
\item \textbf{U7}: the identification of a \emph{computed} table with the graded character of the level-one Demazure
  module $D(1,n\theta)$. It rests on the Demazure character formula (Kumar, Mathieu) and on the identification of
  $D(1,n\theta)$ with a fusion product (Fourier--Littelmann, Chari--Venkatesh); for $A_2$, on Sanderson's and Ion's
  $q$-Whittaker formula. It is not formalised; in Chapters~\ref{ch:L26}--\ref{ch:L27} it appears as the hypothesis
  \lean{DemazureRealizationG} (Input~R of Chapter~\ref{ch:L27}). The star bound $D_N\le c_N$ is Input~S there.
\item \textbf{Demazure module} $D(1,n\theta)$: the subspace of the level-one integrable module $L(\Lambda_0)$ of
  $\widehat{\ltwentyfiveg}$ generated by an extremal vector of weight $n\theta$ under the positive part. Its
  weight-$0$ part is the irreducible $V(n\theta)$, and $D(1,\theta)=\ltwentyfiveg\oplus\C$.
\item \textbf{Kirillov--Reshetikhin module} $\mathrm{KR}(\theta)$: $\ltwentyfiveg\oplus\C$ as a module for the current
  algebra $\ltwentyfiveg[t]$. Fusion products of copies of it give the $D(1,n\theta)$.
\item \textbf{$q$-Whittaker polynomial} $P_\lambda(x;q,0)$: the Macdonald polynomial at $t=0$. \textbf{Gelfand--Tsetlin
  branching}: computing it by restricting $GL_3\to GL_2\to GL_1$ along interlacing patterns. At $t=0$ each step
  contributes a product of $q$-binomials (dictionary of Theorem~\ref{l25:thm:a2star}).
\item \textbf{Rogers--Szeg\H{o} polynomial} $H_M(x)=\sum_{r=0}^M\binom Mr_qx^r$. We write $G_M=H_M(1)$.
\item \textbf{HKOTY fermionic formula}: a sum over quasi-particle configurations, weighted by $q$-binomials. For
  graded fusion products of KR modules it is proved to give the character, for every $\ltwentyfiveg$ (Ardonne--Kedem,
  Di Francesco--Kedem); the $X=M$ form, with one-dimensional sums, is a separate statement.
\item \textbf{Plethystic (product) exponents} of a series $1+\sum_{n\ge1}b_nq^n$: the integers $a_n$ with
  $1+\sum b_nq^n=\prod_n(1-q^n)^{-a_n}$.
\end{itemize}

\subsection{Notation}\label{l25:sec:notation}

Several letters have more than one job in the sources. We fix them as follows; Lean names keep their Lean spelling.
\begin{center}\small
\begin{tabular}{p{0.12\textwidth}p{0.62\textwidth}p{0.18\textwidth}}
\toprule
symbol & meaning in this chapter & where\\
\midrule
$\delta_q$ & the Euler operator $q\,\frac{d}{dq}$ on $q$-series (Lean \lean{qDeriv}, \lean{eulerOp}) & \S\ref{l25:sec:recurrence}\\
$\theta$ & the highest root of $\ltwentyfiveg$ (only); $\theta^\vee$ its coroot & \S\ref{l25:sec:demazure}\\
$\theta_{A_2}$, $\theta_2$, $\Theta(x)$ & the $A_2$ lattice theta series; the Jacobi theta constant; the $D_4$ theta function in $x$ & Thms.~\ref{l25:thm:g2target}, \ref{l25:thm:legendre}, \ref{l25:thm:d4}\\
$\Theta$ (with $u$) & the triple-product sum $\sum(-1)^ju^jq^{j(j-1)/2}$ in the heat-equation proof & Thm.~\ref{l25:thm:ramanujan}\\
$\kappa$ & the constant of the vacuum quartic relation (Lean \lean{κ} in \lean{VacuumQuarticRelation W κ}); the bilinear form in \lean{LieData κ …} is written $\kappa_{\ltwentyfiveg}$ & Def.~\ref{l25:def:vqr}\\
$h$ & the dual Coxeter number $\ltwentyfivehv$ when it parametrises the MLDE; \lean{D.h} is the same number read off a \lean{LieData} & \S\ref{l25:sec:recurrence}\\
$\mathrm{hilb}(n,N)$ & a table of Demazure coefficients $[q^N]\mathrm{Hilb}_qD(1,n\theta)$ & Def.~\ref{l25:def:dsum}\\
$c$ & the central charge; $c_N$ the AK coefficients; $c_0$ (Lean \lean{c₀}) the Casimir ratio of Thm.~\ref{l25:thm:carrierdata} & \\
$z(k,m,m{+}p)$ & Zhu's change-of-coordinate coefficients & Thm.~\ref{l25:thm:sqcomm}\\
$k$ & the level $-1-h/6$; elsewhere a modular weight (said explicitly) & \\
$(a;q)_n$, $(q)_n$ & $\prod_{i=0}^{n-1}(1-aq^i)$; $(q)_n=(q;q)_n$; $n=\infty$ allowed. $(q)_\infty=(q;q)_\infty$ is the Euler product (Lean \lean{euler}) & \S\ref{l25:sec:qseries}\\
$p$, $r$ & $p$: Li degree; in the quintuple product, a nome. $r=q^{1/2}$ in the $D_4$ section; elsewhere $r$ is a summation index & \\
$P_\lambda$, $P_1,P_3$, $P$ & the $q$-Whittaker polynomial; the divisor series $\sum\sigma_1(n)q^n$, $\sum\sigma_3(n)q^n$ (Lean \lean{P1}, \lean{P3}); in the proof of Thm.~\ref{l25:thm:ramanujan} only, $P$ is the triple-product side & \\
$D_N$; $D_i$ & the Demazure sum; Demazure operators (with subscript) & \\
$W$ & the VOA in a Lean statement (Lean \lean{W}); the Lambert series of the $D_4$ proof is $\mathcal W_1$ & \\
$S_M(v)$ & the torus one-point function $\tr_Mo(v)q^{L_0-c/24}$ & \S\ref{l25:sec:block}\\
\bottomrule
\end{tabular}
\end{center}

\subsection{Reading the Lean in this chapter}\label{l25:sec:readlean}

Chapter~\ref{ch:01}, Section~\ref{ch01:sec:lean}, is the general primer. The idioms that the printed code of this
chapter needs are these.
\begin{itemize}
\item \textbf{Truncated blocks.} Proof bodies are omitted. A printed block therefore ends at \lean{:=},
  \lean{:= by} or \lean{where}, and the proof follows in the file. \lean{theorem foo : S where} with nothing after
  it is such a block: the omitted body fills the fields of the structure \lean{S} one by one. It is not a syntax
  error.
\item \textbf{Binders.} \lean{(x : T)} is an explicit argument, \lean{{x : T}} an implicit one inferred by Lean, and
  \lean{[Field K]} an instance argument (a structure found automatically). \lean{(K := ℚ)} supplies the named
  implicit argument \lean{K}. \lean{section} variables are arguments shared by all declarations of a section; we
  list them in the dictionaries.
\item \textbf{\lean{Prop}-valued structures.} \lean{structure S : Prop where eq : …} is a bundled hypothesis. A
  proof \lean{h : S} gives \lean{h.eq}. A theorem that takes \lean{(h : S)} is conditional on \lean{S}, and
  \#\lean{print axioms} does not show this.
\item \textbf{\lean{abbrev} vs \lean{def}}: both define, and \lean{abbrev} unfolds automatically.
\item \textbf{Finite sums.} \lean{∑ i ∈ s, f i} is a sum over the finite set \lean{s}. \lean{Finset.range n}
  $=\{0,\dots,n-1\}$. \lean{Finset.antidiagonal n} $=\{(i,j):i+j=n\}$. \lean{Finset.Icc a b} $=[a,b]$, and
  \lean{Finset.Icc a b with e n = d} is its subset where $e(n)=d$.
\item \textbf{Small notation.} \lean{•} is scalar multiplication. \lean{x⁻¹} is the inverse. \lean{l.getD i 0} is
  entry $i$ of a list (default $0$). \lean{f^[m]} is the $m$-fold iterate of \lean{f}. \lean{x.1.2} is the second
  component of the first component of a pair.
\item \textbf{Algebra.} \lean{(n : ℚ)} is a cast. \lean{MvPolynomial ι ℂ} is $\C[X_i:i\in\iota]$.
  \lean{Submodule.span ℂ s} is the span. \lean{V ⧸ U} is the quotient space and \lean{U.mkQ} the quotient map.
  \lean{PowerSeries R} is $R[[q]]$ (Chapter~\ref{ch:01}).
\item \textbf{Kernel checks.} \lean{decide +kernel} proves a decidable statement by evaluating it in Lean's trusted
  kernel. \lean{native_decide}, which would trust compiled code, is never used. \lean{d4_precompute}\texttt{\%} is a custom
  elaborator that only \emph{proposes} literals; the kernel re-checks them.
\end{itemize}

\subsection{What was machine-checked, and when}\label{l25:sec:machine}

The quotations were re-verified against the state of the library this chapter was written against. The last complete build of the umbrella module
\lean{VOA.Deligne} is older: it was built green at an earlier commit. At that commit it imported, directly or transitively, all $239$ files of this chapter, including the $E_6$, $E_7$, $E_8$ modules
(we checked the import closure). The build took 11099 jobs and reported no
\lean{sorryAx}. The \#\lean{print axioms} commands in those files all reported a subset of
\lean{[propext, Classical.choice, Quot.sound]}. The \emph{aggregate} \lean{lake build VOA} was last run green on
26~October (13868 jobs), at an earlier commit. Neither has been re-run at the state this chapter was written against. Between
that umbrella build and that state the Lean sources changed (the diff of \texttt{VOA} touches
1390 files): certificate files were split to fit the memory budget (among them the $D_4$ divisor and the $A_2$ divisor
certificates), a few files of this chapter were edited (for example \leanpath{VOA/DeligneAudit/DeligneMLDE.lean}, a
docstring-only change), and the unconditional $A_2$ files were added. According to the project's internal build record of
27 October 2026 (not re-run for this chapter) the exported statements of the split files are unchanged and the repaired targets were built green in
targeted builds with standard axioms; every statement this chapter quotes occurs verbatim in the state this chapter was written against. On 2026-10-28 the aggregate build \texttt{bash scripts/campaign\_build.sh VOA} (the full \texttt{lake build VOA} under the project's memory-capped wrapper) was green on 2026-10-28 (a pre-release build): 18{,}901 jobs, and again 18{,}907 jobs after merging the then-current \texttt{origin/main}; the aggregate target contains the umbrella. The ``standard axioms'' of the
status boxes are campaign records of the earlier umbrella build and of the later targeted builds. For this chapter we re-ran
\#\lean{print axioms} on the generic lemmas named in the status boxes, plus the scratch checks reported in
Caveat~\ref{l25:cav:equality} and Example~\ref{l25:ex:a2}. We ran no build.

\textbf{Axioms.} Unless a status box says otherwise, every theorem of this chapter depends only on Lean's three
standard axioms \lean{propext}, \lean{Classical.choice} and \lean{Quot.sound}. The library has no project axiom
(\lean{VOA.pbw} is a theorem, Chapter~\ref{ch:L02}), and nothing here uses it. The status boxes therefore say
``standard axioms''. They separately say whether we re-ran the check for this chapter.

\subsection{Recurring caveats, stated once}\label{l25:sec:caveats}

\begin{enumerate}
\item\label{l25:cav:equality} \textbf{Equality is a Lean theorem only for $A_2$.} $\dim V_N=c_N$ is stated in Lean for
  $A_2$ only (\lean{A2.finrank_eq_mldeCoeff'}). For the other six types only $\ge$ is stated. The equality is
  derivable (warning after Theorem~\ref{l25:thm:alltypes}).
\item\label{l25:cav:U7} \textbf{Demazure tables are algorithms.} \lean{levelOneDemazureHilb} and the $A_2$ formula
  \lean{demazureHilb} are computable tables. Reading them as characters of $D(1,n\theta)$ is input U7, which is not
  formalised.
\item\label{l25:cav:hyp} \textbf{Named inputs are hypotheses, never axioms.} A star identity or bound is a
  \lean{Prop}. A theorem that uses one takes it as an argument, and \#\lean{print axioms} does not show it.
\item\label{l25:cav:docs} \textbf{Stale documentation.} Several docstrings call proved statements ``named'' or
  ``literature'' inputs. They are listed once, in Section~\ref{l25:sec:defects}, and not repeated next to the
  theorems.
\end{enumerate}

%======================================================================
\section{The character side of classical freeness}\label{l25:sec:interface}
%======================================================================

\subsection{Where the characters enter}

Let $V=L_k(\ltwentyfiveg)$ and let $R_V=V/C_2(V)$ be Zhu's $C_2$-algebra (Chapter~\ref{ch:L12}). Classical
freeness asserts that the Li filtration of $V$ has associated graded equal to the coordinate ring of the arc
space $J_\infty X$ of the scheme $X=\operatorname{Spec}R_V$, and that $R_V=\operatorname{Sym}(\ltwentyfiveg)/(I_2^{\rm
ns})$ for an explicit ideal $I_2^{\rm ns}$ of quadrics. Chapter~\ref{ch:L24} reduces this to three
\emph{packages} for a presentation $P$ of $V$ and a ``carrier'' family of vector spaces
$\mathrm{Carrier}_{p,N}$ (bigraded by Li degree $p$ and conformal weight $N$):
current-jet surjectivity, an injective geometric carrier embedding, and the \emph{ordinary carrier character}
identity $\sum_{p\le N}\dim\mathrm{Carrier}_{p,N}=\dim V_N$. The first two give one inequality
for free (link (d) of \eqref{l25:eq:chain}); the third is the character side. Its lower half, links (b) and (c), is
the subject of this chapter.

\begin{definition}[Carrier character lower bound]\label{l25:def:cclb}
For a VOA $V$ and a bigraded family $\mathrm{Carrier}_{p,N}$, the \emph{carrier character lower bound} is the
statement
\begin{equation*}
\forall N\ge0:\qquad \sum_{p=0}^{N}\dim_\C\mathrm{Carrier}_{p,N}\;\le\;\dim_\C V_N .
\end{equation*}
\end{definition}
\noindent File \leanpath{VOA/\allowbreak DeligneAudit/\allowbreak A2Character.lean}.
\begin{leancode}
def CarrierCharacterLowerBound (W : VertexOperatorAlgebra V)
    (Carrier : ℕ → ℕ → Type uC)
    [∀ p N, AddCommGroup (Carrier p N)]
    [∀ p N, Module ℂ (Carrier p N)] : Prop :=
  ∀ N,
    (∑ p ∈ Finset.range (N + 1), Module.finrank ℂ (Carrier p N)) ≤
      Module.finrank ℂ (W.grading N)
\end{leancode}
\begin{dictionary}
\lean{W.grading N} is the weight-$N$ subspace $V_N$ (a \lean{Submodule}; the natural number $N$ is coerced to
$\Z$); \lean{Module.finrank ℂ} is its complex dimension (Mathlib's \lean{finrank}, which is $0$ for an
infinite-dimensional space --- harmless here, since every $V_N$ is finite-dimensional for the VOAs of this
chapter). \lean{Finset.range (N + 1)} is $\{0,\dots,N\}$. The definition is exactly
Definition~\ref{l25:def:cclb}. It is a \lean{Prop}-valued \lean{def}, not a structure, so it is used by
applying it to $N$.
\end{dictionary}
\begin{physics}
$\dim V_N$ counts the states of conformal weight $N$ in the vacuum module of the chiral CFT: the vacuum
character $\tr_V q^{L_0}$. The carrier is a model of the classical (``$\hbar\to0$'') phase space of the
theory, the arc space of the Higgs branch. The inequality says the quantum theory has at least as many
states as its classical limit predicts --- the nontrivial direction, since quantum corrections usually
\emph{remove} states (null vectors).
\end{physics}
\begin{status}
Definition. Unconditional by construction.
\end{status}

\begin{theorem}[The lower bound suffices]\label{l25:thm:ofle}
Given a presentation $P$, current-jet surjectivity and an injective geometric carrier embedding, the carrier
character lower bound implies the ordinary carrier character identity.
\end{theorem}
\noindent File \leanpath{VOA/\allowbreak DeligneAudit/\allowbreak A2Character.lean}.
\begin{leancode}
theorem ordinaryCarrierCharacter_of_le
    (P : CurrentC2Presentation case W)
    (Carrier : ℕ → ℕ → Type uC)
    [∀ p N, AddCommGroup (Carrier p N)]
    [∀ p N, Module ℂ (Carrier p N)]
    [∀ p N, FiniteDimensional ℂ (Carrier p N)]
    (hsurj : CurrentJetSurjectivity case W P)
    (hgeom : GeometricCarrierEmbedding case W P Carrier)
    (hlow : CarrierCharacterLowerBound W Carrier) :
    OrdinaryCarrierCharacter W Carrier
\end{leancode}
\begin{dictionary}
The packages \lean{CurrentC2Presentation}, \lean{CurrentJetSurjectivity}, \lean{GeometricCarrierEmbedding},
\lean{OrdinaryCarrierCharacter} are defined in Chapter~\ref{ch:L24}. \lean{OrdinaryCarrierCharacter W Carrier}
is the \emph{equality} $\sum_p\dim\mathrm{Carrier}_{p,N}=\dim V_N$ for all $N$. The theorem is the
statement for any \lean{case} and any VOA, so the same reduction serves all seven types. A companion,
\lean{carrierCharacterLowerBound_of_ordinary}, gives the converse, so in the presence of the two other
packages the lower bound and the identity are \emph{equivalent}: nothing is lost. \textbf{Implicit data and an
extra hypothesis.} \lean{case : HigherRankDeligneCase} (the Deligne type), the carrier space \lean{V : Type uV} with
its VOA \lean{W}, and the universe \lean{uC} of the carrier are section variables (implicit). The instance
\lean{[∀ p N, FiniteDimensional ℂ (Carrier p N)]} is a hypothesis not visible in
Theorem~\ref{l25:thm:ofle} as stated: it makes \lean{finrank} the true dimension of each carrier space. Without it
\lean{finrank} would be $0$ on an infinite-dimensional carrier and \lean{hlow} would be weaker than it reads.
\end{dictionary}
\begin{proof}[Proof on paper]
Current-jet surjectivity says each Li layer $\mathrm{gr}^{\rm Li}_{p}V_N$ is a quotient of a finite-dimensional
jet space $J_{p,N}$, and the carrier embedding injects the dual $J_{p,N}^*$ into $\mathrm{Carrier}_{p,N}$.
Hence $\dim\mathrm{gr}^{\rm Li}_pV_N\le\dim\mathrm{Carrier}_{p,N}$; summing over $p$ and using that the Li
layers of $V_N$ add up to $\dim V_N$ gives $\dim V_N\le\sum_p\dim\mathrm{Carrier}_{p,N}$ (the library's
\lean{finrank_grading_le_carrier_sum}). Combined with the hypothesis, equality.
\end{proof}
\begin{physics}
The upper bound is ``every state is generated by currents, modulo the classical relations''; the lower bound
is ``no further quantum relations occur''. Together: the chiral algebra is \emph{classically free}, its
quantum states are in bijection with classical arc-space functions.
\end{physics}
\begin{status}
Standard axioms. Unconditional as an implication.
\end{status}

\subsection{The split at the Demazure sum}

The carrier is too complicated to compare with $\dim V_N$ directly. The project interposes a combinatorial
series, the \emph{Demazure sum}, built from Hilbert series of level-one affine Demazure modules $D(1,n\theta)$:
\begin{equation}\label{l25:eq:dsum}
\sum_{N\ge0}D_Nq^N \;=\;\sum_{n\ge0}\frac{q^n\,\mathrm{Hilb}_q D(1,n\theta)}{(q;q)_n},\qquad
(q;q)_n=\prod_{i=1}^n(1-q^i).
\end{equation}
The carrier side (Chapter~\ref{ch:L24} and Chapter~\ref{ch:L26}) proves $\sum_p\dim\mathrm{Carrier}_{p,N}\le D_N$,
link (a) of \eqref{l25:eq:chain}; this chapter must prove $D_N\le\dim V_N$, links (b) and (c).

\begin{definition}[Coefficient arithmetic; the Demazure sum]\label{l25:def:dsum}
For coefficient tables $f,g:\N\to\N$ let $(f\ast g)(N)=\sum_{i=0}^Nf(i)g(N-i)$ (the Cauchy product), and let
$p_{\le n}(N)$ be the number of partitions of $N$ into parts of size at most $n$, i.e.\ $[q^N]1/(q;q)_n$. For a
table $\mathrm{hilb}(n,N)=[q^N]\mathrm{Hilb}_qD(1,n\theta)$, the Demazure sum is
\begin{equation*}
D_N=\sum_{n=0}^{N}\big(\mathrm{hilb}(n,\cdot)\ast p_{\le n}\big)(N-n),
\end{equation*}
which is the $q^N$ coefficient of \eqref{l25:eq:dsum}.
\end{definition}
\noindent Files \leanpath{VOA/\allowbreak DeligneAudit/\allowbreak A2CharacterInterface.lean} and
\leanpath{VOA/\allowbreak DeligneAudit/\allowbreak DeligneCharacter.lean}.
\begin{leancode}
def conv (f g : ℕ → ℕ) (N : ℕ) : ℕ :=
  ∑ i ∈ Finset.range (N + 1), f i * g (N - i)

def partitionsLe : ℕ → ℕ → ℕ
  | 0, N => if N = 0 then 1 else 0
  | n + 1, N => ∑ j ∈ Finset.range (N / (n + 1) + 1), partitionsLe n (N - j * (n + 1))

def demazureSumCoeff (hilb : ℕ → ℕ → ℕ) (N : ℕ) : ℕ :=
  ∑ n ∈ Finset.range (N + 1), conv (hilb n) (partitionsLe n) (N - n)
\end{leancode}
\begin{dictionary}
\lean{partitionsLe n N}: the recursion conditions on the number $j$ of parts equal to $n+1$, which is
exactly how $1/(q;q)_{n+1}=\frac{1}{1-q^{n+1}}\cdot\frac{1}{(q;q)_n}$ expands. So
\lean{partitionsLe n N} $=[q^N]1/(q;q)_n$. In \lean{demazureSumCoeff}, the shift $N-n$ accounts for the
factor $q^n$ in \eqref{l25:eq:dsum}, and terms with $n>N$ vanish, so the finite range $n\le N$ is exact. All of
this is $\N$-valued and computable, so small cases are checked by \lean{decide +kernel} (kernel evaluation,
not \lean{native_decide}). \textbf{What is \lean{hilb}?} Any table; the definition does not know about
Demazure modules. \emph{All} Demazure content enters through the choice of table, and its reading as
$\mathrm{Hilb}_qD(1,n\theta)$ is the input U7 (Caveat~\ref{l25:cav:U7}). The concrete tables are built in
Section~\ref{l25:sec:demazure}, and the $A_2$ table in Theorem~\ref{l25:thm:a2star}; in Lean a table is only what it
computes.
\end{dictionary}
\begin{physics}
\emph{Heuristic.} In the fusion picture (Chapter~\ref{ch:L26}), $D(1,n\theta)$ is the graded fusion product of $n$
copies of the adjoint-plus-singlet representation, and $1/(q;q)_n$ counts the ways of distributing $n$
indistinguishable excitations among levels. The sum \eqref{l25:eq:dsum} then reads as a ``fermionic''
(quasi-particle) form of a character: states built from $n$ quasi-particles of adjoint type, with momenta encoded by
the Demazure grading. As a statement about the vacuum module this is conjectural; its mathematical content is links
(a) and (b) of \eqref{l25:eq:chain}, i.e.\ U5--U7 and the star identity.
\end{physics}
\begin{status}
Definitions; computable; unconditional.
\end{status}

\begin{definition}[Named inputs of the character side, integer form]\label{l25:def:inputs}
For a table \lean{hilb}, a target sequence $t:\N\to\N$ and a VOA $W$:
\begin{itemize}
\item the \emph{star identity} $D_N=t_N$ for all $N$;
\item the \emph{target lower bound} $t_N\le\dim W_N$ for all $N$;
\item the \emph{Demazure bound} $D_N\le\dim W_N$ for all $N$.
\end{itemize}
\end{definition}
\noindent File \leanpath{VOA/\allowbreak DeligneAudit/\allowbreak DeligneCharacter.lean}.
\begin{leancode}
structure StarIdentityFor (hilb : ℕ → ℕ → ℕ) (target : ℕ → ℕ) : Prop where
  eq : ∀ N, demazureSumCoeff hilb N = target N

structure TargetLowerBound (W : VertexOperatorAlgebra V) (target : ℕ → ℕ) : Prop where
  le : ∀ N : ℕ, target N ≤ Module.finrank ℂ (W.grading N)

structure DemazureBoundFor (W : VertexOperatorAlgebra V) (hilb : ℕ → ℕ → ℕ) : Prop where
  le : ∀ N : ℕ, demazureSumCoeff hilb N ≤ Module.finrank ℂ (W.grading N)

theorem demazureBoundFor_of_inputs {W : VertexOperatorAlgebra V} {hilb : ℕ → ℕ → ℕ}
    {target : ℕ → ℕ} (hstar : StarIdentityFor hilb target) (hlow : TargetLowerBound W target) :
    DemazureBoundFor W hilb
\end{leancode}
\begin{dictionary}
Each is a \lean{Prop}-valued \lean{structure} with one field; a proof of \lean{StarIdentityFor hilb t} is a
proof of $\forall N,\ D_N=t_N$, accessed as \lean{h.eq N}. \textbf{These are hypotheses, never axioms}: a
theorem taking \lean{(hstar : StarIdentityFor …)} is conditional on it, and \#\lean{print axioms} does not
show this (Caveat~\ref{l25:cav:hyp}). \lean{demazureBoundFor_of_inputs} is the transitivity
$D_N=t_N\le\dim W_N$.
\end{dictionary}
\begin{physics}
The ``target'' is the vacuum character of the VOA as predicted by representation theory or by the MLDE; the
star identity says the quasi-particle (fermionic) form of the character equals it. Such ``fermionic $=$
bosonic'' identities are the hallmark of integrable structure in 2d CFT (Rogers--Ramanujan being the
prototype for the Lee--Yang model).
\end{physics}
\begin{status}
Definitions and a one-line transitivity; standard axioms.
\end{status}

\begin{theorem}[The glue]\label{l25:thm:glue}
If $\sum_p\dim\mathrm{Carrier}_{p,N}\le D_N$ for all $N$ and the Demazure bound holds, the carrier character lower
bound holds.
\end{theorem}
\noindent File \leanpath{VOA/\allowbreak DeligneAudit/\allowbreak DeligneCharacter.lean}.
\begin{leancode}
theorem carrierCharacterLowerBound_of_demazureSum {W : VertexOperatorAlgebra V}
    {hilb : ℕ → ℕ → ℕ}
    (Carrier : ℕ → ℕ → Type uC)
    [∀ p N, AddCommGroup (Carrier p N)] [∀ p N, Module ℂ (Carrier p N)]
    (hcar : ∀ N, (∑ p ∈ Finset.range (N + 1), Module.finrank ℂ (Carrier p N)) ≤
      demazureSumCoeff hilb N)
    (hchr : DemazureBoundFor W hilb) :
    CarrierCharacterLowerBound W Carrier
\end{leancode}
\begin{dictionary}
Exactly Theorem~\ref{l25:thm:glue}, a transitivity in $\N$. \lean{hcar} is the carrier side
(Chapter~\ref{ch:L26}); \lean{hchr} is this chapter's goal.
\end{dictionary}
\begin{proof}[Proof on paper]
$\sum_p\dim\mathrm{Carrier}_{p,N}\le D_N\le\dim W_N$.
\end{proof}
\begin{physics}
Classical states $\le$ quasi-particle states $\le$ quantum states. The first inequality is a statement about
the classical phase space, the second about the quantum theory.
\end{physics}
\begin{status}
Standard axioms; unconditional as an implication.
\end{status}

\subsection{The rational form and the one-sided form}

The natural targets are rational sequences (the MLDE has rational coefficients), so the library also has
$\Q$-valued versions.

\begin{definition}[Rational inputs]\label{l25:def:Qinputs}
For $t:\N\to\Q$: \lean{StarIdentityQ hilb t} is $D_N=t_N$, \lean{TargetLowerBoundQ W t} is $t_N\le\dim W_N$
(both in $\Q$), and \lean{StarBoundQ hilb t} is the one-sided $D_N\le t_N$.
\end{definition}
\noindent Files \leanpath{VOA/\allowbreak DeligneAudit/\allowbreak DeligneCharacterTypes.lean} and
\leanpath{VOA/\allowbreak DeligneAudit/\allowbreak DeligneCharacterStarBound.lean}.
\begin{leancode}
structure StarIdentityQ (hilb : ℕ → ℕ → ℕ) (target : ℕ → ℚ) : Prop where
  eq : ∀ N, (demazureSumCoeff hilb N : ℚ) = target N

structure TargetLowerBoundQ {V : Type*} [AddCommGroup V] [Module ℂ V]
    (W : VertexOperatorAlgebra V) (target : ℕ → ℚ) : Prop where
  le : ∀ N : ℕ, target N ≤ (Module.finrank ℂ (W.grading N) : ℚ)

theorem demazureBoundFor_of_Q {V : Type*} [AddCommGroup V] [Module ℂ V]
    {W : VertexOperatorAlgebra V} {hilb : ℕ → ℕ → ℕ} {target : ℕ → ℚ}
    (hstar : StarIdentityQ hilb target) (hlow : TargetLowerBoundQ W target) :
    DemazureBoundFor W hilb
\end{leancode}
\begin{leancode}
def StarBoundQ (hilb : ℕ → ℕ → ℕ) (c : ℕ → ℚ) : Prop :=
  ∀ N, (demazureSumCoeff hilb N : ℚ) ≤ c N

theorem demazureBoundFor_of_bound {V : Type*} [AddCommGroup V] [Module ℂ V]
    {W : VertexOperatorAlgebra V} {hilb : ℕ → ℕ → ℕ} {target : ℕ → ℚ}
    (hstar : StarBoundQ hilb target) (hlow : TargetLowerBoundQ W target) :
    DemazureBoundFor W hilb
\end{leancode}
\begin{dictionary}
The casts \lean{(… : ℚ)} send $\N\to\Q$; the inequalities are in $\Q$ and \lean{exact_mod_cast} transports them
back to $\N$. \textbf{Why the one-sided form?} Downstream, the star identity is only ever used through
$D_N=t_N\le\dim W_N$, so $D_N\le t_N$ suffices; \lean{StarIdentityQ.toStarBoundQ} shows the $=$ form implies
it. Whether the $\le$ form is \emph{easier to prove} is a separate question. For proofs by termwise domination it
is not (Section~\ref{l25:sec:rigidity}); for other proofs the question is open.
\end{dictionary}
\begin{physics}
The same three statements as before, with the target allowed to be any rational sequence --- in practice the
MLDE solution, whose coefficients are a priori rational and turn out to be integers.
\end{physics}
\begin{status}
Definitions and transitivities; standard axioms.
\end{status}

\begin{warning}[title={Two different \texttt{mldeCoeff}}]
\lean{VOA.DeligneAudit.DeligneCharacter.mldeCoeff (h N : ℕ) : ℚ} (Definition~\ref{l25:def:mlde} below) is the
$N$-th coefficient of the AK solution. \lean{VOA.mldeCoeff (s : ℕ) (c : ℕ → E46) (h : ℕ → ℕ → E46) (j : ℕ) : E46}
in \leanpath{VOA/\allowbreak Modularity/\allowbreak AKMLDE.lean} (Chapter~\ref{ch:L16}) is the $j$-th coefficient
of a generic monic MLDE over $\C[E_4,E_6]$. They share only the short name (Section~\ref{l25:sec:defects}). In
this chapter \lean{mldeCoeff} always means the first.
\end{warning}

%======================================================================
\section{The Arakawa--Kawasetsu recurrence}\label{l25:sec:recurrence}
%======================================================================

\subsection{A primer on modular linear differential equations}\label{l25:sec:primer}

Write $\ltwentyfiveEu=q\,\frac{d}{dq}$ for the Euler operator, $\sigma_k(n)=\sum_{d\mid n}d^k$, and use the normalised
Eisenstein series $E_2=1-24\sum_{n\ge1}\sigma_1(n)q^n$, $E_4=1+240\sum_{n\ge1}\sigma_3(n)q^n$ and
$E_6=1-504\sum_{n\ge1}\sigma_5(n)q^n$ (Chapter~\ref{ch:L14}). This primer explains the shape of the equation. None of
it is used in the proofs, which are purely formal.

\paragraph{Modular forms and the Serre derivative.} A modular form $f$ of weight $k$ for $SL_2(\Z)$ satisfies
$f(\frac{a\tau+b}{c\tau+d})=(c\tau+d)^kf(\tau)$. Here $q=e^{2\pi i\tau}$, so $\ltwentyfiveEu=\frac1{2\pi i}\frac{d}{d\tau}$.
The ring of such forms is $\C[E_4,E_6]$. In particular there is no nonzero form of weight $2$, the forms of weight $4$
are the multiples of $E_4$, and the space of weight $8$ is spanned by $E_4^2$. The derivative $\ltwentyfiveEu f$ of a
weight-$k$ form is not modular, because differentiating $(c\tau+d)^k$ produces an extra term. The quasimodular series
$E_2$ has exactly the compensating transformation law, so the \emph{Serre derivative}
\begin{equation*}
\vartheta_k=\ltwentyfiveEu-\tfrac{k}{12}E_2
\end{equation*}
maps weight-$k$ forms to weight-$(k+2)$ forms. Ramanujan's identities (Theorem~\ref{l25:thm:ramanujan}) are the formal
shadow of this fact.

\paragraph{MLDEs.} Zhu's theorem says that the characters $f_M=\tr_Mq^{L_0-c/24}$ of a rational, $C_2$-cofinite VOA
form a vector-valued modular function, i.e.\ have weight $0$. Arakawa and Kawasetsu extend this to quasi-lisse VOAs.
For weight $0$, the natural order-$n$ operator is
$\vartheta^n=\vartheta_{2n-2}\circ\dots\circ\vartheta_2\circ\vartheta_0$, which raises the weight by $2n$. A
\emph{monic MLDE of order $n$} is
$\vartheta^nf+\sum_{j=1}^{n}\phi_j\,\vartheta^{n-j}f=0$, where each coefficient $\phi_j$ is a modular form of weight
$2j$. At order two the coefficient $\phi_1$ has weight $2$, so it vanishes, and $\phi_2$ has weight $4$, so it is a
multiple $-\lambda E_4$. Hence every second-order MLDE has the one-parameter form
\begin{equation*}
\vartheta_2\vartheta_0f=\lambda E_4f,\qquad \vartheta_2\vartheta_0=\big(\ltwentyfiveEu-\tfrac16E_2\big)\ltwentyfiveEu .
\end{equation*}

\paragraph{Why order two, and why weight four.} Zhu's recursion (Section~\ref{l25:sec:block}) assigns to each state
$v$ of square-bracket weight $k$ a torus one-point function $S_M(v)$, which is quasimodular of weight $k$. A relation
among states of weight $2n$ that expresses the ``$n$-th power of the stress tensor'' through simpler states becomes an
MLDE of order $n$. The state $L[-2]^2\ltwentyfivevac$, the square $(TT)$ of the stress tensor, has $L_0$-weight $4$.
Its one-point function is $\vartheta_2\vartheta_0f$ up to the Eisenstein terms of Theorem~\ref{l25:thm:tower}, a
weight-four quasimodular object. That the $L_0$-weight of $(TT)$ equals the modular weight of $E_4$ is therefore no
coincidence: Zhu's recursion turns the one into the other. A relation saying that $(TT)$ is, up to states with
vanishing or Eisenstein-type traces, ``classically zero'' (the vacuum quartic relation of
Section~\ref{l25:sec:block}) produces exactly a second-order MLDE. \textbf{No modularity is used in this chapter}: all
equations are identities of formal $q$-series, coefficient by coefficient.

\subsection{The Arakawa--Kawasetsu equation and its recurrence}

Arakawa and Kawasetsu proved (arXiv:1610.05865, \S7) the following. Let $V=L_k(\ltwentyfiveg)$ be in the Deligne series,
with $h=\ltwentyfivehv$ and $k=-1-h/6$. Then the normalised vacuum character $f=q^{-c/24}\sum_Nd_Nq^N$, with
$d_N=\dim V_N$, satisfies
\begin{equation}\label{l25:eq:mlde}
\Big(\ltwentyfiveEu-\tfrac{1}{6}E_2\Big)\ltwentyfiveEu f \;=\; \frac{h^2-1}{144}\,E_4\,f ,
\end{equation}
which is the second-order MLDE above with $\lambda=(h^2-1)/144$.

Substitute $f=q^{e}\sum d_Nq^N$ and read off the coefficient of $q^{N+e}$. Since $\ltwentyfiveEu q^{N+e}=(N+e)q^{N+e}$,
the coefficient of $q^e$ (the \emph{indicial equation}) is $e^2-e/6-(h^2-1)/144=0$. Equivalently
$(12e)^2-2(12e)-(h^2-1)=0$, with roots $12e=1\pm h$. The character has $e=-c/24=(h+1)/12$, the larger root. For
$N\ge1$ the coefficient of $d_N$ becomes $(N+e)^2-(N+e)/6-(h^2-1)/144=N(6N+h)/6$. Multiplying the $q^{N+e}$ equation by
$6$ gives
\begin{equation}\label{l25:eq:rec}
N(6N+h)\,d_N=\sum_{r=1}^{N}\Big[10(h^2-1)\sigma_3(r)-\big(24(N-r)+2(h+1)\big)\sigma_1(r)\Big]\,d_{N-r}.
\end{equation}
The factor $N(6N+h)$ never vanishes for $N\ge1$, so \eqref{l25:eq:rec} with $d_0=1$ determines the sequence. At $N=1$
only $r=1$ contributes, and $(6+h)d_1=10(h^2-1)-2(h+1)=2(h+1)(5h-6)$, so
\begin{equation*}
d_1=\frac{2(h+1)(5h-6)}{h+6}.
\end{equation*}
Deligne's dimension formula for the series is $\dim\ltwentyfiveg=2(h+1)(5h-6)/(h+6)$, so $d_1=\dim\ltwentyfiveg$. The
equation knows the number of currents $J^a_{-1}\ltwentyfivevac$, but only because $k$ is the Deligne level: the value
$c=-2(h+1)$ was used to fix $e$.

\begin{definition}[The AK coefficients]\label{l25:def:mlde}
$\sigma_k(n)=\sum_{d\mid n}d^k$. For $h\in\N$, $c_N=\mathtt{mldeCoeff}\ h\ N\in\Q$ is defined by $c_0=1$
and \eqref{l25:eq:rec} solved for $c_N$.
\end{definition}
\noindent File \leanpath{VOA/\allowbreak DeligneAudit/\allowbreak DeligneCharacterTypes.lean}.
\begin{leancode}
def sigmaK (k n : ℕ) : ℕ := ∑ d ∈ n.divisors, d ^ k

def mldeList (h : ℕ) : ℕ → List ℚ
  | 0 => [1]
  | n + 1 =>
      let l := mldeList h n
      let N := n + 1
      l ++ [(∑ r ∈ Finset.range N,
          ((10 * ((h : ℚ) ^ 2 - 1) * sigmaK 3 (r + 1) -
            (24 * ((N : ℚ) - (r + 1)) + 2 * ((h : ℚ) + 1)) * sigmaK 1 (r + 1)) *
              l.getD (N - (r + 1)) 0)) / ((N : ℚ) * (6 * N + h))]

def mldeCoeff (h N : ℕ) : ℚ := (mldeList h N).getD N 0
\end{leancode}
\begin{dictionary}
\lean{n.divisors} is Mathlib's finset of positive divisors ($\emptyset$ for $n=0$, so $\sigma_k(0)=0$).
\lean{mldeList h n} is the list $[c_0,\dots,c_n]$, built by appending one entry at a time;
\lean{l.getD i 0} reads entry $i$ (default $0$ out of range). The summation index $r\in\{0,\dots,N-1\}$ in Lean is
$r+1\in\{1,\dots,N\}$ on paper. The appended entry is exactly \eqref{l25:eq:rec} divided by $N(6N+h)$.
\lean{mldeCoeff h N} reads the last entry. \textbf{Division by zero?} \lean{(N : ℚ) * (6 * N + h)} with
$N\ge1$ is positive; Lean's $x/0=0$ convention never fires. \textbf{Is it the AK equation?} That the list
satisfies the recurrence is the next theorem; that the recurrence is \eqref{l25:eq:mlde} is the derivation above,
and the library re-derives it from the VOA (Theorem~\ref{l25:thm:kappa}).
\end{dictionary}
\begin{physics}
$c_N$ is the predicted number of states of weight $N$ in the vacuum module: the MLDE ``knows'' the vacuum character
once it knows $c$. The first values are in Table~\ref{l25:tab:cN}. The weight-one value $c_1=\dim\ltwentyfiveg$ counts
the currents $J^a_{-1}\ltwentyfivevac$ (derivation above).
\end{physics}
\begin{status}
Definition; computable.
\end{status}

\begin{table}[ht]\centering\small
\caption{The first AK coefficients $c_N=\mathtt{mldeCoeff}\ h^\vee\ N$. Every entry was computed from
\eqref{l25:eq:rec} in exact rational arithmetic for this chapter. Kernel-checked in the library: $A_2$ through $c_3$
(\lean{A2.mldeCoeff_A2_small}) and $G_2$ through $c_2$ (\lean{mldeCoeff_G2_small}). Kernel-checked in a scratch file for
this chapter, not in the library: $D_4$ through $c_2$.}\label{l25:tab:cN}
\begin{tabular}{llllllll}
\toprule
type & $h^\vee$ & $k$ & $c$ & $c_1$ & $c_2$ & $c_3$ & $c_4$\\
\midrule
$A_2$ & 3 & $-3/2$ & $-8$ & 8 & 36 & 128 & 394\\
$G_2$ & 4 & $-5/3$ & $-10$ & 14 & 92 & 456 & 1848\\
$D_4$ & 6 & $-2$ & $-14$ & 28 & 329 & 2632 & 16380\\
$F_4$ & 9 & $-5/2$ & $-20$ & 52 & 1106 & 14808 & 147239\\
$E_6$ & 12 & $-3$ & $-26$ & 78 & 2509 & 49270 & 698425\\
$E_7$ & 18 & $-4$ & $-38$ & 133 & 7505 & 254885 & 6093490\\
$E_8$ & 30 & $-6$ & $-62$ & 248 & 27249 & 1821002 & 85118126\\
\bottomrule
\end{tabular}
\end{table}

\begin{theorem}[The recurrence, and uniqueness]\label{l25:thm:unique}
Call $d:\N\to K$ ($K$ a field of characteristic $0$) a \emph{normalised AK solution} with parameter $h$ if $d_0=1$
and \eqref{l25:eq:rec} holds for all $N\ge1$. Then $c=\mathtt{mldeCoeff}\ h$ is one, and every normalised AK
solution equals $c$ (cast into $K$).
\end{theorem}
\noindent File \leanpath{VOA/\allowbreak DeligneAudit/\allowbreak DeligneMLDE.lean}.
\begin{leancode}
def akKernel (h N r : ℕ) : ℚ :=
  10 * ((h : ℚ) ^ 2 - 1) * sigmaK 3 r - (24 * ((N : ℚ) - r) + 2 * ((h : ℚ) + 1)) * sigmaK 1 r

structure IsAKSolution {K : Type*} [Field K] [CharZero K] (h : ℕ) (d : ℕ → K) : Prop where
  zero : d 0 = 1
  step : ∀ N : ℕ, 1 ≤ N →
    (N : K) * (6 * N + h) * d N
      = ∑ r ∈ Finset.range N, (akKernel h N (r + 1) : K) * d (N - (r + 1))

theorem mldeCoeff_isAKSolution (h : ℕ) : IsAKSolution (K := ℚ) h (mldeCoeff h) where

theorem eq_mldeCoeff_of_isAKSolution {K : Type*} [Field K] [CharZero K] {h : ℕ} {d : ℕ → K}
    (hd : IsAKSolution h d) : ∀ N, d N = (mldeCoeff h N : K)
\end{leancode}
\begin{dictionary}
\lean{akKernel h N r} is the bracket in \eqref{l25:eq:rec} with $r$ the summation index. \lean{IsAKSolution} is the
definition in Theorem~\ref{l25:thm:unique}: \lean{zero} is $d_0=1$, \lean{step} is \eqref{l25:eq:rec} for $N\ge1$. The
field $K$ is arbitrary of characteristic $0$; the library uses $K=\Q$ (for $q$-series) and $K=\C$ (for graded
dimensions of a complex VOA). The cast \lean{(mldeCoeff h N : K)} is the image of a rational under $\Q\to K$.
Same statement as on paper.
\end{dictionary}
\begin{proof}[Proof on paper]
The first claim is the definition unfolded (the list's last entry is the recurrence's right side divided by
$N(6N+h)$). Uniqueness: by strong induction on $N$, the right side of \eqref{l25:eq:rec} involves only $d_{N-r}$ with
$r\ge1$, equal to $c_{N-r}$ by induction, and $N(6N+h)\ne0$ in characteristic $0$, so $d_N=c_N$.
\end{proof}
\begin{physics}
A second-order ODE with a regular singularity at $q=0$ has a two-dimensional solution space, with characteristic
exponents $(1\pm h)/12$. \textbf{Uniqueness is not delicate here.} For the larger exponent, the vacuum's, the leading
factor of the recurrence is $N(6N+h)$, which is nonzero for every $N\ge1$, so the normalised solution exists and is
unique. A logarithmic solution can arise only from the \emph{smaller} exponent, when the exponent difference $h/6$ is a
positive integer ($h=6,12,18,30$), and that solution is never used. \emph{Heuristic reading:} the smaller exponent
$(1-h)/12$ would be the leading exponent of a second ``character'' of the two-character theory. That character would
belong to a module whose lowest conformal weight differs from the vacuum's by $-h/6$. Nothing in this chapter depends on
this reading.
\end{physics}
\begin{status}
Standard axioms (re-run for this chapter with \lean{lake env lean -DmaxSynthPendingDepth=3}). Unconditional.
\end{status}

\begin{theorem}[Graded dimensions solving the recurrence are the AK coefficients]\label{l25:thm:fromAK}
If the graded dimensions $N\mapsto\dim W_N\in\C$ of a VOA form a normalised AK solution with parameter $h$, then
$\dim W_N=c_N$ for all $N$, and in particular \lean{TargetLowerBoundQ W (mldeCoeff h)}.
\end{theorem}
\noindent File \leanpath{VOA/\allowbreak DeligneAudit/\allowbreak DeligneMLDE.lean}.
\begin{leancode}
theorem finrank_eq_mldeCoeff_of_isAKSolution (W : VertexOperatorAlgebra V) {h : ℕ}
    (hW : IsAKSolution h (fun N : ℕ => (Module.finrank ℂ (W.grading N) : ℂ))) (N : ℕ) :
    (Module.finrank ℂ (W.grading N) : ℚ) = mldeCoeff h N

theorem targetLowerBoundQ_of_isAKSolution (W : VertexOperatorAlgebra V) {h : ℕ}
    (hW : IsAKSolution h (fun N : ℕ => (Module.finrank ℂ (W.grading N) : ℂ))) :
    TargetLowerBoundQ W (mldeCoeff h)
\end{leancode}
\begin{dictionary}
The hypothesis casts $\dim W_N\in\N$ into $\C$; the conclusion is an equality in $\Q$ (the cast $\Q\to\C$ is
injective). The second theorem uses only the $\ge$ half of the equality and discards the rest: this is why the
per-type theorems of Section~\ref{l25:sec:alltypes}, which go through it, export only an inequality.
\end{dictionary}
\begin{proof}[Proof on paper]
Uniqueness (Theorem~\ref{l25:thm:unique}) over $K=\C$, then injectivity of $\Q\to\C$.
\end{proof}
\begin{physics}
``If the character satisfies the MLDE with the right leading exponent, it \emph{is} the MLDE solution.'' The
whole difficulty is moved into proving the MLDE for the VOA.
\end{physics}
\begin{status}
Standard axioms; unconditional as an implication.
\end{status}

\begin{example}[Kernel checks of small values]\label{l25:ex:small}
The library checks by \lean{decide +kernel}:
\begin{leancode}
theorem mldeCoeff_G2_small : (List.range 3).map (mldeCoeff 4) = [1, 14, 92] := by decide +kernel
\end{leancode}
\noindent (file \leanpath{VOA/\allowbreak DeligneAudit/\allowbreak DeligneCharacterTypes.lean}) and
\lean{A2.mldeCoeff_A2_small}: $[c_0,\dots,c_3]=[1,8,36,128]$ for $h=3$ (file \leanpath{DeligneMLDE.lean}). These pin
the conventions (a wrong sign in \lean{akKernel} would already fail at $c_1$, which must be $\dim\ltwentyfiveg$).
\end{example}

\begin{example}[$A_2$ from scratch]\label{l25:ex:a2}
This example runs every layer of the chapter on the smallest case, $\ltwentyfiveg=\mathfrak{sl}_3$, $h=3$, $k=-3/2$,
$c=-8$. All numbers below were recomputed by hand. The statements marked (L) are kernel checks in the library, and
those marked (S) are kernel checks (\lean{decide +kernel}) that we ran in a scratch file for this chapter.

\emph{The recurrence.} With $h=3$ the bracket in \eqref{l25:eq:rec} is $80\,\sigma_3(r)-(24(N-r)+8)\,\sigma_1(r)$.
\begin{itemize}
\item $N=1$: $1\cdot9\cdot d_1=(80-8)\cdot d_0=72$, so $d_1=8=\dim\mathfrak{sl}_3$.
\item $N=2$: the bracket is $80-32=48$ at $r=1$ and $80\cdot9-8\cdot3=696$ at $r=2$, so $2\cdot15\cdot d_2=48\cdot8+696\cdot1=1080$
  and $d_2=36$.
\end{itemize}
So $c_0,c_1,c_2=1,8,36$, as in Table~\ref{l25:tab:cN} and \lean{A2.mldeCoeff_A2_small} (L).

\emph{The constant.} $\kappa=\frac54c^2+\frac{11}2c=80-44=36=(h+1)(5h-6)$, the number in \lean{A2.vacuumQuarticRelation}
(Theorem~\ref{l25:thm:alltypes}).

\emph{The product.} $\prod_{r\ge1}(1+q^r)^8=(1+q)^8(1+q^2)^8\cdots=1+8q+(28+8)q^2+\dots=1+8q+36q^2+\dots$, matching
$c_N$ (\lean{productCoeff_small} (L): $1,8,36,128,394$; Theorem~\ref{l25:thm:a2prod} proves it for all $N$).

\emph{The Demazure sum.} Write $\mathrm{hilb}(n,\cdot)$ for the $A_2$ table (Theorem~\ref{l25:thm:a2star}). It has
$\mathrm{hilb}(0,\cdot)=1$ and $\mathrm{hilb}(1,\cdot)=8+q$ (\lean{demazureHilb_one} (L)), and
$\mathrm{hilb}(2,\cdot)=27+28q+17q^2+8q^3+q^4$ (S). Its total $81=9^2$ is the $q=1$ factorisation of
Theorem~\ref{l25:thm:total}, and its constant term $27=\dim V(2\theta)$ is the weight-$0$ part of $D(1,2\theta)$.
Finally $\mathrm{hilb}(3,0)=64=\dim V(3\theta)$ (S). (For $\mathfrak{sl}_3$, $\dim V(a\omega_1+b\omega_2)=(a+1)(b+1)(a+b+2)/2$
and $\theta=\omega_1+\omega_2$.) Then by Definition~\ref{l25:def:dsum}:
\begin{align*}
D_1&=\mathrm{hilb}(1,0)=8,\\
D_2&=\big[\mathrm{hilb}(1,0)+\mathrm{hilb}(1,1)\big]+\mathrm{hilb}(2,0)=9+27=36,\\
D_3&=\big[\mathrm{hilb}(1,0)+\mathrm{hilb}(1,1)\big]+\big[\mathrm{hilb}(2,0)+\mathrm{hilb}(2,1)\big]+\mathrm{hilb}(3,0)=9+55+64=128 .
\end{align*}
In $D_2$, the term with $n=1$ is $[q^1]\,(8+q)/(1-q)=9$ and the term with $n=2$ is $[q^0]\,\mathrm{hilb}(2,\cdot)/((1-q)(1-q^2))=27$.
These agree with \lean{demazureCoeff_small} (L) and with $c_N$: the star identity holds at $N\le3$. The GT table
$\mathrm{hilb}(2,\cdot)$ also agrees with the affine algorithm \lean{levelOneDemazureHilb cartanA2 thetaA2 wordA2 2} in
degrees $0$--$3$ (S). This is a second instance of the comparison between the two descriptions (Caveat~\ref{l25:cav:U7}).

\emph{The VOA.} \lean{A2.finrank_grading_one} gives $\dim V_1=8$, and \lean{A2.finrank_grading_two_ge} gives
$\dim V_2\ge36$ by a Gram determinant. Both are independent of the MLDE. The chain \eqref{l25:eq:chain} therefore
reads $8\le8\le8$ at $N=1$ and $36\le36\le\dim V_2$ at $N=2$.

\emph{A $D_4$ check.} At $h=6$: $(6+6)d_1=10\cdot35-2\cdot7=336$, so $d_1=28=\dim\mathfrak{so}_8$. The closed form of
Theorem~\ref{l25:thm:closed} gives $[q^{N+1}]\,\big(\tfrac1{240}\ltwentyfiveEu E_4\big)\,(q)_\infty^{-10}$. Here
$\frac1{240}\ltwentyfiveEu E_4=q+18q^2+\dots$ and $(q)_\infty^{-10}=1+10q+\dots$, giving $1$ at $N=0$ and $18+10=28$ at $N=1$.
\end{example}

%======================================================================
\section{From one weight-four trace identity to the MLDE}\label{l25:sec:block}
%======================================================================

\subsection{The Virasoro tower over the vacuum}

\paragraph{Torus one-point functions.} Put a chiral CFT on a torus with modular parameter $\tau$, $q=e^{2\pi i\tau}$,
and insert one field $v$ at a point. Translation invariance along the torus makes the answer independent of the point.
It equals the trace
\begin{equation*}
S_M(v)=\tr_M\,o(v)\,q^{L_0-c/24},
\end{equation*}
where $M$ is the module propagating around the cycle and $o(v)=v_{(\wt v-1)}$ is the zero mode of $v$, the part of the
field that preserves $L_0$-weight (Chapter~\ref{ch:L15}, Definition~\ref{l15:def:torusonepoint}). Three examples fix
the meaning.
\begin{itemize}
\item $v=\ltwentyfivevac$: $o(\ltwentyfivevac)=\mathrm{id}$, so $S_M(\ltwentyfivevac)=\tr_Mq^{L_0-c/24}=:f_M$ is the
  character. For $M=V$ it is the normalised vacuum character $f=q^{-c/24}\sum_N\dim V_N\,q^N$.
\item $v=\omega$, the conformal vector: $o(\omega)=L_0$, so $S_M(\omega)=\tr_ML_0q^{L_0-c/24}$, which is
  $\ltwentyfiveEu f_M+\frac c{24}f_M$.
\item $v=(TT)$, the normal-ordered square of the stress tensor: its one-point function is what the MLDE is about.
\end{itemize}
The library stores $S_M(v)$ as a formal $q$-series with complex exponents, i.e.\ a function $\C\to\C$ sending an
exponent to its coefficient. The exponents lie in $-c/24+\N$ for the modules used here.

\paragraph{Zhu's recursion and square-bracket modes.} The one idea of this section is Zhu's: there is a second set of
modes $a[n]$, the \emph{square-bracket modes}, defined from the change of coordinate $z\mapsto e^z-1$ from the plane to
the cylinder (Chapter~\ref{ch:L15}, Definition~\ref{l15:def:sqmode}). For them the trace of a descendant reduces to
traces of simpler states with \emph{Eisenstein series as coefficients}. For the stress tensor, with $L[n]$ the
square-bracket Virasoro modes (\lean{Lsq W n}), it reads
\begin{equation}\label{l25:eq:zhurec}
S_M\big(L[-2]v\big)=\ltwentyfiveEu\,S_M(v)+\sum_{m\ge1}G_{2m}\star S_M\big(L[2m-2]v\big).
\end{equation}
Here $\star$ is the product of $q$-series (Lean \lean{qConv}, the Cauchy product of an $\N$-indexed series with a
formal series), and $G_{2m}$ is Zhu's normalisation of the Eisenstein series (\lean{zhuG m}). These have
$G_2=-E_2/12$ and $G_4=E_4/720$: constant term $-B_{2m}/(2m)!$ and $q^n$-coefficient $\frac{2}{(2m-1)!}\sigma_{2m-1}(n)$.
\lean{qDeriv} multiplies the coefficient at exponent $e$ by $e$, i.e.\ it is $\ltwentyfiveEu$. Physically,
\eqref{l25:eq:zhurec} is the torus Ward identity for $T$. Integrating $T$ around the cycle gives the derivative in
$\tau$, and the OPE of $T$ with $v$, weighted by the Weierstrass function of the torus, gives the Eisenstein
corrections.

\emph{Check at $v=\ltwentyfivevac$.} $L[n]\ltwentyfivevac=0$ for $n\ge-1$, so every term of the sum vanishes, and
$S_M(L[-2]\ltwentyfivevac)=\ltwentyfiveEu f_M$. This agrees with the second example above, because
$L[-2]\ltwentyfivevac=\omega-\frac c{24}\ltwentyfivevac$: the shift by $c/24$ is the Schwarzian of the change to the
cylinder. Write $f=S_V(\ltwentyfivevac)$ (\lean{torusPartition}) for the normalised vacuum character of $V$.

\begin{theorem}[The second Virasoro descendant of the vacuum]\label{l25:thm:tower}
For every VOA $W$ and every ordinary module $M$,
\begin{equation*}
S_M\big(L[-2]^2\ltwentyfivevac\big)=\ltwentyfiveEu^2f_M+G_2\star(2\ltwentyfiveEu f_M)+G_4\star\big(\tfrac c2f_M\big),
\qquad f_M=S_M(\ltwentyfivevac).
\end{equation*}
\end{theorem}
\noindent File \leanpath{VOA/\allowbreak DeligneAudit/\allowbreak DeligneMLDEBlock.lean}.
\begin{leancode}
theorem torusOnePoint_Lsq_sq_vacuum (N : OrdVModule W M) :
    torusOnePoint N (Lsq W (-2) (Lsq W (-2) W.vacuum))
      = qDeriv (qDeriv (torusPartition N))
        + (qConv (zhuG 1) ((2 : ℂ) • qDeriv (torusPartition N))
          + qConv (zhuG 2) ((W.centralCharge / 2) • torusPartition N))
\end{leancode}
\begin{dictionary}
\lean{OrdVModule W M}: an ordinary module (Chapter~\ref{ch:L12}). \lean{Lsq W (-2) (Lsq W (-2) W.vacuum)} is
$L[-2]L[-2]\ltwentyfivevac$. \lean{zhuG 1} is $G_2$ and \lean{zhuG 2} is $G_4$ (index $j$ means weight $2j$). The
statement is exactly the display. No hypothesis on $W$ or $M$ beyond their structures. The traces make sense
because an \lean{OrdVModule} carries, as a field, that each weight space is finite-dimensional, and that the weights are
bounded below (Chapter~\ref{ch:L12}).
\end{dictionary}
\begin{proof}[Proof on paper]
Apply Zhu's recursion \eqref{l25:eq:zhurec} to
$v=L[-2]\ltwentyfivevac$: the square-bracket Virasoro relations give $L[0]L[-2]\ltwentyfivevac=2L[-2]\ltwentyfivevac$,
$L[2]L[-2]\ltwentyfivevac=\frac c2\ltwentyfivevac$ and $L[2m-2]L[-2]\ltwentyfivevac=0$ for $m\ge3$ (the exact tower
\lean{Lsq_tower} of Chapter~\ref{ch:L16}); and $S(L[-2]\ltwentyfivevac)=\ltwentyfiveEu f$ by the check above. So
$S(L[-2]^2\ltwentyfivevac)=\ltwentyfiveEu(\ltwentyfiveEu f)+G_2\star S(2L[-2]\ltwentyfivevac)+G_4\star S(\frac c2\ltwentyfivevac)$,
which is the display.
\end{proof}
\begin{physics}
$L[-2]^2\ltwentyfivevac$ is the state of the normal-ordered square $(TT)$ of the stress tensor on the
torus (in Zhu's coordinates). Its torus one-point function is the second $\tau$-derivative of the partition
function, corrected by the Eisenstein series that implement the change from the plane to the torus. This is the
standard physics derivation of MLDEs (Mathur--Mukhi--Sen; Beem--Rastelli for 4d theories): a null-type relation for
$(TT)$ becomes a differential equation for the character.
\end{physics}
\begin{status}
Standard axioms. Unconditional.
\end{status}

\subsection{The vacuum quartic relation and the indicial equation}

\begin{definition}[Vacuum quartic relation; formal MLDE]\label{l25:def:vqr}
For a VOA $W$ and $\kappa\in\C$, the \emph{vacuum quartic relation} with constant $\kappa$ is the trace identity
$S_W(L[-2]^2\ltwentyfivevac)=\kappa\,G_4\star f$ on the adjoint module. The \emph{formal MLDE} with constant $\kappa$ is
$\ltwentyfiveEu^2f+G_2\star(2\ltwentyfiveEu f)+G_4\star(\frac c2f)=\kappa\,G_4\star f$.
\end{definition}
\noindent File \leanpath{VOA/\allowbreak DeligneAudit/\allowbreak DeligneMLDEBlock.lean}.
\begin{leancode}
structure VacuumQuarticRelation (W : VertexOperatorAlgebra 𝒱) (κ : ℂ) : Prop where
  eq : torusOnePoint (adjointModule W) (Lsq W (-2) (Lsq W (-2) W.vacuum))
    = κ • qConv (zhuG 2) (torusPartition (adjointModule W))

def FormalMLDE (κ : ℂ) : Prop :=
  qDeriv (qDeriv (torusPartition (adjointModule W)))
      + (qConv (zhuG 1) ((2 : ℂ) • qDeriv (torusPartition (adjointModule W)))
        + qConv (zhuG 2) ((W.centralCharge / 2) • torusPartition (adjointModule W)))
    = κ • qConv (zhuG 2) (torusPartition (adjointModule W))
\end{leancode}
\begin{dictionary}
\lean{adjointModule W} is $W$ as a module over itself; \lean{torusPartition (adjointModule W)} is the formal series
with coefficient $\dim W_n$ at exponent $n-c/24$ (\lean{torusPartition_adjoint}). By
Theorem~\ref{l25:thm:tower}, the relation implies the formal MLDE with the same $\kappa$
(\lean{formalMLDE_of_vacuumQuarticRelation}, the only direction used); the converse follows from
Theorem~\ref{l25:thm:tower} by the same rewriting but is not stated in Lean. \textbf{Section variables.}
\lean{W : VertexOperatorAlgebra 𝒱} is a section variable; in the \lean{Coefficients} and \lean{Algebraic} sections of
\leanpath{DeligneMLDEBlock.lean} it is \emph{explicit} (\lean{FormalMLDE W κ}, \lean{kappa_eq_of_formalMLDE W h …},
\lean{isAKSolution_of_vacuumQuarticRelation W h …}), in the \lean{Tower} section implicit. Rewriting with
$G_2=-E_2/12$, $G_4=E_4/720$, the formal MLDE reads $(\ltwentyfiveEu-\frac16E_2)\ltwentyfiveEu f=\frac{\kappa-c/2}{720}E_4f$, which
is \eqref{l25:eq:mlde} exactly when $\kappa-c/2=5(h^2-1)$. \textbf{Formal, not analytic:} equality of functions
$\C\to\C$ of exponents, i.e.\ coefficientwise; no convergence or modularity enters.
\end{dictionary}
\begin{physics}
Why a state of the form $a[-2]b$ contributes \emph{exactly} a multiple of $E_4$: Zhu's recursion for $a[-2]$ has
coefficients $(2m-1)G_{2m}$ multiplying the traces of $a[2m-2]b$. The $m=1$ term is the trace of a zero mode, which
vanishes. When $a[2]b$ is a multiple $\lambda\ltwentyfivevac$ of the vacuum and higher positive modes kill $b$, only
$m=2$ survives, and $S(a[-2]b)=-3\lambda G_4\star f$ (Definition~\ref{l25:def:sqc2}). In CFT terms, the OPE $a(z)b(0)$ has
a single vacuum term at order $z^{-4}$ in the cylinder coordinate, and the torus propagator turns that term into
$E_4$. The relation says that $(TT)$ behaves like such a state on the torus. Its one-point function is then $E_4$ times
the partition function, which is the MLDE.
\end{physics}
\begin{status}
Definitions (a named \lean{Prop}); satisfiable --- it is proved for seven VOAs in
Section~\ref{l25:sec:alltypes}.
\end{status}

\begin{theorem}[The constant is forced; the AK recurrence]\label{l25:thm:kappa}
Let $W$ have central charge $c=-2(h+1)$ and $\ltwentyfivevac\ne0$. If the formal MLDE (or the vacuum
quartic relation, which implies it) holds with \emph{some} constant $\kappa$, then $\kappa=(h+1)(5h-6)$, and the graded dimensions
$N\mapsto\dim W_N$ form a normalised AK solution with parameter $h$. Hence $\dim W_N=c_N$ and
\lean{TargetLowerBoundQ W (mldeCoeff h)}.
\end{theorem}
\noindent File \leanpath{VOA/\allowbreak DeligneAudit/\allowbreak DeligneMLDEBlock.lean}.
\begin{leancode}
theorem kappa_eq_of_formalMLDE (h : ℕ) (hc : W.centralCharge = -2 * ((h : ℂ) + 1))
    (hvac : W.vacuum ≠ 0) {κ : ℂ} (hE : FormalMLDE W κ) :
    κ = ((h : ℂ) + 1) * (5 * h - 6)

theorem isAKSolution_of_vacuumQuarticRelation (h : ℕ)
    (hc : W.centralCharge = -2 * ((h : ℂ) + 1)) (hvac : W.vacuum ≠ 0) {κ : ℂ}
    (hrel : VacuumQuarticRelation W κ) :
    IsAKSolution h (fun N : ℕ => (Module.finrank ℂ (W.grading N) : ℂ))

theorem targetLowerBoundQ_of_vacuumQuarticRelation (W : VertexOperatorAlgebra 𝒱) (h : ℕ)
    (hc : W.centralCharge = -2 * ((h : ℂ) + 1)) (hvac : W.vacuum ≠ 0)
    {κ : ℂ} (hrel : VacuumQuarticRelation W κ) :
    DeligneCharacter.TargetLowerBoundQ W (DeligneCharacter.mldeCoeff h)
\end{leancode}
\begin{dictionary}
Three hypotheses: \lean{hc} ($c=-2(h+1)$, true for $L_k(\ltwentyfiveg)$ in the Deligne series because
$\dim\ltwentyfiveg=2(h+1)(5h-6)/(h+6)$ and $c=k\dim\ltwentyfiveg/(k+h)$), \lean{hvac} ($W$ is not the zero space; it
makes $\dim W_0=1$, because the library's \lean{VertexOperatorAlgebra} has as a field \lean{vacuum_grade} that
$W_0=\C\ltwentyfivevac$, i.e.\ CFT type, Chapter~\ref{ch:01}), and the relation with an \emph{existentially free} $\kappa$. \textbf{Why is $\kappa$ free?}
Because it is determined: Lean need not be told its value, which removes a possible source of error (a wrong
$\kappa$ in a hypothesis would make it unsatisfiable, hence the theorem vacuous). \textbf{Could the conclusion be
vacuous?} No: \lean{TargetLowerBoundQ} asserts $\dim W_N\ge c_N$ with $c_1=\dim\ltwentyfiveg>0$ etc., which fails for
the zero VOA; and the hypotheses are discharged for seven concrete VOAs below.
\end{dictionary}
\begin{proof}[Proof on paper]
Read the formal MLDE at the exponent $e+N$, $e=-c/24$, where the coefficient of $f$ is $\dim W_N$. At $N=0$ only
the constant terms of $G_2,G_4$ contribute, and $\dim W_0=1$:
$e^2-\frac{1}{12}\cdot2e+\frac1{720}\cdot\frac c2=\frac{\kappa}{720}$, i.e.\ $\kappa=720e^2-120e+\frac c2
=\frac54c^2+\frac{11}2c$. At $c=-2(h+1)$ this is $5(h+1)^2-11(h+1)=(h+1)(5h-6)$. Then
$\kappa-c/2=5c(c+4)/4=5(h^2-1)$, so the MLDE is \eqref{l25:eq:mlde}, and the $q^{e+N}$ coefficient for $N\ge1$ is
\eqref{l25:eq:rec} after multiplying by $6$ (the computation of Section~\ref{l25:sec:recurrence}).
Theorem~\ref{l25:thm:fromAK} concludes.
\end{proof}
\begin{physics}
The indicial equation says that the vacuum (weight $-c/24$ on the torus) is a solution's leading term; this fixes
the coefficient of $E_4$ in terms of $c$ alone. In physics language this is the ``holomorphic modular bootstrap''
with the vacuum as input: the one free parameter of a second-order MLDE is fixed by the central charge.
\emph{Observation (heuristic, not used):} Mathur, Mukhi and Sen classified the second-order MLDEs at positive central
charge whose vacuum solution has nonnegative integer coefficients, and level-one WZW models of the Deligne series
appear among them. That the same Lie algebras reappear here at negative central charge, in the 4d chiral algebras,
is a known coincidence of the two families. It is not a theorem of either source, and nothing here depends on it.
\end{physics}
\begin{status}
Standard axioms for \lean{isAKSolution_of_vacuumQuarticRelation} and \lean{kappa_eq_of_formalMLDE} (re-run for this
chapter). Conditional on \lean{hc}, \lean{hvac},
\lean{hrel} as stated.
\end{status}

\begin{warning}[title={No modular forms, no Etingof-finiteness: what this replaces}]
Chapter~\ref{ch:L16} shows that the library's general AK theorem
(\lean{IsSqPrimary.exists_monicMLDEOn_all_analyticBlocks_of_etingofFinite_uncond}) is conditional on
\lean{EtingofFinite} only (its side inputs \lean{AKWeightGrading} and \lean{E46SerreClosed} are now the theorems
\lean{VOA.akWeightGrading} and \lean{VOA.e46SerreClosed}) but bounds neither the order nor the coefficients of its
MLDE, and that the unconditional \lean{character_formalODE_of_etingofFinite} is near-vacuous. Those files are
\emph{not used} here. The route of this section needs only Zhu's formal recursion and the exact Virasoro tower; the
target is coefficientwise, so the formal version suffices, and $\kappa$ comes from the indicial equation. \textbf{Dead end recorded:} do not route a coefficient identity through
\lean{AKModularODE.character_formalODE_of_etingofFinite}; it does not determine the character.
\end{warning}

\subsection{The algebraic form of the relation}

The trace identity is a statement about traces. Arakawa and Kawasetsu derive it from an algebraic statement,
$\bar\omega^2=0$ in $R_V$, and the library has a square-bracket version of that statement. To motivate the definitions,
start from Zhu's recursion for a square-bracket mode $a[-2]$, with $a$ homogeneous:
\begin{equation}\label{l25:eq:zhuminus2}
S\big(a[-2]b\big)=-\sum_{m\ge1}(2m-1)\,G_{2m}\star S\big(a[2m-2]b\big).
\end{equation}
In general the right side involves infinitely many unknown traces. It collapses to a single known term in exactly
the following situation. The $m=1$ term is $S(a[0]b)$, the trace of a commutator-type state, which vanishes; we
therefore call such states \emph{trace-null}. The $m=2$ term is $3G_4\star S(a[2]b)$, which is a multiple of $G_4\star f$
precisely when $a[2]b$ is a multiple of the vacuum. The terms $m\ge3$ vanish if $a[2m-2]b=0$. The definition records
exactly these conditions.

\begin{definition}[Trace-null vectors; square-bracket $C_2$ terms]\label{l25:def:sqc2}
The \emph{trace-null} subspace is spanned by the $a[0]v$ ($a$ homogeneous of weight $\ge1$) and by the $L[-1]v$. A
\emph{square-bracket $C_2$ term of weight four} is a state $a[-2]b$ with $a$ homogeneous of weight $k_a\ge1$,
$a[2]b=\lambda\ltwentyfivevac$, and $a[2m]b=0$ for $m\ge2$. By \eqref{l25:eq:zhuminus2} its one-point function is
$-3\lambda\,G_4\star f$. The \emph{square-bracket $C_2$ quartic relation} with constant $\kappa$ says that
$L[-2]^2\ltwentyfivevac\equiv\sum_ia_i[-2]b_i$ modulo trace-null vectors, for finitely many such terms, with
$\kappa=-3\sum_i\lambda_i$.
\end{definition}
\noindent File \leanpath{VOA/\allowbreak DeligneAudit/\allowbreak DeligneMLDEBlock.lean}.
\begin{leancode}
noncomputable def traceNull : Submodule ℂ 𝒱 :=
  Submodule.span ℂ ({x | ∃ (ka : ℕ) (a v : 𝒱), 1 ≤ ka ∧ a ∈ W.grading (ka : ℤ) ∧ x = vSqZ W a ka 0 v}
    ∪ Set.range (Lsq W (-1)))

structure SqC2Term (W : VertexOperatorAlgebra 𝒱) where
  ka : ℕ
  a : 𝒱
  b : 𝒱
  lam : ℂ
  one_le : 1 ≤ ka
  mem : a ∈ W.grading (ka : ℤ)
  two : vSqZ W a ka 2 b = lam • W.vacuum
  high : ∀ m : ℕ, 2 ≤ m → vSqZ W a ka (2 * (m : ℤ)) b = 0

structure SqC2QuarticRelation (W : VertexOperatorAlgebra 𝒱) (κ : ℂ) : Prop where
  rel : ∃ (s : Finset ℕ) (t : ℕ → SqC2Term W),
    Lsq W (-2) (Lsq W (-2) W.vacuum) - ∑ i ∈ s, vSqZ W (t i).a (t i).ka (-2) (t i).b
        ∈ traceNull W
      ∧ κ = -3 * ∑ i ∈ s, (t i).lam

theorem vacuumQuarticRelation_of_sqC2 {κ : ℂ} (h : SqC2QuarticRelation W κ) :
    VacuumQuarticRelation W κ
\end{leancode}
\begin{dictionary}
\lean{vSqZ W a ka n} is the square-bracket mode $a[n]$ of a state $a$ of $L_0$-weight \lean{ka}
(Chapter~\ref{ch:L15}, Definition~\ref{l15:def:sqmode}); the weight must be supplied because square modes depend
on it. \lean{Set.range (Lsq W (-1))} is the image of $L[-1]$. \lean{SqC2Term} is data (a \lean{structure} in
\lean{Type}), carrying the two vanishing conditions as fields.
\end{dictionary}
\begin{proof}[Proof on paper]
Every one-point function kills trace-null vectors (cyclicity of the trace for $a[0]$; $S(L[-1]v)=0$). For one term,
Zhu's recursion for $a[-2]$ gives $S(a[-2]b)=-\sum_{m\ge1}(2m-1)G_{2m}\star S(a[2m-2]b)$; the $m=1$ term is a
zero-mode trace and vanishes, the $m=2$ term is $3G_4\star S(\lambda\ltwentyfivevac)$, and $m\ge3$ vanish by
\lean{high}. So $S(a[-2]b)=-3\lambda G_4\star f$; sum.
\end{proof}
\begin{physics}
Modulo zero modes and derivatives, which have vanishing torus one-point functions, $(TT)$ is a sum of
``$\partial^2$-descendant'' states, each contributing a multiple of $E_4$. This is the trace version of
``$\bar\omega^2=0$ in $R_V$'', i.e.\ the square of the Sugawara Casimir vanishes on the Higgs branch.
\end{physics}
\begin{status}
Standard axioms. Unconditional as an implication. Not itself used by the per-type instances, which go through the
current calculus of the next section.
\end{status}

%======================================================================
\section{The square-bracket current calculus}\label{l25:sec:zhu}
%======================================================================

To prove the vacuum quartic relation for $L_k(\ltwentyfiveg)$ one needs to compute with square-bracket modes of the
currents $J^a\in V_1$. The ``obvious'' route --- Zhu's isomorphism $(V,Y)\cong(V,Y[\,])$ --- is not in the
library and would be a large project (recorded as a dead end at the end of this section). Instead, the files
\leanpath{VOA/\allowbreak Modularity/\allowbreak ZhuCoordinate*.lean} prove directly the three relations needed.

\begin{theorem}[Square-bracket current algebra]\label{l25:thm:sqcomm}
Let $W$ be any VOA. For $a,b\in W_1$ with $a_{(1)}b=\beta\ltwentyfivevac$, and all $m,n\in\Z$,
\begin{equation*}
[a[m],b[n]]=(a_{(0)}b)[m+n]+m\beta\,\delta_{m+n,0}.
\end{equation*}
If moreover $a$ is primary ($L_1a=0$), then $[L[m],a[n]]=-n\,a[m+n]$. For $a\in W_{P+1}$,
$(L_{-1}a)[n]=-n\,a[n-1]-(P+1)\,a[n]$ (the translation rule).
\end{theorem}
\noindent File \leanpath{VOA/\allowbreak Modularity/\allowbreak ZhuCoordinateCurrent.lean}.
\begin{leancode}
theorem sqCurrent_comm {a b : 𝒱} (ha : a ∈ W.grading 1) (hb : b ∈ W.grading 1)
    {β : ℂ} (hp : W.smul a 1 b = β • W.vacuum) (m n : ℤ) (x : 𝒱) :
    vSqZ W a 1 m (vSqZ W b 1 n x) - vSqZ W b 1 n (vSqZ W a 1 m x)
      = vSqZ W (W.smul a 0 b) 1 (m + n) x + (if m + n = 0 then ((m : ℂ) * β) • x else 0)

theorem Lsq_sqCurrent_comm {a : 𝒱} (ha : a ∈ W.grading 1) (ha1 : W.L 1 a = 0) (m n : ℤ)
    (x : 𝒱) :
    Lsq W m (vSqZ W a 1 n x) - vSqZ W a 1 n (Lsq W m x) = (-(n : ℂ)) • vSqZ W a 1 (m + n) x

theorem vSqZ_L_neg_one {a : 𝒱} (P : ℕ) (ha : a ∈ W.grading ((P + 1 : ℕ) : ℤ)) (n : ℤ) :
    vSqZ W (W.L (-1) a) (P + 2) n
      = (-(n : ℂ)) • vSqZ W a (P + 1) (n - 1) - ((P : ℂ) + 1) • vSqZ W a (P + 1) n
\end{leancode}
\begin{dictionary}
\lean{W.smul a n b} is the ordinary mode product $a_{(n)}b$ (Chapter~\ref{ch:01}); so \lean{hp} is
$a_{(1)}b=\beta\ltwentyfivevac$ and \lean{W.smul a 0 b} is $a_{(0)}b$. \lean{vSqZ W a 1 m} is $a[m]$ for a weight-one
state. \lean{W.L 1 a = 0}: $a$ is primary. The ``if'' encodes the Kronecker delta. All three are the displays,
for an arbitrary VOA.
\end{dictionary}
\begin{proof}[Proof on paper]
\emph{Square modes as residues.} Zhu's definition is $a[n]=\Res_z Y(a,e^z-1)e^{z\,\wt a}z^{n}$. Substitute
$w=e^z-1$, so that $z=L(w):=\log(1+w)$ and $dz=dw/(1+w)$. Then
\begin{equation*}
a[n]=\Res_w\,Y(a,w)\,(1+w)^{\wt a-1}L(w)^n .
\end{equation*}
Since $L(w)=w\,u(w)$ with $u=\log(1+w)/w=1-\frac w2+\dots$ invertible, this is a finite combination
$a[n]=\sum_{p\ge0}z(\wt a,n,n+p)\,a_{(n+p)}$ of ordinary modes, with $z(k,m,m+p)=[w^p]\big((1+w)^{k-1}u^m\big)$. For
$\wt a=1$ the factor $(1+w)^0$ disappears, and $a[n]=\Res_wY(a,w)L(w)^n$.

\emph{The bracket, in residue form.} For weight-one $a,b$ the ordinary commutator formula is
$[a_{(i)},b_{(j)}]=(a_{(0)}b)_{(i+j)}+i\,(a_{(1)}b)_{(i+j-1)}$, since $a_{(k)}b=0$ for $k\ge2$ by weight. With
$a_{(1)}b=\beta\ltwentyfivevac$ and $\ltwentyfivevac_{(n)}=\delta_{n,-1}$, the second term is $i\beta\,\delta_{i+j,0}$.
Write $a[m]=\sum_ix_ia_{(i)}$ and $b[n]=\sum_jy_jb_{(j)}$, where $\sum_ix_iw^i=L(w)^m$ and $\sum_jy_jw^j=L(w)^n$.
\begin{itemize}
\item \emph{Bracket term.} $\sum_{i,j}x_iy_j(a_{(0)}b)_{(i+j)}$ is the square mode of $a_{(0)}b$ whose
  coefficient series is the product $L^m\cdot L^n=L^{m+n}$. That is $(a_{(0)}b)[m+n]$, the identity
  $u^mu^n=u^{m+n}$.
\item \emph{Central term, worked out completely.} $\sum_{i,j}x_iy_j\,i\,\beta\,\delta_{i+j,0}=\beta\sum_iix_iy_{-i}$.
  For two Laurent series $F=\sum x_iw^i$ and $G=\sum y_jw^j$, the sum $\sum_i i x_iy_{-i}$ is $\Res_w F'(w)G(w)$. So the
  central term is $\beta\,\Res_w\,\frac{d}{dw}\big(L^m\big)\,L^n=\beta\,m\,\Res_w L^{m+n-1}\,dL$. If $m+n\ne0$, the
  integrand is $\frac{m}{m+n}\,d\big(L^{m+n}\big)$, the derivative of a Laurent series in $w$, whose residue is $0$. If
  $m+n=0$, it is $m\,dL/L=m\,\frac{dw}{(1+w)\,w\,u(w)}$, whose residue is $m/u(0)=m$. Hence the central term is
  $m\beta\,\delta_{m+n,0}$, as claimed.
\end{itemize}
\emph{The Virasoro and translation rules.} These are proved the same way, as residue identities for the coefficient
series. For $[L[m],a[n]]$ one needs the transformation of $\omega$ under $w=e^z-1$ (the Schwarzian shift
$-c/24$) and the primarity of $a$. For $(L_{-1}a)[n]$ one uses $Y(L_{-1}a,w)=\partial_wY(a,w)$ and integrates by parts,
which produces the two terms $-n\,a[n-1]$ and $-(P+1)\,a[n]$. In every case the needed coefficient identity comes down
to the one-line differential equation $(1+w)(u+wu')=1$, i.e.\ $\frac{d}{dw}\log(1+w)=\frac1{1+w}$. The proofs are in
\leanpath{ZhuCoordinateCoeff.lean}.
\end{proof}
\begin{physics}
Zhu's coordinates are the map from the plane to the cylinder $w=e^{2\pi i z}-1$. Currents transform as primaries of
weight one, so their modes in the new coordinate obey the same affine Kac--Moody algebra at the same level; the
stress tensor picks up the Schwarzian ($c/24$ shift), which is why $L[n]$ differs from $L_n$ but still acts on
currents as $-n$.
\end{physics}
\begin{status}
Standard axioms. Unconditional.
\end{status}

\begin{definition}[The weighted square-bracket $C_2$ span; the current $C_2$ span]\label{l25:def:spans}
$\mathrm{sqC2Span}(N)$ is spanned by $a[-2]b$ with $a\in W_p$, $p\ge1$, $b$ of $L_0$-degree $\le N-p-1$.
For a family of currents $g_i$, $\mathrm{currentC2}$ is spanned by $g_i[-n]y$ with $n\ge2$.
\end{definition}
\noindent Files \leanpath{VOA/\allowbreak Modularity/\allowbreak ZhuCoordinateGrading.lean}.
\begin{leancode}
def sqC2Span (N : ℤ) : Submodule ℂ 𝒱 :=
  Submodule.span ℂ {x | ∃ (p : ℕ) (a b : 𝒱), 1 ≤ p ∧ a ∈ W.grading (p : ℤ) ∧
    b ∈ gradeLE W (N - p - 1) ∧ x = vSqZ W a p (-2) b}

def currentC2 : Submodule ℂ 𝒱 :=
  Submodule.span ℂ {x | ∃ (i : ι) (n : ℕ) (y : 𝒱), 2 ≤ n ∧ x = vSqZ W (g i) 1 (-(n : ℤ)) y}
\end{leancode}
\begin{dictionary}
\lean{gradeLE W d} $=\bigoplus_{j\le d}W_j$ (\leanpath{ChangeCoord.lean}). The weight bound on $b$ makes every
$a[2m]b$ with $m\ge2$ vanish and $a[2]b\in W_0=\C\ltwentyfivevac$ (\lean{vSqZ_mem_gradeLE}), so every generator of
\lean{sqC2Span W 4} is a \lean{SqC2Term}. The translation rule shows $a[-n]b\in\mathrm{sqC2Span}$ for all $n\ge2$
(\lean{vSqZ_neg_mem_sqC2Span}).
\end{dictionary}
\begin{physics}
A state $g[-n]y$ with $n\ge2$ is a field with at least one derivative, $\partial^{n-1}g$, multiplying $y$. On
the torus its one-point function is a pure Eisenstein correction, with no $\ltwentyfiveEu$ term, because
\eqref{l25:eq:zhuminus2} has no derivative term. Quotienting by these states is passing to the ``classical'' algebra
$R_V$ (no derivatives) in the cylinder coordinate. The computation in which these spans appear is the reduction of
$(TT)$ to $C_2$ terms.
\end{physics}
\begin{status}
Definitions.
\end{status}

\begin{theorem}[$\kappa$ from a membership]\label{l25:thm:membership}
If $\ltwentyfivevac\ne0$ and $L[-2]^2\ltwentyfivevac\in\mathrm{sqC2Span}(4)$, then the vacuum quartic relation holds with
$\kappa=\frac54c^2+\frac{11}2c$.
\end{theorem}
\noindent File \leanpath{VOA/\allowbreak DeligneAudit/\allowbreak DeligneVacuumRelation.lean}.
\begin{leancode}
theorem vacuumQuarticRelation_of_mem_sqC2Span (hvac : W.vacuum ≠ 0)
    (hx : Lsq W (-2) (Lsq W (-2) W.vacuum) ∈ sqC2Span W 4) :
    VacuumQuarticRelation W (5 / 4 * W.centralCharge ^ 2 + 11 / 2 * W.centralCharge)
\end{leancode}
\begin{dictionary}
The membership needs no coefficients: \lean{exists_kappa_of_mem_sqC2Span} shows every element of
\lean{sqC2Span W 4} has one-point function $\kappa'G_4\star f$ for \emph{some} $\kappa'$, and
\lean{kappa_eq_of_vacuumQuarticRelation} pins $\kappa'$ by the indicial computation of
Theorem~\ref{l25:thm:kappa}.
\end{dictionary}
\begin{proof}[Proof on paper]
Linearity of $S$ and the computation for one \lean{SqC2Term} (proof of \lean{vacuumQuarticRelation_of_sqC2}) give
existence; the lowest coefficient fixes $\kappa$.
\end{proof}
\begin{physics}
If $(TT)$ is, in the cylinder coordinate, a sum of states of the form $\partial g\cdot y$ of total weight four,
then each of them contributes a multiple of $E_4$ times the partition function, by \eqref{l25:eq:zhuminus2}. The MLDE
follows, and its constant is then fixed by the vacuum's exponent, so no coefficient of the decomposition needs to be
computed.
\end{physics}
\begin{status}
Standard axioms. Unconditional as an implication.
\end{status}

\begin{definition}[Affine current data]\label{l25:def:sqcd}
\lean{SqCurrentData W ι} packages: primary weight-one currents $g_i$ ($i\in\iota$), structure constants $T_{ij}^c$
with $g_{i(0)}g_j=\sum_cT_{ij}^cg_c$, a form $F$ and level $k$ with $g_{i(1)}g_j=kF_{ij}\ltwentyfivevac$, an inverse form
$K$, a dual Coxeter number $h^\vee$ with $k+h^\vee\ne0$, the Sugawara formula
$\omega=\frac{1}{2(k+h^\vee)}\sum K_{ij}g_{i(-1)}g_j$, the trace condition $\sum K_{ij}T_{ij}^c=0$, and
$c(k+h^\vee)=k\sum K_{ij}F_{ij}$.
\end{definition}
\noindent File \leanpath{VOA/\allowbreak DeligneAudit/\allowbreak DeligneVacuumRelation.lean}.
\begin{leancode}
structure SqCurrentData (ι : Type*) [Fintype ι] where
  g : ι → 𝒱
  mem : ∀ i, g i ∈ W.grading 1
  primary : ∀ i, W.L 1 (g i) = 0
  T : ι → ι → ι → ℂ
  bracket : ∀ i j, W.smul (g i) 0 (g j) = ∑ c, T i j c • g c
  F : ι → ι → ℂ
  level : ℂ
  form : ∀ i j, W.smul (g i) 1 (g j) = (level * F i j) • W.vacuum
  K : ι → ι → ℂ
  hdual : ℂ
  ne : level + hdual ≠ 0
  sugawara : W.conformalVector
    = (2 * (level + hdual))⁻¹ • ∑ i, ∑ j, K i j • W.smul (g i) (-1) (g j)
  bracket_trace : ∀ c, ∑ i, ∑ j, K i j * T i j c = 0
  central : W.centralCharge * (level + hdual) = level * ∑ i, ∑ j, K i j * F i j
\end{leancode}
\begin{dictionary}
Every field is an identity checkable in an affine VOA; \lean{affineCurrentData D k hk}
(\leanpath{DeligneVacuumRelationAffine.lean}) constructs an instance for $L_k(\ltwentyfiveg)$ from any \lean{LieData}
(Chapter~\ref{ch:L23}), with $T$ the structure constants and $F$ the form table. The structure does not require
the $g_i$ to generate $W$; it is used only to compute $L[-2]\ltwentyfivevac$ and its descendants.
\end{dictionary}
\begin{physics}
A WZW-type current algebra with its Sugawara stress tensor. \lean{bracket_trace} is the statement that the
Casimir is invariant, so the Sugawara vector has no weight-one correction term.
\end{physics}
\begin{status}
Definition; instantiated for every simple $\ltwentyfiveg$.
\end{status}

\begin{theorem}[Square Sugawara; the symbol calculus]\label{l25:thm:symbols}
Given affine current data, $L[-2]\ltwentyfivevac=\frac{1}{2(k+h^\vee)}\sum K_{ij}g_i[-1]g_j[-1]\ltwentyfivevac$. In
$Q=V/\mathrm{currentC2}$ the operators $g_i[-1]$ commute, so each polynomial $p\in\C[X_i]$ acts; put
$\mathrm{sym}(p)=p(g[-1])\ltwentyfivevac\in Q$ and $\mathrm{Cas}=\sum K_{ij}X_iX_j$. If $\mathrm{sym}(X_i\,\mathrm{Cas})=0$
for every $i$, and $\ltwentyfivevac\ne0$, then the vacuum quartic relation holds with $\kappa=\frac54c^2+\frac{11}2c$.
Moreover $\mathrm{sym}$ vanishes on the ideal generated by polynomials with vanishing symbol, and
$\mathrm{sym}(\mathrm{ad}_b\,p)=g_b[0]\,\mathrm{sym}(p)$ for the coadjoint derivation $\mathrm{ad}_b$.
\end{theorem}
\noindent File \leanpath{VOA/\allowbreak DeligneAudit/\allowbreak DeligneVacuumRelation.lean}.
\begin{leancode}
theorem Lsq_neg_two_vacuum_eq :
    Lsq W (-2) W.vacuum = (2 * (D.level + D.hdual))⁻¹ •
      ∑ i, ∑ j, D.K i j • D.sq i (-1) (D.sq j (-1) W.vacuum)

noncomputable def sym (p : MvPolynomial ι ℂ) : 𝒱 ⧸ currentC2 W D.g :=
  D.ev p ((currentC2 W D.g).mkQ W.vacuum)

noncomputable def casPoly : MvPolynomial ι ℂ :=
  ∑ i, ∑ j, MvPolynomial.C (D.K i j) * (X i * X j)

theorem sym_eq_zero_of_mem_span {κ : Type*} (r : κ → MvPolynomial ι ℂ)
    (hr : ∀ e, D.sym (r e) = 0) {p : MvPolynomial ι ℂ} (hp : p ∈ Ideal.span (Set.range r)) :
    D.sym p = 0

theorem sym_adG (c : ι) (p : MvPolynomial ι ℂ) : D.sym (adG D.T c p) = D.Zq c (D.sym p)

theorem vacuumQuarticRelation_of_symbols (hvac : W.vacuum ≠ 0)
    (h : ∀ i, D.sym (X i * D.casPoly) = 0) :
    VacuumQuarticRelation W (5 / 4 * W.centralCharge ^ 2 + 11 / 2 * W.centralCharge)
\end{leancode}
\begin{dictionary}
\lean{D.sq i n} is $g_i[n]$. \lean{D.ev p} is the evaluation of a polynomial at the commuting operators
\lean{D.Tq i} $=g_i[-1]$ on the quotient (built by \lean{MvPolynomial.aeval} into the commutative subalgebra they
generate). \lean{D.Zq c} $=g_c[0]$ on the quotient. \lean{adG T b} is the derivation $X_i\mapsto\sum_cT_{bi}^cX_c$ of
$\C[X]$. \lean{X i * D.casPoly} is $X_i\cdot\mathrm{Cas}$. The key hypothesis says the cubic polynomials
$X_i\mathrm{Cas}$ act trivially on the vacuum modulo $\mathrm{currentC2}$. \textbf{Section variables.} In
\leanpath{DeligneVacuumRelation.lean} the VOA \lean{W : VertexOperatorAlgebra 𝒱} is implicit here and
\lean{D : SqCurrentData W ι} (with \lean{[Fintype ι]}) is an explicit section variable: the ambient square-current data,
which is why the declarations are written \lean{D.sym}, \lean{D.casPoly}. Its fields \lean{D.level}, \lean{D.hdual},
\lean{D.K} are $k$, $h^\vee$ and the inverse form $K_{ij}$.
\end{dictionary}
\begin{proof}[Proof on paper]
\emph{Square Sugawara.} Expand $g_i[-1]g_j[-1]\ltwentyfivevac=g_{i(-1)}g_j+\frac12g_{i(0)}g_j-\frac1{12}g_{i(1)}g_j$,
using the coefficients $z(1,-1,0)=\frac12$ and $z(1,-1,1)=-\frac1{12}$, read off from
$u^{-1}=1+\frac w2-\frac{w^2}{12}+\dots$. Then use \lean{bracket_trace}, which kills the $g_{i(0)}g_j$ sum, and \lean{central}, which
turns the $g_{i(1)}g_j$ sum into the $c/24$ shift between $\omega$ and $L[-2]\ltwentyfivevac$.

\emph{Stability of $\mathrm{currentC2}$ (needed for everything that follows).} For $m\le0$ the operator $g_i[m]$ maps
$\mathrm{currentC2}$ into itself (\lean{sq_mem_currentC2}). Take a generator $g_c[-n]y$ with $n\ge2$. By
Theorem~\ref{l25:thm:sqcomm},
\begin{equation*}
g_i[m]\,g_c[-n]y=g_c[-n]\,\big(g_i[m]y\big)+\sum_eT_{ic}^e\,g_e[m-n]y ,
\end{equation*}
and there is no central term, because $m-n\le-2\ne0$. The first term is a generator, and so is each $g_e[m-n]y$, since
$m-n\le-2$. So $g_i[-1]$ and $g_i[0]$ descend to operators \lean{D.Tq i}, \lean{D.Zq i} on $Q=V/\mathrm{currentC2}$.
The $g_i[-1]$ commute on $Q$, since their commutator $(g_{i(0)}g_j)[-2]$ maps into $\mathrm{currentC2}$. So
$\mathrm{sym}$ is well defined, and $[g_c[0],g_i[-1]]=(g_{c(0)}g_i)[-1]$ gives \lean{sym_adG}. The span property is
linearity.

\emph{The reduction.} Put $\omega'=L[-2]\ltwentyfivevac=s\sum K_{ij}g_i[-1]g_j[-1]\ltwentyfivevac$, with
$s=(2(k+h^\vee))^{-1}$. The rule $[L[m],g[n]]=-n\,g[m+n]$ at $m=-2$, $n=-1$ reads
$L[-2]\,g_i[-1]=g_i[-1]\,L[-2]+g_i[-3]$ (\lean{Lsq_neg_two_sq}). Applying it twice,
\begin{equation*}
L[-2]^2\ltwentyfivevac=L[-2]\,\omega'=s\sum_{i,j}K_{ij}\Big(g_i[-1]\big(g_j[-1]\,\omega'+g_j[-3]\ltwentyfivevac\big)
+g_i[-3]\,g_j[-1]\ltwentyfivevac\Big).
\end{equation*}
Here $g_i[-3]g_j[-1]\ltwentyfivevac$ is a generator of $\mathrm{currentC2}$. Also $g_j[-3]\ltwentyfivevac$ is a
generator, and $g_i[-1]$ preserves $\mathrm{currentC2}$ by stability, so $g_i[-1]g_j[-3]\ltwentyfivevac$ lies in it too.
It remains to show $g_j[-1]\,\omega'\in\mathrm{currentC2}$ for every $j$. In $Q$, $\omega'$ has class
$s\,\mathrm{sym}(\mathrm{Cas})$, so $g_j[-1]\omega'$ has class $s\,\mathrm{sym}(X_j\mathrm{Cas})$, which is $0$ by
hypothesis. Hence $L[-2]^2\ltwentyfivevac\in\mathrm{currentC2}$. Each term has square-bracket weight four, so it lies in
the weight-four part, which is contained in $\mathrm{sqC2Span}(4)$ (\lean{currentC2At_le_sqC2Span}).
Theorem~\ref{l25:thm:membership} concludes.
\end{proof}
\begin{physics}
This is the statement ``the Higgs branch relation $J^a\cdot\mathrm{Cas}=0$ holds in $R_V$'' lifted to Zhu's
coordinates. \emph{Heuristic, not Lean:} in the Deligne series the Higgs branch is expected to be the minimal
nilpotent orbit closure (for types other than $A$, its ideal is the associated graded of the Joseph ideal), and
$J^a\,\mathrm{Cas}$ vanishes on it; in particular $\mathrm{Cas}^2=0$ in $R_V$ while $\mathrm{Cas}\ne0$ ($R_V$ is
non-reduced). Nilpotency of order two of the stress tensor's symbol is what makes the MLDE second order.
\end{physics}
\begin{status}
Standard axioms. Unconditional as an implication.
\end{status}

\begin{warning}[title={Dead end: Zhu's change of coordinates as a VOA isomorphism}]
A first attempt tried to transport ``$\omega_{(-1)}\omega\in C_2(V)$'' through Zhu's isomorphism
$(V,Y)\cong(V,Y[\,])$. The repository's \lean{changeCoord} has no transformation law and there is no square Borcherds
identity; formalising them is a large project. The symbol calculus above avoids it entirely. A Lean trap recorded
on the way: in $L_{-3/2}(\mathfrak{sl}_3)$, \lean{rw} with equations between quotient expressions times out in
\lean{isDefEq}; proving the combination as one \lean{mkQ} of a vector (\lean{sym_three}) fixes it.
\end{warning}

%======================================================================
\section{The vacuum quartic relation for every Deligne type}\label{l25:sec:alltypes}
%======================================================================

\begin{theorem}[From carrier data to the target bound]\label{l25:thm:carrierdata}
Suppose given the following data and hypotheses.
\begin{enumerate}
\item \emph{Carrier data} $CD$ (Chapter~\ref{ch:L24}): a number $n$ of currents (for the intended use,
  $n=\dim\ltwentyfiveg$); quadrics $\mathrm{rel}_r\in\operatorname{Sym}^2$ spanning $I_2^{\rm ns}$; a Casimir quadric $Q$;
  and a certificate of U3, i.e.\ $X_aQ\in(\mathrm{rel}_r)_r$ for every $a$.
\item A Lie algebra record $D$ (a \lean{LieData}, Chapter~\ref{ch:L23}) on the index set $\{0,\dots,n-1\}$, with dual
  Coxeter number $D.h$ in the normalisation of its invariant form.
\item A level $k\in\C$ with $k+D.h\ne0$, so that the Sugawara construction and $L_k=$\lean{simpleVOA D k hk} exist.
\item \lean{hcas}: a number $c_0\in\C$ with $\mathrm{Cas}=c_0\,Q$. Here $\mathrm{Cas}=\sum K_{ij}X_iX_j$ is the Sugawara
  quadric of $D$, and $Q$ is the certificate's Casimir.
\item \lean{hrel}: every relation $\mathrm{rel}_r$ has vanishing symbol in $L_k$ (Theorem~\ref{l25:thm:symbols}).
\end{enumerate}
Then the vacuum quartic relation holds for $L_k$ with $\kappa=\frac54c^2+\frac{11}2c$. If moreover
\begin{enumerate}\setcounter{enumi}{5}
\item \lean{hc}: $h\in\N$ satisfies $nk/(k+D.h)=-2(h+1)$, i.e.\ the central charge is $-2(h+1)$,
\end{enumerate}
then $\dim L_k(\ltwentyfiveg)_N\ge c_N=\mathtt{mldeCoeff}\ h\ N$ for all $N$.
\end{theorem}
\noindent File \leanpath{VOA/\allowbreak DeligneAudit/\allowbreak DeligneVacuumRelationCarrier.lean}.
\begin{leancode}
theorem vacuumQuarticRelation_of_carrierData (CD : DeligneCarrierData) (D : LieData κ (Fin CD.n))
    (k : ℂ) (hk : k + D.h ≠ 0) (c₀ : ℂ)
    (hcas : (affineCurrentData D k hk).casPoly = MvPolynomial.C c₀ * CD.Q)
    (hrel : ∀ r, (affineCurrentData D k hk).sym (quadraticRelation (CD.rel r)) = 0) :
    VacuumQuarticRelation (simpleVOA D k hk)
      (5 / 4 * (simpleVOA D k hk).centralCharge ^ 2 + 11 / 2 * (simpleVOA D k hk).centralCharge) :=

theorem targetLowerBound_of_carrierData (CD : DeligneCarrierData) (D : LieData κ (Fin CD.n))
    (k : ℂ) (hk : k + D.h ≠ 0) (h : ℕ)
    (hc : (CD.n : ℂ) * k / (k + D.h) = -2 * ((h : ℂ) + 1)) (c₀ : ℂ)
    (hcas : (affineCurrentData D k hk).casPoly = MvPolynomial.C c₀ * CD.Q)
    (hrel : ∀ r, (affineCurrentData D k hk).sym (quadraticRelation (CD.rel r)) = 0) :
    DeligneCharacter.TargetLowerBoundQ (simpleVOA D k hk) (DeligneCharacter.mldeCoeff h)
\end{leancode}
\begin{dictionary}
\lean{CD.Q} $=$ \lean{quadraticRelation CD.cas}, the Casimir as a quadratic polynomial; \lean{quadraticRelation q} is
$\sum_{x,y}q(x,y)X_xX_y$. \lean{simpleVOA D k hk} is $L_k(\ltwentyfiveg)$ for the Lie algebra encoded by $D$
(Chapter~\ref{ch:L23}). \lean{hc} is the central charge condition ($c=k\dim\ltwentyfiveg/(k+h^\vee)$ by
\lean{simple_centralCharge'}). \lean{hcas} identifies the Sugawara quadric with the certificate Casimir (needed because
U3 is stated for the latter). \lean{hrel} is the vanishing of the relation symbols. \textbf{$c_0\ne0$ is not required}:
if $c_0=0$, U3 is not even used. The conclusion's $W$ is the concrete simple VOA. \textbf{Notation clash.} The
\lean{κ} in \lean{LieData κ (Fin CD.n)} is a bilinear form \lean{κ : LinearMap.BilinForm ℂ 𝔤} on a finite-dimensional
Lie algebra \lean{𝔤} (implicit section variables, with \lean{[FiniteDimensional ℂ 𝔤]}), not the constant $\kappa$ of
Definition~\ref{l25:def:vqr}. \lean{D.h} is the dual Coxeter number \emph{in the normalisation of that form} (the
adjoint Casimir acts by $2\,$\lean{D.h}), and $k$ is the level relative to the same form; the concrete records of
Chapter~\ref{ch:L23} are normalised so that \lean{D.h} is the usual $h^\vee$. The theorem assumes neither that
$\ltwentyfiveg$ is simple nor that \lean{D.h} equals the ambient $h$: the two are related only through \lean{hc}.
\lean{simpleVOA D k hk} is the simple quotient of the universal affine VOA for whatever $D$ encodes. The first
theorem is the same statement stopped one step earlier, before \lean{hc} is used; it is the input to the equality
route recorded in the warning after Theorem~\ref{l25:thm:alltypes}.
\end{dictionary}
\begin{proof}[Proof on paper]
By U3 and \lean{hcas}, $X_a\mathrm{Cas}=c_0X_aQ$ lies in the ideal generated by the relations, whose symbols vanish;
\lean{sym_eq_zero_of_mem_span} gives $\mathrm{sym}(X_a\mathrm{Cas})=0$; Theorem~\ref{l25:thm:symbols} gives the vacuum
quartic relation; Theorem~\ref{l25:thm:kappa} the bound. How \lean{hrel} is proved per type: the symbol of the
highest-weight relation vanishes because the weight-two singular vector dies in $L_k$ and two finite ``shadow''
checks hold (\lean{sym_singular}: the bracket and form contractions of the relation vanish). All other relations are
obtained from it by the coadjoint action, and vanishing \emph{propagates} as follows. Since
$\mathrm{sym}(\mathrm{ad}_b R)=g_b[0]\,\mathrm{sym}(R)$ (Theorem~\ref{l25:thm:symbols}), if $\mathrm{sym}(R)=0$ and
$\lambda\,\mathrm{ad}_bR=\mu\,R'$ with $\mu\ne0$, then $\mathrm{sym}(R')=0$. Think of the relations as the vertices of a
graph, with an edge $R\to R'$ for each such identity, labelled by $(b,\lambda,\mu)$. A breadth-first search from the
highest-weight relation that reaches every vertex is then a proof that all symbols vanish (\lean{sym_step}). The
certificates of Section~\ref{l25:sec:data} are such searches: 102 steps for $D_4$, the 1539 relations of $E_7$, and
3874 steps for the 3875 relations of $E_8$. Each step is an identity of explicit quadrics, checked by
\lean{decide +kernel}.
\end{proof}
\begin{physics}
\emph{Heuristic:} the quadrics of $I_2^{\rm ns}$ are the Higgs branch chiral-ring relations; they hold in $R_V$ because the weight-two singular vector
(the 4d ``Higgs branch chiral ring relation'') is null in the simple quotient. The theorem says these classical
relations, plus invariance of the Casimir, already force the full vacuum character.
\end{physics}
\begin{status}
Standard axioms (re-run for this chapter). Unconditional as an
implication.
\end{status}

\begin{theorem}[The character bound for every Deligne type]\label{l25:thm:alltypes}
For each row of the table, the simple affine VOA $W=L_k(\ltwentyfiveg)$ satisfies, with no hypothesis,
$\dim W_N\ge c_N=\mathtt{mldeCoeff}\ h^\vee\ N$ for all $N$ (\lean{TargetLowerBoundQ W (mldeCoeff h∨)}). The last column
says what else is stated as a Lean theorem.
\begin{center}\small
\begin{tabular}{p{0.05\textwidth}p{0.06\textwidth}p{0.04\textwidth}p{0.05\textwidth}p{0.27\textwidth}p{0.38\textwidth}}
\toprule
type & $k$ & $h^\vee$ & $\kappa$ & bound $\dim W_N\ge c_N$ & also stated\\
\midrule
$A_2$ & $-3/2$ & 3 & 36 & via \lean{finrank_eq_mldeCoeff'} & the equality $\dim W_N=c_N$; \lean{vacuumQuarticRelation} ($\kappa=36$)\\
$G_2$ & $-5/3$ & 4 & 70 & \lean{g2_targetLowerBound} & \lean{g2_vacuumQuarticRelation} (as $\frac54c^2+\frac{11}2c$); $c$ and $\ltwentyfivevac\ne0$\\
$D_4$ & $-2$ & 6 & 168 & \lean{targetLowerBoundQ4} & \lean{vacuumQuarticRelation4} ($\kappa=168$)\\
$F_4$ & $-5/2$ & 9 & 390 & \lean{f4_targetLowerBound} & \lean{f4_vacuumQuarticRelation} (as $\frac54c^2+\frac{11}2c$); $c$ and $\ltwentyfivevac\ne0$\\
$E_6$ & $-3$ & 12 & 702 & \lean{e6_targetLowerBound} & only the inputs: \lean{e6_hc}, \lean{e6_casimir_match}, \lean{e6_hrel}\\
$E_7$ & $-4$ & 18 & 1596 & \lean{e7_targetLowerBound} & only the inputs: \lean{e7_hc}, \lean{e7_casimir_match}, \lean{e7_hrel}\\
$E_8$ & $-6$ & 30 & 4464 & \lean{targetLowerBoundQ8} & \lean{vacuumQuarticRelation8} ($\kappa=4464$)\\
\bottomrule
\end{tabular}
\end{center}
In every row $\kappa=(h^\vee+1)(5h^\vee-6)$, as Theorem~\ref{l25:thm:kappa} forces.
\end{theorem}
\noindent Files \leanpath{VOA/\allowbreak DeligneAudit/\allowbreak DeligneVacuumRelationA2.lean},
\leanpath{DeligneVacuumRelationD4.lean}, \leanpath{E6Vacuum.lean}, \leanpath{E7LieVacuum.lean}, \leanpath{E7LieZhu.lean},
\leanpath{DeligneVacuumRelationE8.lean}, \leanpath{G2LieVacuum.lean}, \leanpath{F4LieVacuum.lean} (all in
\leanpath{VOA/\allowbreak DeligneAudit/}). Proof bodies are omitted: each block ends at \lean{:=} or \lean{:= by}. The
line above each block names the type and its \lean{W}.

\noindent$A_2$, namespace \lean{VOA.DeligneAudit.A2}, \lean{W} $=L_{-3/2}(\mathfrak{sl}_3)$:
\begin{leancode}
theorem vacuumQuarticRelation : VacuumQuarticRelation W 36 := by

theorem finrank_eq_mldeCoeff' (N : ℕ) :
    (Module.finrank ℂ (W.grading N) : ℚ) = DeligneCharacter.mldeCoeff 3 N :=
\end{leancode}
\noindent$D_4$, namespace \lean{VOA.DeligneAudit.D4Lie}, \lean{W4 := simpleVOA d4CertData (-2) hk4}:
\begin{leancode}
theorem vacuumQuarticRelation4 : VacuumQuarticRelation W4 168 := by

theorem targetLowerBoundQ4 : DeligneCharacter.TargetLowerBoundQ W4 (DeligneCharacter.mldeCoeff 6) :=
\end{leancode}
\noindent$E_6$, namespace \lean{VOA.DeligneAudit.DeligneCarrier.E6}, \lean{W6 := simpleVOA e6LieData k6 hk6} with
$k_6=-3$; \lean{E6cd} is its \lean{affineCurrentData}:
\begin{leancode}
theorem e6_hc : ((E6data.n : ℕ) : ℂ) * k6 / (k6 + e6LieData.h) = -2 * (((12 : ℕ) : ℂ) + 1) := by

theorem e6_casimir_match : E6cd.casPoly = MvPolynomial.C (2 / 3 : ℂ) * E6data.Q :=

theorem e6_hrel (r : Fin 650) : E6cd.sym (quadraticRelation (E6data.rel r)) = 0 := by

theorem e6_targetLowerBound : DeligneCharacter.TargetLowerBoundQ W6 (DeligneCharacter.mldeCoeff 12) :=
\end{leancode}
\noindent$E_7$, namespace \lean{VOA.DeligneAudit.E7Lie}, \lean{W7 := simpleVOA e7CertData (-4) hk7}:
\begin{leancode}
theorem e7_hc : ((DeligneCarrier.E7.E7data.n : ℕ) : ℂ) * (-4) / (-4 + e7CertData.h) =
    -2 * (((18 : ℕ) : ℂ) + 1) := by

theorem e7_casimir_match :
    (VacuumRelation.affineCurrentData (ι := Fin DeligneCarrier.E7.E7data.n) e7CertData (-4)
      hk7).casPoly = MvPolynomial.C (1 / 2 : ℂ) * DeligneCarrier.E7.E7data.Q :=

theorem e7_hrel (r' : Fin 1539) :
    (affineCurrentData (ι := Fin DeligneCarrier.E7.E7data.n) e7CertData (-4) hk7).sym
      (quadraticRelation (DeligneCarrier.E7.E7data.rel r')) = 0 :=

theorem e7_targetLowerBound : DeligneCharacter.TargetLowerBoundQ W7 (DeligneCharacter.mldeCoeff 18) :=
\end{leancode}
\noindent$E_8$, namespace \lean{VOA.DeligneAudit.E8Lie}, \lean{W8 := simpleVOA e8CertData (-6) hk8}:
\begin{leancode}
theorem targetLowerBoundQ8 :
    DeligneCharacter.TargetLowerBoundQ W8 (DeligneCharacter.mldeCoeff 30) :=

theorem vacuumQuarticRelation8 : VacuumQuarticRelation W8 4464 := by
\end{leancode}
\noindent$G_2$, namespace \lean{VOA.DeligneAudit.G2Lie}, \lean{g2W := simpleVOA g2CertData (-5/3) g2_level_ne}:
\begin{leancode}
theorem g2_vacuumQuarticRelation :
    DeligneMLDE.VacuumQuarticRelation g2W
      (5 / 4 * g2W.centralCharge ^ 2 + 11 / 2 * g2W.centralCharge) :=

theorem g2W_centralCharge : g2W.centralCharge = -2 * (((4 : ℕ) : ℂ) + 1) := by

theorem g2W_vacuum_ne_zero : g2W.vacuum ≠ 0 := simple_vacuum_ne_zero' _ _ _

theorem g2_targetLowerBound :
    DeligneCharacter.TargetLowerBoundQ g2W (DeligneCharacter.mldeCoeff 4) :=
\end{leancode}
\noindent$F_4$, namespace \lean{VOA.DeligneAudit.F4Lie}, \lean{f4W := simpleVOA f4CertData (-5/2) f4_level_ne}:
\begin{leancode}
theorem f4_vacuumQuarticRelation :
    DeligneMLDE.VacuumQuarticRelation f4W
      (5 / 4 * f4W.centralCharge ^ 2 + 11 / 2 * f4W.centralCharge) :=

theorem f4W_centralCharge : f4W.centralCharge = -2 * (((9 : ℕ) : ℂ) + 1) := by

theorem f4W_vacuum_ne_zero : f4W.vacuum ≠ 0 := simple_vacuum_ne_zero' _ _ _

theorem f4_targetLowerBound :
    DeligneCharacter.TargetLowerBoundQ f4W (DeligneCharacter.mldeCoeff 9) :=
\end{leancode}
\begin{dictionary}
The \lean{LieData} records are constructed and checked in Chapter~\ref{ch:L23}. \textbf{All these theorems have no
hypotheses.} \textbf{Same as, stronger or weaker?} For $A_2$ the Lean statement is the mathematical equality. For the
other six types the exported conclusion is \lean{TargetLowerBoundQ}, which by definition is the one-sided
$\forall N,\ c_N\le\dim W_N$. That is \emph{weaker} than the equality, and Theorem~\ref{l25:thm:alltypes} claims only
that (Caveat~\ref{l25:cav:equality}). The proofs pass through \lean{targetLowerBoundQ_of_isAKSolution}, which discards
the other half of an equality it has in hand (Theorem~\ref{l25:thm:fromAK}). The numbers: $36=4\cdot9$, $70=5\cdot14$,
$168=7\cdot24$, $390=10\cdot39$, $702=13\cdot54$, $1596=19\cdot84$, $4464=31\cdot144$, each $(h+1)(5h-6)$. For $G_2$ and
$F_4$ the library states $\kappa$ as $\frac54c^2+\frac{11}2c$, which at $c=-10,-20$ gives $70$ and $390$.

\textbf{What is known about the Lie algebras.} For each record, Chapter~\ref{ch:L23} proves in Lean that the
algebra is a simple Lie algebra of the right dimension, with the invariant form normalised so that \lean{D.h} is the
usual $h^\vee$ (Theorem~\ref{l23:thm:g2f4e6} and the lattice theorems there). It does \emph{not} prove that, say,
\lean{e6LieData} is isomorphic to the Bourbaki $\mathfrak e_6$, because the Cartan--Killing classification is not
formalised. For dimensions $14,28,52,133,248$ the identification follows on paper from the classification and the
dimension. Dimension $78$ is ambiguous ($B_6$, $C_6$, $E_6$), and for the structure-constant $\mathfrak e_6$ used here
the ambiguity is settled on paper by the value of the Killing form on a long coroot (Remark~\ref{l23:rem:type}). The
library has three Lean models of $\mathfrak e_6$, and no isomorphism between them is formalised; the $E_6$ theorem is
about \lean{e6LieData}. So each theorem is literally about $L_k$ of the algebra the record encodes, and that algebra is
the intended one by an argument on paper.\end{dictionary}
\begin{proof}[Proof on paper]
Theorem~\ref{l25:thm:carrierdata} with each type's carrier data. The work is in \lean{hrel}: one singular-vector
computation per type and a propagation search through all relations of $I_2^{\rm ns}$, as described in the proof of
Theorem~\ref{l25:thm:carrierdata} (certificates in Section~\ref{l25:sec:data}). The Casimir match \lean{hcas} is checked
entrywise (\lean{quadraticRelation_eq_of_symm}). Then the central charge is computed from \lean{hc}, and
Theorem~\ref{l25:thm:kappa} applies.
\end{proof}
\begin{physics}
This is the vacuum character of the rank-one Minahan--Nemeschansky theories ($E_6,E_7,E_8$), of $SU(2)$ with four
flavours ($D_4$) and of the $(A_1,D_4)$ Argyres--Douglas theory ($A_2$): their Schur indices satisfy the second-order
MLDE \eqref{l25:eq:mlde}, as Beem--Rastelli predicted and Arakawa--Kawasetsu proved. $G_2$ at $k=-5/3$ and $F_4$ at
$-5/2$ have no known 4d origin, but the VOA statement is the same. What the library contains, as Lean theorems: the
equality $\dim L_N=c_N$ for $A_2$, and the lower bound $\dim L_N\ge c_N$ (the half that classical freeness needs) for
all seven types (Caveat~\ref{l25:cav:equality}). For $A_2$ the further equality with the product $\prod(1+q^r)^8$ is
derivable but not stated (Remark~\ref{l25:rem:a2cons}).
\end{physics}
\begin{status}
\textbf{Verdict.} Every theorem in the blocks above is a Lean theorem with no hypotheses and standard axioms, and the
$\ge$ bound is unconditional for all seven types; the equality only for $A_2$. \textbf{Machine evidence}
(Section~\ref{l25:sec:machine}). These files were built green in the \lean{VOA.Deligne} umbrella build at
an earlier commit; the quoted statements are unchanged in the state this chapter was written against, and a pre-release aggregate build, which contains the umbrella,
was green on 2026-10-28 (18{,}901 jobs; 18{,}907 after a merge). The axiom lists are the \#\lean{print axioms} commands in the files, as run in that build
and in the later targeted builds. For this chapter we re-ran \#\lean{print axioms} on the generic lemmas
\lean{targetLowerBound_of_carrierData}, \lean{vacuumQuarticRelation_of_carrierData},
\lean{isAKSolution_of_vacuumQuarticRelation} and \lean{finrank_eq_mldeCoeff_of_isAKSolution} (standard), and we did not
import the $E_6$, $E_7$, $E_8$ modules.
\end{status}

\begin{warning}[title={LEAN DEFECT: equality is not a Lean theorem for six types}]
The equality $\dim L_k(\ltwentyfiveg)_N=c_N$ is \textbf{not stated} in the library for $G_2$, $D_4$, $F_4$, $E_6$, $E_7$,
$E_8$; only $A_2$ has it (\lean{A2.finrank_eq_mldeCoeff'}). It is derivable from exported lemmas, in one line per type:
\begin{itemize}
\item $G_2$, $F_4$: \lean{finrank_eq_mldeCoeff_of_isAKSolution _ (isAKSolution_of_vacuumQuarticRelation _ h
  xW_centralCharge xW_vacuum_ne_zero x_vacuumQuarticRelation) N} with $h=4$, $9$.
\item $D_4$, $E_8$: the same with \lean{vacuumQuarticRelation4} / \lean{vacuumQuarticRelation8}, $h=6$ / $30$; the
  central-charge hypothesis is \lean{simple_centralCharge'} plus \lean{d4CertData_h} / \lean{e8CertData_h}, and
  \lean{hvac} is \lean{simple_vacuum_ne_zero'}.
\item $E_6$, $E_7$: no \lean{VacuumQuarticRelation} theorem is exported, so first apply
  \lean{vacuumQuarticRelation_of_carrierData} to \lean{(E6data, e6LieData, k6, hk6, 2/3, e6_casimir_match, e6_hrel)}
  and \lean{(E7data, e7CertData, -4, hk7, 1/2, e7_casimir_match, e7_hrel)}, then proceed as above with
  \lean{e6_hc}/\lean{e7_hc} turned into the central-charge hypothesis by \lean{simple_centralCharge'}.
\end{itemize}
\emph{What the machine confirmed for this chapter:} the $G_2$, $F_4$ and $D_4$ compositions typecheck in a scratch
file (\lean{lake env lean -DmaxSynthPendingDepth=3}, no errors), namely
\begin{plaincode}
example (N : ℕ) : (Module.finrank ℂ (G2Lie.g2W.grading N) : ℚ) = DeligneCharacter.mldeCoeff 4 N :=
  DeligneMLDE.finrank_eq_mldeCoeff_of_isAKSolution _
    (DeligneMLDE.isAKSolution_of_vacuumQuarticRelation _ 4 G2Lie.g2W_centralCharge
      G2Lie.g2W_vacuum_ne_zero G2Lie.g2_vacuumQuarticRelation) N
\end{plaincode}
and its $F_4$ and $D_4$ analogues. The $E_6$, $E_7$, $E_8$ compositions were \emph{not} checked (heavy modules not
imported). The scratch file is not part of the library. Recommended fix: a generic
\lean{finrank_eq_mldeCoeff_of_carrierData} beside \lean{targetLowerBound_of_carrierData}, and per-type
\lean{*_finrank_eq_mldeCoeff} corollaries. Until then, cite the equality for these six types as ``derivable, not
stated''. This is the authoritative statement of Caveat~\ref{l25:cav:equality}; elsewhere the chapter refers here.
\end{warning}

\begin{remark}[The $A_2$ consequences]\label{l25:rem:a2cons}
What exists in Lean for $A_2$, exactly:
\begin{itemize}
\item \lean{A2.finrank_eq_mldeCoeff'}: $\dim V_N=c_N$ (with \lean{mldeCoeff 3});
\item \lean{A2.productCoeff_eq_mldeCoeff'}: $[q^N]\prod_{r\ge1}(1+q^r)^8=c_N$ (from \lean{A2.a2ProductMLDE},
  Theorem~\ref{l25:thm:a2prod});
\item \lean{A2.kacWakimotoLowerBound_of_quartic}: the Kac--Wakimoto lower half $[q^N]\prod(1+q^r)^8\le\dim V_N$
  (structure \lean{KacWakimotoLowerBound}), from any vacuum quartic relation and \lean{a2ProductMLDE}, with no
  literature input;
\item \lean{A2.finrank_grading_two_ge}: $36\le\dim V_2$, by an explicit $36\times36$ Gram determinant at $k=-3/2$
  (\leanpath{A2CharacterGram.lean}), an independent low-weight check.
\end{itemize}
The \emph{character formula} $\dim V_N=[q^N]\prod_{r\ge1}(1+q^r)^8$ itself is \textbf{not stated} as one Lean theorem.
It follows in one line from the first two items, and the proof of \lean{kacWakimotoLowerBound_of_isAKSolution} even
derives it internally before weakening it to $\le$. It is a further instance of ``derivable, not stated'', recorded in
Section~\ref{l25:sec:defects}.
\end{remark}

%======================================================================
\section{Formal \texorpdfstring{$q$}{q}-series: the toolkit and the classical identities}\label{l25:sec:qseries}
%======================================================================

All identities in this section are equalities in a formal power series ring $A[[q]]$ (Mathlib's
\lean{PowerSeries A}), proved from scratch. Nothing analytic is used except one appeal to Mathlib's level-one
Sturm bound (Theorem~\ref{l25:thm:ramanujan}). The infrastructure lives in
\leanpath{VOA/\allowbreak Commutative/\allowbreak JacobiTripleCore.lean}. We use the $q$-Pochhammer symbols
\begin{equation*}
(a;q)_n=\prod_{i=0}^{n-1}(1-aq^i),\qquad (a;q)_\infty=\prod_{i\ge0}(1-aq^i),\qquad (q)_n=(q;q)_n,
\end{equation*}
and $\binom nk_q=(q)_n/((q)_k(q)_{n-k})$ for the Gaussian binomial.

\paragraph{Which identity is used where.}
\begin{center}\small
\begin{tabular}{p{0.3\textwidth}p{0.6\textwidth}}
\toprule
identity & used for\\
\midrule
Jacobi triple product & everything below: Ramanujan (via the heat equation), Legendre (via Gauss), Lorenz\\
pentagonal, quintuple product & side results; the quintuple product feeds Lorenz\\
Legendre & the last step of the $A_2$ star identity (critical)\\
Ramanujan & the $A_2$ product (critical), the closed forms of Section~\ref{l25:sec:closed} ($D_4$ critical), $G_2$\\
$A_2$ product & $c_N$ for $A_2$ as a product; Kac--Wakimoto lower half (critical for $A_2$)\\
Lorenz, Borwein, $G_2$ target & the target side of the open $G_2$ star identity (side result)\\
\bottomrule
\end{tabular}
\end{center}

\begin{definition}[Agreement, the Euler operator, infinite products, sums over $\Z$]\label{l25:def:toolkit}
$f$ and $g$ \emph{agree to order $K$} if $q^K\mid f-g$. The Euler operator is $\ltwentyfiveEu f=\sum_nn\,a_nq^n$. For factors
$u_k$ with $u_k\equiv1\pmod{q^{k+1}}$, the infinite product $\prod_{k\ge0}u_k$ has $q^d$-coefficient that of the
finite product $\prod_{k<d}u_k$. For $c:\Z\to A$, $e:\Z\to\N$ with finite fibres inside $[-d-1,d+1]$,
$\sum_{n\in\Z}c(n)q^{e(n)}$ has $q^d$-coefficient $\sum_{e(n)=d}c(n)$.
\end{definition}
\noindent File \leanpath{VOA/\allowbreak Commutative/\allowbreak JacobiTripleCore.lean}.
\begin{leancode}
def AgreeTo (K : ℕ) (f g : PowerSeries R) : Prop := (X : PowerSeries R) ^ K ∣ f - g

def eulerOp (f : PowerSeries R) : PowerSeries R := mk fun n => (n : R) * coeff n f

def ProdFactors (u : ℕ → PowerSeries R) : Prop := ∀ k, AgreeTo (k + 1) (u k) 1

def infProd (u : ℕ → PowerSeries R) : PowerSeries R :=
  mk fun d => coeff d (∏ k ∈ Finset.range d, u k)

def zSeries (c : ℤ → R) (e : ℤ → ℕ) : PowerSeries R :=
  mk fun d => ∑ n ∈ Finset.Icc (-(d : ℤ) - 1) (d + 1) with e n = d, c n
\end{leancode}
\begin{dictionary}
\textbf{Junk values.} \lean{infProd u} is defined for \emph{every} sequence, but it is the infinite product only
when \lean{ProdFactors u} holds; every theorem that uses it carries or proves that condition
(e.g.\ \lean{prodFactors_tripleFactor}). Likewise \lean{zSeries c e} truncates the sum to $|n|\le d+1$ and is
the intended sum only when every $n$ with $e(n)=d$ satisfies this (\lean{natAbs_le_pentExp} etc.). The Euler
product $(q;q)_\infty$ is Mathlib-style \lean{euler} of \leanpath{QPochhammer.lean} (Chapter~\ref{ch:L14}), with
inverse \lean{reciprocalEuler}.
\end{dictionary}
\begin{physics}
Formal $q$-series are partition functions counted state by state; infinite products are free-field (Fock space)
partition functions, sums over $\Z$ are lattice (momentum) sums. Every identity below is a boson--fermion or
product--sum identity of 2d CFT.
\end{physics}
\begin{status}
Definitions.
\end{status}

\begin{theorem}[Jacobi triple product]\label{l25:thm:jtp}
For any commutative ring $A$ and unit $u\in A^\times$, in $A[[q]]$:
\begin{equation*}
(1-u)\prod_{m\ge1}(1-uq^m)(1-u^{-1}q^m)(1-q^m)=\sum_{j\in\Z}(-1)^ju^jq^{j(j-1)/2}.
\end{equation*}
\end{theorem}
\noindent File \leanpath{VOA/\allowbreak Commutative/\allowbreak JacobiTriple.lean}.
\begin{leancode}
def pentExp (j : ℤ) : ℕ := (j * (j - 1) / 2).toNat

def tripleCoeff (u : Aˣ) (j : ℤ) : A := negSign j * ((u ^ j : Aˣ) : A)

def tripleSum (u : Aˣ) : PowerSeries A := zSeries (tripleCoeff u) pentExp

def tripleFactor (u : Aˣ) (m : ℕ) : PowerSeries A :=
  (1 - C (u : A) * X ^ (m + 1)) * (1 - C ((u⁻¹ : Aˣ) : A) * X ^ (m + 1)) * (1 - X ^ (m + 1))

theorem jacobiTriple (u : Aˣ) :
    (1 - C (u : A)) * infProd (tripleFactor u) = tripleSum u
\end{leancode}
\begin{dictionary}
\lean{negSign j} $=(-1)^{|j|}=(-1)^j$. \lean{pentExp j} $=j(j-1)/2\ge0$ (\lean{toNat} is harmless: the value is
a nonnegative integer). \lean{tripleFactor u m} is the $(m+1)$-st factor. $u$ is any unit, e.g.\ a monomial
$x^n$ in $\Z[x^{\pm1}]$, giving the two-variable identity coefficientwise (\lean{twoVariableJTP},
\leanpath{JacobiTripleD4.lean}).
\end{dictionary}
\begin{proof}[Proof on paper]
\emph{Step 1: Cauchy's finite $q$-binomial theorem.} For every $m$,
$\prod_{i=0}^{m-1}(1+zq^i)=\sum_{k=0}^m\binom mk_qq^{k(k-1)/2}z^k$. This is proved by induction on $m$ from the
$q$-Pascal rule.

\emph{Step 2: the finite triple product.} Pull a factor $-q^i/u$ out of each $1-q^i/u$:
\begin{equation*}
(u;q)_n\,(q/u;q)_n=\prod_{i=0}^{n-1}(1-uq^i)\prod_{i=1}^{n}(1-q^i/u)
=(-1)^nu^{-n}q^{n(n+1)/2}\prod_{i=0}^{2n-1}\big(1-uq^{-n}q^{i}\big).
\end{equation*}
Apply Step 1 with $m=2n$ and $z=-uq^{-n}$, and write $k=n+j$. The exponent of $q$ becomes
$\binom{n+j}{2}-n(n+j)+\binom{n+1}{2}=\binom j2$, so
\begin{equation*}
(u;q)_n(q/u;q)_n=\sum_{j=-n}^{n}(-1)^ju^jq^{j(j-1)/2}\binom{2n}{n+j}_q .
\end{equation*}

\emph{Step 3: the limit.} For fixed $j$, $\binom{2n}{n+j}_q=(q)_{2n}/((q)_{n+j}(q)_{n-j})$ agrees with
$(q)_\infty/((q)_\infty(q)_\infty)=1/(q)_\infty$ to order $n-|j|$. Similarly the left side agrees with
$(u;q)_\infty(q/u;q)_\infty$ to order $n$. Letting $n\to\infty$ coefficientwise and multiplying by $(q)_\infty$ gives
$(u;q)_\infty(q/u;q)_\infty(q)_\infty=\sum_j(-1)^ju^jq^{j(j-1)/2}$, which is the statement, since
$(u;q)_\infty=(1-u)\prod_{m\ge1}(1-uq^m)$.
\end{proof}
\begin{physics}
Boson--fermion correspondence: the partition function of a complex fermion with a fugacity $u$ for the $U(1)$ charge
equals the sum over charge sectors of a free boson (charge $j$, energy $j^2/2$).
\end{physics}
\begin{status}
Standard axioms (re-run for this chapter). Unconditional. Not in Mathlib.
\end{status}

\noindent\emph{Side result} (feeds Lorenz's formula only).
\begin{theorem}[Euler's pentagonal theorem; the sum side of the quintuple product]\label{l25:thm:pent}
$(q;q)_\infty=\sum_{r\in\Z}(-1)^rq^{r(3r-1)/2}$. With $J_1(u)=\sum(-1)^nu^np^{n(n-1)/2}$ and $J_2(u)=\sum u^mp^{m^2}$,
$J_1(u)J_2(u)=(p;p)_\infty\sum_j\big(u^{3j}p^{j(3j-1)}-u^{3j+1}p^{j(3j+1)}\big)$ (the sum side of the quintuple
product identity).
\end{theorem}
\noindent Files \leanpath{VOA/\allowbreak Commutative/\allowbreak JacobiTriplePentagonal.lean},
\leanpath{VOA/\allowbreak Commutative/\allowbreak JacobiTripleQuintuple.lean}.
\begin{leancode}
def pentSeries : PowerSeries R := zSeries (fun r => negSign r) pentagonal

theorem euler_eq_pentagonal : (euler : PowerSeries R) = pentSeries := by

theorem quintuple_sum (u : Rˣ) :
    fsr (fun n : ℤ => negSign n * ((u ^ n : Rˣ) : R)) pentExp *
        fsr (fun m : ℤ => ((u ^ m : Rˣ) : R)) sqExp =
      euler * (fsr (fun j : ℤ => ((u ^ (3 * j) : Rˣ) : R)) (fun j => 2 * pentagonal j) -
        fsr (fun j : ℤ => ((u ^ (3 * j + 1) : Rˣ) : R)) (fun j => 2 * pentagonal (-j)))
\end{leancode}
\begin{dictionary}
\lean{pentagonal} is Mathlib's $k\mapsto k(3k-1)/2$ (Mathlib defines the numbers but does not prove the theorem).
\lean{fsr c e} $=\sum_ic_iq^{e(i)}$ coefficientwise as a \lean{finsum} (\leanpath{JacobiTripleFiber.lean}).
\lean{sqExp n} $=n^2$. With $q=p^2$ the second statement is the sum-side form of the classical quintuple product identity
(the sum $J_1J_2$ rewritten as $(p;p)_\infty$ times a two-term theta-type sum); the product form of the quintuple
product identity was not needed and is not formalised, so the library does not contain ``the quintuple product
identity'' in its usual product form.
\end{dictionary}
\begin{proof}[Proof on paper]
\emph{Pentagonal.} Formally this is the triple product with base $q^3$ and $u=q$:
$(q;q^3)_\infty(q^2;q^3)_\infty(q^3;q^3)_\infty=\sum_j(-1)^jq^{j}q^{3j(j-1)/2}=\sum_j(-1)^jq^{j(3j-1)/2}$, and the left side
is $(q)_\infty$, because every $n\ge1$ is $\equiv0,1,2\pmod3$. Since $u=q$ is not a unit of $A[[q]]$,
Theorem~\ref{l25:thm:jtp} cannot be applied directly. One reruns Steps 1--3 of its proof with $z=-q\cdot q^{-3n}$ at
base $q^3$, which gives the finite identity
$\prod_{m<n}(1-q^{3m+2})\prod_{i<n}(1-q^{3i+1})=\sum_{k\le2n}(-1)^{k-n}q^{(k-n)(3(k-n)-1)/2}\binom{2n}{k}_{q^3}$, and lets
$n\to\infty$.

\emph{Quintuple.} The coefficient of $u^k$ in $J_1(u)J_2(u)$ is $\sum_{n}(-1)^np^{E(n)}$ with
$E(n)=\binom n2+(k-n)^2$. Write $k=3j+t$ with $t\in\{0,1,2\}$ and complete the square.
\begin{itemize}
\item $t=0$, $n=r+2j$: $E=\frac{r(3r-1)}2+j(3j-1)$, and the sign is $(-1)^r$. By the pentagonal theorem the sum is
  $p^{j(3j-1)}(p;p)_\infty$.
\item $t=1$, $n=r+2j+1$: $E=\frac{r(3r+1)}2+j(3j+1)$, and the sign is $-(-1)^r$. Replacing $r$ by $-r$ gives
  $-p^{j(3j+1)}(p;p)_\infty$.
\item $t=2$, $n=r+2j+1$: $E=\frac{3r(r-1)}2+\text{const}$, and the sign is $-(-1)^r$. The terms $r$ and $1-r$ have the
  same exponent and opposite signs, so the sum is $0$.
\end{itemize}
Summing over $k$ gives the display.
\end{proof}
\begin{physics}
$(q;q)_\infty=q^{-1/24}\eta$ is the inverse partition function of one boson; the pentagonal theorem is the
character of the trivial $c=0$ minimal model written as a Virasoro Verma-module alternating sum.
\end{physics}
\begin{status}
Standard axioms. Unconditional.
\end{status}

\begin{theorem}[Legendre's four-triangular-numbers theorem]\label{l25:thm:legendre}
In $\Z[[q]]$: $\displaystyle\sum_{v\ge0}\frac{(2v+1)q^v}{1-q^{2v+1}}=(q;q)_\infty^4\prod_{k\ge1}(1+q^k)^8$.
\end{theorem}
\noindent File \leanpath{VOA/\allowbreak Commutative/\allowbreak FourTriangular.lean}.
\begin{leancode}
def lambertS : PowerSeries R :=
  powerShiftSum 1 fun v => ((2 * v + 1 : ℕ) : PowerSeries R) * geomInvSeries (2 * v + 1)

theorem lambertS_eq :
    (lambertS : PowerSeries ℤ) = euler ^ 4 * infProd fun k => (1 + X ^ (k + 1)) ^ 8
\end{leancode}
\begin{dictionary}
\lean{geomInvSeries m} $=1/(1-q^m)$; \lean{powerShiftSum 1 f} $=\sum_vq^vf(v)$. The right side is
$\psi(q)^4$ with $\psi=\sum q^{n(n+1)/2}=(q;q)_\infty(-q;q)_\infty^2$ (Gauss), and the left side is
$\sum_n\sigma(2n+1)q^n$ --- so this is Legendre's count of representations by four triangular numbers.
\end{dictionary}
\begin{proof}[Proof on paper]
Work in $\Z[[x]]$ with $q=x^2$ and $\ltwentyfiveEu=x\frac d{dx}$. Put $\varphi(x)=\sum_{n\in\Z}x^{n^2}$,
$Q_\pm=(\mp x;x^2)_\infty$, $P_2=(x^2;x^2)_\infty$, and let $S(q)$ be the left side of the theorem.
\begin{enumerate}
\item \emph{Gauss} (the triple product at $u=\mp x$, base $x^2$): $\varphi(x)=P_2Q_+^2$ and $\varphi(-x)=P_2Q_-^2$.
\item \emph{Log-derivative of the products.} $Q_+=\prod_k(1+x^{2k+1})$ and $Q_-=\prod_k(1-x^{2k+1})$, and
  $\ltwentyfiveEu\log\frac{1+y}{1-y}=\frac{2\ltwentyfiveEu y}{1-y^2}$ with $y=x^{2k+1}$, $\ltwentyfiveEu y=(2k+1)y$.
  Summing over $k$ gives $\frac{\ltwentyfiveEu Q_+}{Q_+}-\frac{\ltwentyfiveEu Q_-}{Q_-}=2xS(x^2)$.
\item \emph{The Wronskian as a product.} Let $W=\ltwentyfiveEu\varphi(x)\cdot\varphi(-x)-\varphi(x)\cdot\ltwentyfiveEu\varphi(-x)$.
  By items 1--2, $W=\varphi(x)\varphi(-x)\cdot2\big(\frac{\ltwentyfiveEu Q_+}{Q_+}-\frac{\ltwentyfiveEu Q_-}{Q_-}\big)
  =4xS(x^2)\,P_2^2Q_+^2Q_-^2$.
\item \emph{The Wronskian as a sum.} $W=\sum_{m,n}(m^2-n^2)(-1)^nx^{m^2+n^2}$. Pairs of equal parity cancel under
  $(m,n)\mapsto(n,m)$. Opposite-parity pairs are $(m,n)=(a+b+1,a-b)$, up to the swap, with $m^2-n^2=(2a+1)(2b+1)$ and
  $m^2+n^2=2a(a+1)+2b(b+1)+1$. The sum factorises: $W=xZ^2$, where $Z=\sum_{a\in\Z}(-1)^a(2a+1)x^{2a(a+1)}$.
\item \emph{Jacobi's cube identity} $(q)_\infty^3=\sum_{a\ge0}(-1)^a(2a+1)q^{a(a+1)/2}$, at $q=x^4$, gives $Z=2(x^4;x^4)_\infty^3$.
\item Equate items 3 and 4 and cancel $4x$: $S(x^2)P_2^2Q_+^2Q_-^2=(x^4;x^4)^6_\infty$. Undo $q=x^2$: $Q_+Q_-=(q;q^2)_\infty$
  and $P_2=(q)_\infty$. Euler's $(q;q^2)_\infty(-q;q)_\infty=1$ and $(q^2;q^2)_\infty=(q)_\infty(-q;q)_\infty$ then give
  $S=(q)_\infty^4(-q;q)_\infty^8$.
\end{enumerate}
\end{proof}
\begin{physics}
$\psi(q)^4$ is, up to normalisation, the fourth power of the Jacobi theta constant $\theta_2(\tau)=2q^{1/8}\psi(q)$, the \emph{numerator} of the
Ramond-sector character of eight Majorana--Weyl fermions ($SO(8)$ level one). The character itself is
$\theta_2^4/(2\eta^4)=8q^{1/3}\psi^4/(q;q)_\infty^4=8q^{1/3}\prod(1+q^r)^8$. The identity is the $A_2$ star identity's
last step, where $\prod(1+q^r)^8=q^{-1/3}(\eta(2\tau)/\eta(\tau))^8$ is the character of
$L_{-3/2}(\mathfrak{sl}_3)$.
\end{physics}
\begin{status}
Standard axioms (re-run for this chapter). Unconditional. Discharges \lean{A2.FourTriangular}
(\lean{A2.fourTriangular}, \leanpath{FourTriangularDeligne.lean}).
\end{status}

\subsection{Ramanujan's identities from the heat equation}

\begin{theorem}[Ramanujan's three identities]\label{l25:thm:ramanujan}
With $E_2=1-24\sum\sigma_1q^n$, $E_4=1+240\sum\sigma_3q^n$, $E_6=1-504\sum\sigma_5q^n$ in $\Q[[q]]$:
\begin{equation*}
12\ltwentyfiveEu E_2=E_2^2-E_4,\qquad 3\ltwentyfiveEu E_4=E_2E_4-E_6,\qquad 2\ltwentyfiveEu E_6=E_2E_6-E_4^2 .
\end{equation*}
\end{theorem}
\noindent Files \leanpath{VOA/\allowbreak DeligneAudit/\allowbreak DeligneCharacterD4MLDE.lean},
\leanpath{DeligneMLDERamanujan.lean}, \leanpath{DeligneMLDERamanujanFull.lean},
\leanpath{DeligneMLDERamanujanAdapters.lean}, \leanpath{DeligneMLDEEightSq.lean}.
\begin{leancode}
def sigS (k : ℕ) : PowerSeries ℚ := PowerSeries.mk fun n => if n = 0 then 0 else (sigmaK k n : ℚ)

def E2s : PowerSeries ℚ := 1 - 24 * sigS 1
def E4s : PowerSeries ℚ := 1 + 240 * sigS 3
def E6s : PowerSeries ℚ := 1 - 504 * sigS 5

structure RamanujanIdentities : Prop where
  e2 : 12 * eulerOp E2s = E2s ^ 2 - E4s
  e4 : 3 * eulerOp E4s = E2s * E4s - E6s
  e6 : 2 * eulerOp E6s = E2s * E6s - E4s ^ 2

theorem ramanujanIdentities : RamanujanIdentities := by
\end{leancode}
\begin{leancode}
theorem ramanujan : C 12 * (P1 * P1) = C 5 * P3 + P1 - C 6 * eulerOp P1 := by

theorem eight_sq (n : ℕ) :
    ∑ p ∈ Finset.antidiagonal n, e4 p.1 * e4 p.2 = e8 n := by
\end{leancode}
\begin{dictionary}
Namespace \lean{VOA.DeligneAudit.DeligneCharacter.D4MLDE} for the first block; \lean{P1}, \lean{P3} are
$\sum\sigma_1q^n$, $\sum\sigma_3q^n$ built from the divisor series \lean{lamQ}; \lean{ramanujan} is
$12P_1^2=5P_3+P_1-6\ltwentyfiveEu P_1$, which is $12\ltwentyfiveEu E_2=E_2^2-E_4$ after substituting $E_2=1-24P_1$ and
$E_4=1+240P_3$. \lean{eight_sq} is $E_4^2=E_8$ coefficientwise ($e_4(0)=1$, $e_4(n)=240\sigma_3(n)$, $e_8(n)=480\sigma_7(n)$).
The structure \lean{RamanujanIdentities} is a \lean{Prop}, and \lean{ramanujanIdentities} proves it.
\end{dictionary}
\begin{proof}[Proof on paper]
\emph{Setting.} Let $A=\Q[[x]]$ and $u=\varepsilon e^x\in A^\times$ with $\varepsilon=\pm1$. Write $\partial=d/dx$, acting on
coefficients, and note $\partial u=u$. The triple product (Theorem~\ref{l25:thm:jtp}) in $A[[q]]$ reads
$\Theta=(1-u)P$, where
\begin{equation*}
\Theta=\sum_{j\in\Z}(-1)^ju^jq^{j(j-1)/2},\qquad P=\prod_{m\ge1}(1-uq^m)(1-u^{-1}q^m)(1-q^m).
\end{equation*}

\emph{Two derivatives of the two sides.} On the sum side, $\partial u^j=ju^j$ and $\ltwentyfiveEu q^{j(j-1)/2}=\frac{j(j-1)}2q^{j(j-1)/2}$,
so $\Theta$ satisfies the \emph{heat equation} $\partial^2\Theta-\partial\Theta=2\ltwentyfiveEu\Theta$. On the product side,
the logarithmic derivatives are
\begin{equation*}
R:=\frac{\partial P}{P}=\sum_{m\ge1}\Big(\frac{u^{-1}q^m}{1-u^{-1}q^m}-\frac{uq^m}{1-uq^m}\Big),\qquad
T:=\frac{\ltwentyfiveEu P}{P}=-\sum_{m\ge1}m\Big(\frac{uq^m}{1-uq^m}+\frac{u^{-1}q^m}{1-u^{-1}q^m}+\frac{q^m}{1-q^m}\Big).
\end{equation*}

\emph{The master identity.} Substitute $\Theta=(1-u)P$ into the heat equation. Since $\partial(1-u)=-u$,
\begin{equation*}
\partial^2\Theta-\partial\Theta=P\big[(1-u)(\partial R+R^2)-(1+u)R\big],\qquad 2\ltwentyfiveEu\Theta=2(1-u)PT .
\end{equation*}
Dividing by the unit $P$ gives
\begin{equation*}
(1-u)\big(\partial R+R^2-2T\big)=(1+u)R. \tag{$*$}
\end{equation*}

\emph{Reading off the identities.} At $\varepsilon=1$, expand $R$ and $T$ in powers of $x$. The coefficient of $x^{k}$ in
$\frac{e^{\pm x}q^m}{1-e^{\pm x}q^m}$ is $\frac{(\pm1)^k}{k!}\sum_{d\ge1}d^kq^{md}$. So the $x$-coefficients of $R$ and $T$ are,
up to constants, the divisor series $\sum_n\sigma_k(n)q^n$ and $\sum_nn\sigma_{k}(n)q^n=\ltwentyfiveEu\sum_n\sigma_k(n)q^n$.
Comparing the coefficients of $x^1$ in $(*)$ gives a tautology (the involution $d\leftrightarrow n/d$ of divisors). The
coefficient of $x^3$ gives $12P_1^2=5P_3+P_1-6\ltwentyfiveEu P_1$, i.e.\ $12\ltwentyfiveEu E_2=E_2^2-E_4$. The coefficient
of $x^5$ gives $3\ltwentyfiveEu E_4=E_2E_4-E_6$. The coefficient of $x^7$ gives a relation of weight eight,
$2\ltwentyfiveEu E_6=E_2E_6-E_8$, with $E_8=1+480\sum\sigma_7(n)q^n$.

\emph{The step $E_4^2=E_8$.} This does \emph{not} come out of the heat equation (warning below). The space $M_8$ of
level-one modular forms of weight $8$ is one-dimensional, spanned by $E_4^2$, and contains no cusp form. So $E_4^2-E_8$,
being in $M_8$ with constant term $1-1=0$, is zero. This is Mathlib's level-one Sturm bound
(\lean{ModularForm.sturm_bound_levelOne_nat}) applied at weight $8$, where it says that a form whose constant term vanishes is
zero. It is the one analytic input of the section.
\end{proof}
\begin{physics}
Ramanujan's identities say that $\vartheta_k=\ltwentyfiveEu-\frac k{12}E_2$ (the Serre derivative) maps modular forms to modular
forms: they are the formal shadow of modularity, and they are exactly what is needed to verify that an explicit
$q$-series solves an MLDE.
\end{physics}
\begin{status}
Standard axioms for \lean{ramanujanIdentities} (re-run for this chapter). Unconditional. Uses Mathlib's modular forms
only through \lean{eight_sq}.
\end{status}

\begin{warning}[title={Dead end: the heat equation alone does not give $E_4^2=E_8$}]
The single-variable heat equation for the theta function gives one relation per weight; at weight 8 the $x^6$
relation contains $\sigma_7$, and no combination eliminates it. Do not retry; the cheapest source is Mathlib's
dimension formula, as used.
\end{warning}

\begin{theorem}[$A_2$: the product solves the $h=3$ recurrence]\label{l25:thm:a2prod}
$f=\prod_{r\ge1}(1+q^r)^8$ satisfies $6\ltwentyfiveEu^2f+3\ltwentyfiveEu f+24P_1\ltwentyfiveEu f+8P_1f-80P_3f=0$; hence
$[q^N]f=\mathtt{mldeCoeff}\ 3\ N$.
\end{theorem}
\noindent Files \leanpath{VOA/\allowbreak DeligneAudit/\allowbreak DeligneMLDEProduct.lean},
\leanpath{VOA/\allowbreak DeligneAudit/\allowbreak DeligneMLDE.lean}.
\begin{leancode}
def fQ : PowerSeries ℚ := infProd fun k => (1 + X ^ (k + 1)) ^ 8

theorem product_mlde :
    6 * eulerOp (eulerOp fQ) + 3 * eulerOp fQ + 24 * (P1 * eulerOp fQ) + 8 * (P1 * fQ)
      - 80 * (P3 * fQ) = 0 := by

structure A2ProductMLDE : Prop where
  sol : IsAKSolution (K := ℚ) 3 (fun N => (productCoeff N : ℚ))

theorem a2ProductMLDE : A2ProductMLDE :=
\end{leancode}
\begin{dictionary}
\lean{productCoeff N} $=[q^N]\prod(1+q^r)^8$ (computable, \leanpath{A2CharacterInterface.lean}). The displayed ODE is
\eqref{l25:eq:mlde} at $h=3$ for $g=q^{1/3}f$ (exponent $(h+1)/12=1/3$), multiplied by $6$ and rewritten with
$E_2=1-24P_1$, $E_4=1+240P_3$: indeed $(\ltwentyfiveEu-\frac16E_2)\ltwentyfiveEu g=\frac1{18}E_4g$ becomes
$\ltwentyfiveEu^2f+\frac12\ltwentyfiveEu f+4P_1\ltwentyfiveEu f+\frac43P_1f-\frac{40}3P_3f=0$.
\end{dictionary}
\begin{proof}[Proof on paper]
\emph{The log-derivative.} Since $\frac{q^r}{1+q^r}=\frac{q^r}{1-q^r}-\frac{2q^{2r}}{1-q^{2r}}$,
\begin{equation*}
\ltwentyfiveEu\log f=8\sum_r\frac{rq^r}{1+q^r}=8P_1(q)-16P_1(q^2)=\tfrac13\big(2E_2(q^2)-E_2(q)-1\big).
\end{equation*}
Put $g=q^{1/3}f$ and $F=2E_2(q^2)-E_2(q)$. Then $\ltwentyfiveEu g=\frac13Fg$ and
$\ltwentyfiveEu^2g=\big(\frac19F^2+\frac13\ltwentyfiveEu F\big)g$.

\emph{Reduction to an Eisenstein identity.} The MLDE $(\ltwentyfiveEu-\frac16E_2)\ltwentyfiveEu g=\frac1{18}E_4g$, multiplied by
$18/g$, becomes $6\ltwentyfiveEu F+2F^2-E_2F-E_4=0$. By Ramanujan at $q$ and at $q^2$,
$\ltwentyfiveEu F=\frac13\big(E_2(q^2)^2-E_4(q^2)\big)-\frac1{12}\big(E_2^2-E_4\big)$. Substituting, the equation becomes
\begin{equation*}
F^2=\tfrac15\big(E_4(q)+4E_4(q^2)\big),
\end{equation*}
the classical level-two identity $\big(2E_2(2\tau)-E_2(\tau)\big)^2=\frac15\big(E_4(\tau)+4E_4(2\tau)\big)$.

\emph{The level-two identity.} It comes from the master identity $(*)$ of Theorem~\ref{l25:thm:ramanujan}, now at
$\varepsilon=-1$, i.e.\ $u=-e^x$. There the geometric series of $R$ and $T$ carry signs: their $x$-coefficients
are the alternating divisor sums $\sum_{dk=n}(-1)^kd^j$, which are combinations of $\sigma_j(n)$ and $\sigma_j(n/2)$,
i.e.\ of series in $q$ and in $q^2$. The coefficients of $x^0$ and $x^2$ of $(*)$ give the needed identity
(\lean{level_two}). Uniqueness (Theorem~\ref{l25:thm:unique}) then gives $[q^N]f=c_N$.
\end{proof}
\begin{physics}
$q^{1/3}\prod(1+q^r)^8=(\eta(2\tau)/\eta(\tau))^8$: the $A_2$ vacuum character is a $\Gamma_0(2)$ eta quotient, the
Schur index of the $(A_1,D_4)$ theory.
\end{physics}
\begin{status}
Standard axioms. Unconditional.
\end{status}

\subsection{Level three: Lorenz and Borwein, for \texorpdfstring{$G_2$}{G2}}

\noindent\emph{Side result} at this commit: it proves the target side of the open $G_2$ star identity.
\begin{theorem}[Lorenz's formula; Borwein's cubic identity; the $G_2$ target]\label{l25:thm:g2target}
Let $a=\theta_{A_2}=\sum_{x,y\in\Z}q^{x^2-xy+y^2}$, $b^3=(q;q)^9_\infty/(q^3;q^3)^3_\infty$, $c^3=27q(q^3;q^3)^9_\infty/(q;q)^3_\infty$.
Then $a=1+6\sum_{n\ge1}\big(\sum_{d\mid n}\chi_{-3}(d)\big)q^n$ (Lorenz), $a^3=b^3+c^3$ (Borwein), $2a^2=3E_2(q^3)-E_2$,
$10a^4=E_4+9E_4(q^3)$. Consequently the Kac--Wakimoto form $T=\theta_{A_2}(q)(q^3;q^3)^6_\infty/(q;q)^8_\infty$ has
$[q^N]T=\mathtt{mldeCoeff}\ 4\ N$.
\end{theorem}
\noindent Files \leanpath{VOA/\allowbreak DeligneAudit/\allowbreak G2StarRiccati.lean},
\leanpath{G2StarMLDE.lean}, \leanpath{G2StarLorenz.lean}, \leanpath{G2StarReduction.lean},
\leanpath{G2StarUnconditional.lean}.
\begin{leancode}
def lambertA2 : PowerSeries ℚ :=
  PowerSeries.mk fun n => if n = 0 then 1 else 6 * ∑ d ∈ n.divisors, chi3 d

def LorenzA2 : Prop := AQ = lambertA2

theorem lorenzA2 : LorenzA2 := aq_eq_one_add_six_L0.trans lambertA2_eq_L0.symm

structure G2Level3Core : Prop where
  l1 : AQ ^ 3 = BQ + CQ
  l2 : 2 * AQ ^ 2 = 3 * E2q3 - E2q
  l3 : 10 * AQ ^ 4 = E4q + 9 * E4q3

theorem g2Level3Core : G2Level3Core := g2Level3Core_of_eisLambert g2Level3EisLambert

def g2KW : PowerSeries ℤ := thetaA2 * euler3 ^ 6 * reciprocalEuler ^ 8

structure G2KWEqMLDE : Prop where
  eq : ∀ N, ((coeff N g2KW : ℤ) : ℚ) = mldeCoeff 4 N

theorem g2KWEqMLDE : G2KWEqMLDE := g2KWEqMLDE_of_eisLambert g2Level3EisLambert
\end{leancode}
\begin{dictionary}
Namespace \lean{VOA.DeligneAudit.G2Star}. \lean{chi3 k} is $\chi_{-3}(k)$ ($1,-1,0$ for $k\equiv1,2,0$ mod 3). \lean{AQ} is
\lean{thetaA2} over $\Q$; \lean{thetaA2} counts $(x,y)\in[-N,N]^2$ with $x^2-xy+y^2=N$ (a correct finite box: $x^2-xy+y^2\ge\frac34\max(x^2,y^2)$, so every solution has
$|x|,|y|\le\sqrt{4N/3}\le N$ for $N\ge1$). \lean{BQ},
\lean{CQ} are $b^3$, $c^3$; \lean{E2q3} etc.\ are $E_k(q^3)$ via \lean{expand 3}. \lean{LorenzA2} and
\lean{G2KWEqMLDE} are \lean{Prop}s, and \lean{lorenzA2} and \lean{g2KWEqMLDE} prove them.
\end{dictionary}
\begin{proof}[Proof on paper]
\emph{Where $T$ comes from.} $T$ was found by a numerical search, as a closed form matching the first coefficients of
$\mathtt{mldeCoeff}\ 4$. It has the shape that Kac and Wakimoto's admissible-character formula would give. The theorem
does not depend on that origin: it proves that $T$ solves the $h=4$ equation.

\emph{Lorenz.} Split the $A_2$ lattice by parity ($4(m^2-mn+n^2)=(2m-n)^2+3n^2$), compute the $u$-Wronskian of the triple
product at the cube roots of unity, and identify two weight-$3/2$ sums by $z$-derivatives of the quintuple product
(\lean{quintuple_deriv₀}). The Lambert-form relations $l_2,l_3$ (and a weight-8 $l_6$) come from the heat-equation identity
(*) at $u=\omega e^x$, orders $x^0$--$x^4$, splitting over $\Q(\sqrt{-3})$ (\leanpath{DeligneMLDECubic*.lean}). Borwein's
cubic identity is then proved by a Riccati-type uniqueness argument. First, $d:=E_4-9E_4(q^3)+8a^4-16ac^3$ satisfies a
first-order equation $48a^2\,\ltwentyfiveEu d=d\,\big(d+32a^4+16a^2E_2+16ac^3\big)$ (an explicit polynomial identity),
and $d\equiv0\bmod q^2$. If $d\ne0$ had lowest term $\gamma q^n$, comparing the coefficients of $q^n$ would give
$48n\gamma=48\gamma$ (the bracket has constant term $32+16=48$), so $n=1$, contradicting $n\ge2$. Hence $d=0$. Then
$\psi=b^3+c^3-a^3$ satisfies $4\ltwentyfiveEu\psi=(E_2-a^2)\psi$ and $\psi(0)=0$, so $\psi=0$ (\lean{borweinCubic_of_eis}).
These give the level-three relations of Ramanujan's ``theory of signature 3'':
\begin{equation*}
E_4=a^4+8ac^3,\qquad E_6=a^6-20a^3c^3-8c^6,\qquad 48\,\ltwentyfiveEu a=a\big(4E_2-6a^2+2E_4a^{-2}\big),
\end{equation*}
the last from $l_2$, $l_3$ and Ramanujan for $E_2$.

\emph{$T$ solves the $h=4$ equation.} Put $f=q^{5/12}T$ (exponent $(h+1)/12=5/12$). From
$\ltwentyfiveEu\log(q)_\infty=(E_2-1)/24$ and the analogous formula at $q^3$,
\begin{equation*}
\ltwentyfiveEu\log f=\ltwentyfiveEu\log a+\tfrac34E_2(q^3)-\tfrac13E_2 .
\end{equation*}
Eliminate $E_2(q^3)$ by $l_2$ and $\ltwentyfiveEu\log a$ by the formula above. The $E_2$ terms cancel, leaving
\begin{equation*}
\ltwentyfiveEu\log f=:g=\frac{5a^3+4c^3}{12a}=\frac38a^2+\frac{E_4}{24a^2}.
\end{equation*}
The MLDE \eqref{l25:eq:mlde} at $h=4$ is $(\ltwentyfiveEu-\frac16E_2)\ltwentyfiveEu f=\frac{15}{144}E_4f$, i.e.\ the Riccati
equation $\ltwentyfiveEu g+g^2-\frac16E_2g=\frac{15}{144}E_4$. Compute $\ltwentyfiveEu g$ with the formula for
$\ltwentyfiveEu a$ and $3\ltwentyfiveEu E_4=E_2E_4-E_6$, and express $E_4$, $E_6$ through $a$ and $c^3$. The equation
becomes a polynomial identity in $a$, $c^3$, $E_2$ that holds identically. We checked this by computer algebra for this
chapter; in Lean it is a \lean{linear_combination} certificate (\lean{mldeExpr_TQ}). Uniqueness
(Theorem~\ref{l25:thm:unique}) gives $[q^N]T=c_N$.
\end{proof}
\begin{physics}
$\theta_{A_2}/\eta^2$ is the character of the $A_2$ lattice VOA, i.e.\ of the $SU(3)$ level-one WZW model. Here it
appears twisted by the level-three eta quotient $(q^3;q^3)^6_\infty/(q;q)^6_\infty$. \emph{Heuristic:} $T$ should be the
Kac--Wakimoto admissible character of $L_{-5/3}(G_2)$, written as a theta function over free bosons. The library does
not derive $T$ from the Kac--Wakimoto formula; it only proves that $T$ solves the MLDE.
\end{physics}
\begin{status}
Standard axioms. Unconditional.
\end{status}

%======================================================================
\section{Closed forms of the MLDE solution}\label{l25:sec:closed}
%======================================================================

For $h=6(t+1)$ ($D_4,E_6,E_7,E_8$ at $t=0,1,2,4$) the exponent $-c/24=(h+1)/12$ and the weight work out so that
$\eta^{2h-2}\cdot\mathrm{ch}$ is a quasimodular form $H$ of weight $h$ and depth $\le1$. Here $\mathrm{ch}$ is the
normalised vacuum character and $\eta=q^{1/24}(q)_\infty$. The $D_4$ case is on the critical path, since it feeds
Theorem~\ref{l25:thm:d4}. The $E_6$, $E_7$, $E_8$ cases are \emph{side results}: the open star bounds compare with $c_N$
directly.

\begin{theorem}[Generic closed form]\label{l25:thm:HOde}
Let $h=6t+6$, and let $H\in\Q[[q]]$ satisfy
$144\ltwentyfiveEu^2H-24hE_2\ltwentyfiveEu H+h(h-1)(E_2^2-E_4)H=0$. Suppose $G=H\cdot(q)_\infty^{-(12t+10)}$ vanishes
below $q^{t+1}$, with $[q^{t+1}]G=\gamma\ne0$. Then $[q^{N+t+1}]G/\gamma=\mathtt{mldeCoeff}\ (6t+6)\ N$ for all $N$.
\end{theorem}
\noindent File \leanpath{VOA/\allowbreak DeligneAudit/\allowbreak DeligneCharacterMLDEClosed.lean}.
\begin{leancode}
def Qk (k : ℕ) : PowerSeries ℚ := reciprocalEuler ^ k

def HOde (t : ℕ) (H : PowerSeries ℚ) : Prop :=
  144 * eulerOp (eulerOp H) - 24 * (6 * (t : PowerSeries ℚ) + 6) * E2s * eulerOp H +
    (6 * (t : PowerSeries ℚ) + 6) * (6 * (t : PowerSeries ℚ) + 5) * (E2s ^ 2 - E4s) * H = 0

theorem eq_mldeCoeff_of_HOde (t : ℕ) (H : PowerSeries ℚ)
    (hH : HOde t H) (hlow : ∀ n < t + 1, coeff n (H * Qk (12 * t + 10)) = 0) (c : ℚ)
    (hc : coeff (t + 1) (H * Qk (12 * t + 10)) = c) (hc0 : c ≠ 0) (N : ℕ) :
    coeff (N + (t + 1)) (H * Qk (12 * t + 10)) / c = mldeCoeff (6 * t + 6) N
\end{leancode}
\begin{dictionary}
Namespace \lean{VOA.DeligneAudit.DeligneCharacter.MLDEClosed}. Ramanujan's identities are used inside the proof
(\lean{ramanujanIdentities}), not as a hypothesis. \lean{Qk k} is $(q)_\infty^{-k}$, and the Lean \lean{c} is
$\gamma$. The shift by $t+1$ is the order of vanishing of $H$: $\eta^{2h-2}=q^{(h-1)/12}(q)_\infty^{2h-2}$ and the
character starts at $q^{-c/24}=q^{(h+1)/12}$, so $H=\eta^{2h-2}\mathrm{ch}$ starts at $q^{2h/12}=q^{t+1}$.
\end{dictionary}
\begin{proof}[Proof on paper]
Put $k=12t+10=2h-2$, $G=H(q)_\infty^{-k}$ and $f=q^{-(h-1)/12}G$. Since $\ltwentyfiveEu\log(q)_\infty=(E_2-1)/24$,
the prefactor $\Phi=q^{-(h-1)/12}(q)_\infty^{-k}$ of $f=\Phi H$ has $\ltwentyfiveEu\log\Phi=\varphi:=-\frac{h-1}{12}E_2$.
Then $\ltwentyfiveEu f=\Phi(\ltwentyfiveEu H+\varphi H)$ and
$\ltwentyfiveEu^2f=\Phi\big(\ltwentyfiveEu^2H+2\varphi\ltwentyfiveEu H+(\ltwentyfiveEu\varphi+\varphi^2)H\big)$, where Ramanujan
gives $\ltwentyfiveEu\varphi=-\frac{h-1}{144}(E_2^2-E_4)$. Substituting into \eqref{l25:eq:mlde} and multiplying by
$144/\Phi$, the coefficient of $\ltwentyfiveEu H$ is $144(2\varphi-\frac16E_2)=-24hE_2$, and the coefficient of $H$ is
\begin{equation*}
-(h-1)(E_2^2-E_4)+(h-1)^2E_2^2+2(h-1)E_2^2-(h^2-1)E_4=h(h-1)(E_2^2-E_4).
\end{equation*}
So \lean{HOde} is exactly \eqref{l25:eq:mlde} for $f$. Moreover $f=q^{(h+1)/12}\cdot(G/q^{t+1})$, because
$(h-1)/12+(h+1)/12=t+1$. Hence $G/(\gamma q^{t+1})$ is a normalised solution of \eqref{l25:eq:rec}, and uniqueness
(Theorem~\ref{l25:thm:unique}) concludes.
\end{proof}
\begin{physics}
$\eta^{2h-2}\mathrm{ch}=H$: the character is a quasimodular form divided by the partition function $\eta^{2h-2}$ of
$2h-2$ free bosons. For $E_6$ one finds $\eta^{22}\mathrm{ch}=(\sum n\sigma_9(n)q^n-\Delta)/1050$ (numerical evidence
through $q^{30}$, not machine-checked).
\end{physics}
\begin{status}
Standard axioms. Unconditional as an implication.
\end{status}

\begin{theorem}[$D_4$, $E_6$, $E_7$, $E_8$ closed forms]\label{l25:thm:closed}
With $H_{E_6}=7E_6^2+5E_4^3-12E_2E_4E_6$, $H_{E_7}=211E_2E_4E_6^2+77E_2E_4^4-91E_6^3-197E_4^3E_6$,
$H_{E_8}=318910E_2E_4^4E_6^2+145105E_2E_4E_6^4+33649E_2E_4^7-43225E_6^5-139129E_4^6E_6-315310E_4^3E_6^3$:
\begin{align*}
\mathtt{mldeCoeff}\ 12\ N&=[q^{N+2}]\,H_{E_6}(q)_\infty^{-22}/3991680,\\
\mathtt{mldeCoeff}\ 18\ N&=[q^{N+3}]\,H_{E_7}(q)_\infty^{-34}/42343741440,\\
\mathtt{mldeCoeff}\ 30\ N&=[q^{N+5}]\,H_{E_8}(q)_\infty^{-58}/13909248338175590400,\\
\mathtt{mldeCoeff}\ 6\ N&=[q^{N+1}]\,\tfrac1{240}\big(\ltwentyfiveEu E_4\big)(q)_\infty^{-10}.
\end{align*}
The last is the coefficientwise form of $[q^N]\,q^{-7/12}\,\ltwentyfiveEu E_4/(240\,\eta^{10})$: since
$\eta^{10}=q^{10/24}(q)_\infty^{10}$, the prefactor is $q^{-7/12-5/12}=q^{-1}$. Check: $\frac1{240}\ltwentyfiveEu E_4=\sum
n\sigma_3(n)q^n=q+18q^2+\dots$, so $c_0=1$ and $c_1=18+10=28$ (Example~\ref{l25:ex:a2}).
\end{theorem}
\noindent Files \leanpath{VOA/\allowbreak DeligneAudit/\allowbreak DeligneCharacterE6MLDE.lean},
\leanpath{DeligneCharacterE78MLDE.lean}, \leanpath{DeligneCharacterD4MLDE.lean}.
\begin{leancode}
def H6 : PowerSeries ℚ := 7 * E6s ^ 2 + 5 * E4s ^ 3 - 12 * E2s * E4s * E6s

def e6Target (N : ℕ) : ℚ := coeff N g6 / 3991680

theorem e6Target_eq_mldeCoeff (N : ℕ) : e6Target N = mldeCoeff 12 N := by
\end{leancode}
\begin{leancode}
theorem e7Target_eq_mldeCoeff (N : ℕ) :
    coeff (N + 3) (HE7 * Qk 34) / 42343741440 = mldeCoeff 18 N := by

theorem e8Target_eq_mldeCoeff (N : ℕ) :
    coeff (N + 5) (HE8 * Qk 58) / 13909248338175590400 = mldeCoeff 30 N := by
\end{leancode}
\begin{leancode}
def d4E4Target (N : ℕ) : ℚ := coeff N gser / 240

theorem d4E4Target_eq_mldeCoeff (hR : RamanujanIdentities) (N : ℕ) :
    d4E4Target N = mldeCoeff 6 N := by
\end{leancode}
\begin{dictionary}
\lean{g6} is $N\mapsto[q^{N+2}](H_6\cdot P^{-22})$ (\lean{G6 := H6 * Q22}). \lean{HE7}, \lean{HE8} are defined in
\leanpath{DeligneCharacterE78MLDE.lean} with the integer coefficients displayed above (written as nested products).
\lean{gser} is the shift by one of $\ltwentyfiveEu E_4\cdot(q)_\infty^{-10}$. The $D_4$ statement carries \lean{hR}, discharged by
\lean{ramanujanIdentities}. The integer polynomials $H$ were found by a computer fit outside the library. The Lean proof
only \emph{verifies} \lean{HOde} by a certificate and the low coefficients by \lean{decide}, so the provenance of the
numbers does not matter.
\end{dictionary}
\begin{proof}[Proof on paper]
Theorem~\ref{l25:thm:HOde} at $t=1,2,4$ (and the analogous direct proof at $t=0$), with \lean{HOde} verified by an
explicit linear combination of Ramanujan's identities and the low coefficients computed exactly.
\end{proof}
\begin{physics}
These are the closed forms of the Schur indices of the rank-one $E_6,E_7,E_8$ Minahan--Nemeschansky theories and of
$SU(2)$, $N_f=4$, as quasimodular forms over $\eta^{2h-2}$.
\end{physics}
\begin{status}
Standard axioms (re-run for this chapter for \lean{e7Target_eq_mldeCoeff} and \lean{e8Target_eq_mldeCoeff}).
Unconditional.
\end{status}

\begin{remark}[No closed form of this shape for $A_2$, $G_2$, $F_4$; side remark]
When $h/6\notin\Z$ no such level-one form exists: a search over $\eta^k\mathrm{ch}$, $k<80$, against level-one
quasimodular forms of weight $<40$ found nothing (a computer search; evidence, not proof). $A_2$ and $G_2$ have closed forms of
level $2$ and $3$ instead (Theorems~\ref{l25:thm:a2prod}, \ref{l25:thm:g2target}); $F_4$ has none in the library.
\end{remark}

%======================================================================
\section{Level-one affine Demazure characters}\label{l25:sec:demazure}
%======================================================================

\begin{definition}[The Demazure algorithm]\label{l25:def:demazure}
Fix a Cartan matrix $A$ (rows $A_{ij}=\langle\alpha_j,\alpha_i^\vee\rangle$), the highest root $\theta$ in simple-root
coordinates, and a reduced word for $s_\theta$. Characters are finite sums of $z^\mu q^e$. The finite Demazure operator
$D_i e^\lambda=(e^\lambda-e^{-\alpha_i}e^{s_i\lambda})/(1-e^{-\alpha_i})$ acts by the $\alpha_i$-string formula, and the
level-one affine $D_0$ uses $\langle\Lambda_0+\mu,\alpha_0^\vee\rangle=1-\langle\mu,\theta^\vee\rangle$ and
$e^{-\alpha_0}=q^{-1}z^\theta$. Then $C_m=q^{m^2}(D_{s_\theta}D_0)^m e^{\Lambda_0}$, and
$\mathrm{hilb}(m,N)$ is the sum of coefficients of $C_m$ in $q$-degree $N$ after setting $z=1$, clamped at $0$.
\end{definition}
\noindent File \leanpath{VOA/\allowbreak DeligneAudit/\allowbreak DeligneCharacterDemazure.lean}.
\begin{leancode}
def demazureStep (A : List (List ℤ)) (θ : List ℤ) (word : List ℕ) (f : LPoly) : LPoly :=
  word.foldr (fun i g => demazureFin A i g) (demazureZero A θ f)

def demazureChar (A : List (List ℤ)) (θ : List ℤ) (word : List ℕ) (m : ℕ) : LPoly :=
  (demazureStep A θ word)^[m] [((List.replicate θ.length 0, 0), 1)]

def levelOneDemazureHilb (A : List (List ℤ)) (θ : List ℤ) (word : List ℕ) (m N : ℕ) : ℕ :=
  (((demazureChar A θ word m).filter fun (t : Term) => decide (t.1.2 + (m : ℤ) * m = N)).map
    fun (t : Term) => t.2).sum.toNat
\end{leancode}
\begin{dictionary}
\lean{LPoly} is a list of terms $((\mu,e),\text{coefficient})$, merged by \lean{normalize}. \lean{demazureFin A i}
and \lean{demazureZero A θ} apply the string formulas termwise (definitions in the file). \lean{f^[m]} is $m$-fold
iteration starting from $e^{\Lambda_0}$ (the zero weight with energy $0$). \lean{levelOneDemazureHilb} selects terms of
energy $N-m^2$ and \textbf{clamps} the integer sum with \lean{toNat}. \textbf{Fidelity caveats.} (1) Reading
\lean{levelOneDemazureHilb A θ word m} as $\mathrm{Hilb}_qD(1,m\theta)$ needs the Demazure character formula
(Kumar, Mathieu) and the identification of $D(1,m\theta)$ with the Demazure module $V_{(s_\theta s_0)^m}(\Lambda_0)$
(Fourier--Littelmann, Chari--Venkatesh). This is U7 (Caveat~\ref{l25:cav:U7}); in Lean it is an algorithm. (2) The clamp
hides negative values; the ``unclamped'' integer version \lean{hilbIntG} is used where positivity matters
(Section~\ref{l25:sec:counts}). (3) The algorithm is correct for the non-simply-laced $G_2$ and $F_4$ as well. The
affine step needs $\langle\mu,\theta^\vee\rangle$, and the code computes $\sum_i\langle\theta,\alpha_i^\vee\rangle\mu_i$
from the Cartan matrix. In simple-root coordinates the two differ termwise by the factor
$(\theta,\theta)/(\alpha_i,\alpha_i)$, multiplying $(\theta,\alpha_i)$. For $G_2$ and $F_4$ every simple root is either
long, so the factor is $1$, or orthogonal to the long root $\theta$, so $(\theta,\alpha_i)=0$. Hence they agree
(\leanpath{DeligneCharacterExceptional.lean} explains this).
\end{dictionary}
\begin{physics}
The algorithm is the affine Weyl-group computation of a level-one integrable character (the $\widehat{\ltwentyfiveg}_1$
WZW vacuum module) truncated to a Schubert cell, so it is a piece of a genuine CFT character. \emph{Heuristic:}
$D(1,m\theta)$ is the graded fusion product of $m$ copies of $\ltwentyfiveg\oplus\C$ (the KR module
$\mathrm{KR}(\theta)$), and its $q$-grading is the energy of the corresponding quasi-particle configuration. That reading
is a theorem in the literature (Fourier--Littelmann, Chari--Venkatesh), but here it is input U7.
\end{physics}
\begin{status}
Definitions; computable; small cases kernel-checked: $\mathrm{Hilb}D(1,\theta)=\dim\ltwentyfiveg+q$ for $A_2$, $G_2$,
$F_4$, $E_6$ (\lean{levelOneDemazureHilb_A2_one}, \lean{g2DemazureHilb_one}, \lean{f4DemazureHilb_one},
\lean{e6DemazureHilb_one}), all in the library. For $E_7$ the corresponding check lies outside the library
(Section~\ref{l25:sec:defects}). For $E_8$ a single \lean{decide} at $m=1$ exceeded 10 minutes (a dead end); $C_1$ is
instead pinned by the chain files of Section~\ref{l25:sec:data}.
\end{status}

\begin{definition}[Per-type data and the named star identities]\label{l25:def:typedata}
For $X\in\{G_2,F_4,E_6,E_7,E_8\}$ the library fixes Bourbaki Cartan matrices, $\theta$ and reduced words of lengths
$5,15,21,33,57$, defines \lean{xDemazureHilb}, and names the star identity and target bound.
\end{definition}
\noindent Files \leanpath{VOA/\allowbreak DeligneAudit/\allowbreak DeligneCharacterExceptionalParameters.lean} (the Cartan
data, words and \lean{g2DemazureHilb}), \leanpath{VOA/\allowbreak DeligneAudit/\allowbreak DeligneCharacterExceptionalData.lean}
(the star identities) and \leanpath{VOA/\allowbreak DeligneAudit/\allowbreak DeligneCharacterExceptional.lean} (the finite
guards).
\begin{leancode}
def cartanG2 : List (List ℤ) := [[2, -3], [-1, 2]]
def thetaG2 : List ℤ := [3, 2]
def wordG2 : List ℕ := [1, 0, 1, 0, 1]

def g2DemazureHilb (m N : ℕ) : ℕ := levelOneDemazureHilb cartanG2 thetaG2 wordG2 m N

abbrev G2StarIdentity : Prop := StarIdentityQ g2DemazureHilb (mldeCoeff 4)
abbrev F4StarIdentity : Prop := StarIdentityQ f4DemazureHilb (mldeCoeff 9)
abbrev E6StarIdentity : Prop := StarIdentityQ e6DemazureHilb (mldeCoeff 12)
abbrev E7StarIdentity : Prop := StarIdentityQ e7DemazureHilb (mldeCoeff 18)
abbrev E8StarIdentity : Prop := StarIdentityQ e8DemazureHilb (mldeCoeff 30)

theorem g2_star_upto_two : ∀ N ∈ List.range 3,
    (demazureSumCoeff g2DemazureHilb N : ℚ) = mldeCoeff 4 N := by
\end{leancode}
\begin{dictionary}
Indices are $0$-based (Bourbaki node $i$ is list index $i-1$). For $G_2$, $\theta=3\alpha_1+2\alpha_2$ with $\alpha_1$
short: \lean{cartanG2} row $0$ is $\langle\cdot,\alpha_1^\vee\rangle$. The other types are analogous (quoted in the
file). The \lean{XStarIdentity} are \lean{abbrev}s of \lean{Prop}s, never axioms: a theorem that uses one takes it
as an explicit hypothesis. The finite guards \lean{g2_star_upto_two}, \lean{f4_star_upto_one} are kernel checks of the star identity at
$N\le2$, $N\le1$.
\end{dictionary}
\begin{physics}
\emph{Heuristic:} each named identity says ``quasi-particle (fermionic) character $=$ MLDE character'' for the
corresponding chiral algebra. Its mathematical content is the equality of two explicit $q$-series.
\end{physics}
\begin{status}
Definitions and small kernel checks. The identities themselves are \textbf{named inputs} (open), see
Section~\ref{l25:sec:star}.
\end{status}

%======================================================================
\section{The star identities}\label{l25:sec:star}
%======================================================================

\subsection{\texorpdfstring{$A_2$}{A2}: proved by an elementary route}

\begin{theorem}[The $A_2$ star identity]\label{l25:thm:a2star}
With $\mathrm{Hilb}_qD(1,n\theta)=P_{(2n,n,0)}(1,1,1;q,0)$ (the $q$-Whittaker polynomial, by Gelfand--Tsetlin
branching at $t=0$),
\begin{equation*}
\sum_{n\ge0}\frac{q^nP_{(2n,n,0)}(1,1,1;q,0)}{(q;q)_n}=\prod_{r\ge1}(1+q^r)^8 .
\end{equation*}
\end{theorem}
\noindent Files \leanpath{VOA/\allowbreak DeligneAudit/\allowbreak A2CharacterInterface.lean},
\leanpath{A2CharacterStarChain.lean}, \leanpath{A2Headline.lean}.
\begin{leancode}
def demazureHilb (n N : ℕ) : ℕ := gl3Hilb (2 * n) n 0 N

structure StarIdentity : Prop where
  eq : ∀ N, demazureCoeff N = productCoeff N

structure StarReduction : Prop where
  eq : (euler : PowerSeries ℤ) * starF = powerShiftSum 1 fun v => shiftedU (2 * v)

structure ShiftedMehler : Prop where
  eq : ∀ m : ℕ, shiftedU m * (euler : PowerSeries ℤ) ^ 3 * (1 - q ^ (m + 1)) = (m + 1 : ℕ)

structure FourTriangular : Prop where
  eq : (powerShiftSum 1 fun v => ((2 * v + 1 : ℕ) : PowerSeries ℤ) * geomInv (2 * v + 1)) =
    (euler : PowerSeries ℤ) ^ 4 * fermionInf

theorem starIdentity : StarIdentity := starIdentity_of_fourTriangular fourTriangular
\end{leancode}
\begin{dictionary}
Namespace \lean{VOA.DeligneAudit.A2}. \textbf{The table.} The $q$-Whittaker polynomial
$P_\lambda(x_1,x_2,x_3;q,0)$ at $x=(1,1,1)$ is computed by Gelfand--Tsetlin branching: restrict from $GL_3$ to $GL_2$ by
summing over partitions $\mu=(m_1,m_2)$ that interlace $\lambda$ ($\lambda_1\ge m_1\ge\lambda_2\ge m_2\ge\lambda_3$), then
from $GL_2$ to $GL_1$ likewise. At $t=0$ each step contributes the branching coefficient
$\psi_{\lambda/\mu}=\prod_i\binom{\lambda_i-\lambda_{i+1}}{\lambda_i-\mu_i}_q$. \lean{gl3Hilb l₁ l₂ l₃ N} is the $q^N$
coefficient of the result (\leanpath{A2CharacterInterface.lean}), and \lean{demazureHilb n} takes $\lambda=(2n,n,0)$.
Unwinding gives the closed form $\mathrm{Hilb}_qD(1,n\theta)=\sum_{i,j\le n}\binom ni_q\binom nj_qG_{n-i+j}$, with
$G_M=\sum_r\binom Mr_q$ the Rogers--Szeg\H{o} polynomial at $1$. \lean{demazureCoeff} $=$
\lean{demazureSumCoeff demazureHilb} by \lean{rfl}.

\textbf{The proof structures.} \lean{starF} and \lean{fermionInf} are the two sides as series, and
$\lean{shiftedU m}=U_m=\sum_sq^sG_sG_{s+m}/(q)_s$. The three structures are the three steps of the proof below, each a
\lean{Prop}, and each proved (\lean{starReduction}, \lean{shiftedMehler}, \lean{fourTriangular}).

\textbf{Reading.} The reading of \lean{demazureHilb} as $\mathrm{Hilb}_qD(1,n\theta)$ is U7 (Sanderson, Ion;
Caveat~\ref{l25:cav:U7}). It was compared with the affine algorithm \lean{levelOneDemazureHilb} for small $n$ and $N$;
the extent and location of those checks are in Section~\ref{l25:sec:defects}, and Example~\ref{l25:ex:a2} adds $n=2$,
$N\le3$.
\end{dictionary}
\begin{proof}[Proof on paper]
Write $F=\sum_nq^n\,\mathrm{Hilb}_qD(1,n\theta)/(q)_n$ for the left side.

\emph{Step 1} (\lean{StarReduction}: $(q)_\infty F=\sum_{v\ge0}q^vU_{2v}$). Use the $q$-Vandermonde identity in its
``overlapping subsets'' form,
\begin{equation*}
\binom ni_q\binom nj_q=\sum_kq^{(i-k)(j-k)}\frac{(q)_n}{(q)_k(q)_{i-k}(q)_{j-k}(q)_{n-i-j+k}} .
\end{equation*}
It counts pairs of subspaces of $\mathbb F_q^n$ (an $i$-dimensional and a $j$-dimensional one) by the dimension $k$ of
their intersection. Substituting into the closed form of $\mathrm{Hilb}_qD(1,n\theta)$ and reindexing turns $F$ into
\begin{equation*}
F=\sum_{k,u,v,w\ge0}\frac{q^{k+u+v+w+uv}\,G_{k+2v+w}}{(q)_k(q)_u(q)_v(q)_w}.
\end{equation*}
Euler's sum $\sum_u\frac{q^{u(1+v)}}{(q)_u}=\frac{(q)_v}{(q)_\infty}$ removes $u$. The convolution
$\sum_{k+w=s}\frac1{(q)_k(q)_w}=\frac{G_s}{(q)_s}$, which holds because $\binom sk_q=\frac{(q)_s}{(q)_k(q)_{s-k}}$,
collapses $k,w$ into $s=k+w$. The
result is $(q)_\infty F=\sum_vq^v\sum_s\frac{q^sG_sG_{s+2v}}{(q)_s}=\sum_vq^vU_{2v}$.

\emph{Step 2} (\lean{ShiftedMehler}: $U_m(q)_\infty^3(1-q^{m+1})=m+1$). A second use of $q$-Vandermonde gives
$U_m=\sum_{p\le m}\binom mp_qW_p$, where $W_p=\sum_{l,u}\frac{q^{l+u(p+1)}G_{l+u}}{(q)_l(q)_u}$. Rogers' Mehler formula for
the Rogers--Szeg\H{o} polynomials,
\begin{equation*}
\sum_{M\ge0}\frac{H_M(x)H_M(y)\,t^M}{(q)_M}=\frac{(xyt^2;q)_\infty}{(t;q)_\infty(xt;q)_\infty(yt;q)_\infty(xyt;q)_\infty},
\end{equation*}
at $x=1$, $y=q^p$, $t=q$ gives $W_p=\frac{(q)_p}{(1-q^{p+1})(q)_\infty^3}$. Since $\binom mp_q(q)_p=(q)_m/(q)_{m-p}$, it
remains to prove the finite identity
\begin{equation*}
\sum_{p=0}^m\frac{(q)_m}{(q)_{m-p}\,(1-q^{p+1})}=\frac{m+1}{1-q^{m+1}},
\end{equation*}
an identity of rational functions in $q$, valid for every $m$ (the Lean proof of \lean{shiftedMehler} is in
\leanpath{A2CarrierMehler.lean}).

\emph{Step 3} (\lean{FourTriangular}). By Steps 1 and 2,
$(q)_\infty^4F=(q)_\infty^3\sum_vq^vU_{2v}=\sum_v\frac{(2v+1)q^v}{1-q^{2v+1}}$. This is
$(q)_\infty^4\prod_r(1+q^r)^8$ by Legendre (Theorem~\ref{l25:thm:legendre}). Cancel $(q)_\infty^4$.

\emph{On novelty.} The paper's proof uses flavoured Macdonald polynomials at $t=0$. This route avoids them. It was found
in this project, and we do not know it in the literature.
\end{proof}
\begin{physics}
Fermionic sum over fusion products of adjoints $=$ the $(A_1,D_4)$ Schur index $(\eta(2\tau)/\eta(\tau))^8$ (the fermionic
reading being the heuristic of Definition~\ref{l25:def:dsum}). In free-field terms, $\prod(1+q^r)^8$ is the partition
function of eight Ramond-sector free fermions with their zero modes removed. So the identity says that the
quasi-particle description of the vacuum module counts states like eight free fermions.
\end{physics}
\begin{status}
Standard axioms; \textbf{unconditional} as a $q$-series identity. Its use as ``$D_N\le\dim V_N$'' for the true
Demazure modules needs U7.
\end{status}

\subsection{\texorpdfstring{$D_4$}{D4}: the elliptic evaluation}

\begin{theorem}[$D_4$: the series \lean{d4SchurQ} $=$ MLDE solution; Demazure $=$ \lean{d4SchurQ} under positivity]\label{l25:thm:d4}
Work over $\Bbbk=\Q((r))$, $r=q^{1/2}$. Let
\begin{equation*}
\Theta(x)=\sum_{k\in\Z}(-1)^k\big(q^{\binom k2}+q^{\binom{k+1}2}\big)x^{2k},\qquad
M(s,x,z)=\frac{(s^2;q)_\infty(s^2x^2;q)_\infty(s^2x^{-2};q)_\infty}
{\prod_{\epsilon=\pm1}\prod_{i=1}^4\prod_{\sigma=\pm1}(s\,x^\epsilon z_i^\sigma;q)_\infty},
\end{equation*}
a power series in $s$ with coefficients Laurent polynomials in $x,z_1,\dots,z_4$ over $\Bbbk$; let
$I(s,z)=\mathrm{CT}_x\,\Theta(x)M(s,x,z)$ (constant term in $x$), and let
\begin{equation*}
\mathtt{d4SchurQ}_N=\sum_{m=0}^{N}\Big[r^{2(N-m)}\Big]\ \tfrac12\,[s^{2m}]\,I(s,z)\big|_{z=1}.
\end{equation*}
Then, unconditionally, $\mathtt{d4SchurQ}_N=[q^{N+1}]\,\frac1{240}(\ltwentyfiveEu E_4)\,(q)_\infty^{-10}=c_N$ for $h=6$; in
modular notation this is $[q^N]\,q^{-7/12}\,\ltwentyfiveEu E_4/(240\eta^{10})$ (Theorem~\ref{l25:thm:closed}). If moreover
\lean{D4DemazurePositivity} holds, the $D_4$ Demazure sum equals $\mathtt{d4SchurQ}_N$.
\end{theorem}
\noindent Files \leanpath{VOA/\allowbreak DeligneAudit/\allowbreak DeligneCharacterD4MLDE.lean},
\leanpath{DeligneCharacterD4Wp12.lean}, \leanpath{DeligneCharacterD4Schur.lean}, \leanpath{DeligneCharacterD4Final.lean},
\leanpath{D4CharacterTop.lean}.
\begin{leancode}
def D4SchurEqE4 : Prop := ∀ N, D4K.d4SchurQ N = d4E4Target N

theorem d4SchurEqE4 : D4MLDE.D4SchurEqE4 := by

theorem d4SchurQ_eq_mldeCoeff_holds (N : ℕ) : D4K.d4SchurQ N = mldeCoeff 6 N :=
\end{leancode}
\noindent The defining objects, verbatim: \leanpath{VOA/\allowbreak DeligneAudit/\allowbreak DeligneCharacterD4Euler.lean}
($\Bbbk$), \leanpath{DeligneCharacterD4Series.lean} ($R$, $S$, $M$), \leanpath{DeligneCharacterD4Functional.lean}
($\Theta$), \leanpath{DeligneCharacterD4Assembly2.lean} ($I$), \leanpath{DeligneCharacterD4Schur.lean} (\lean{ev1},
\lean{d4SchurQ}).
\begin{leancode}
abbrev Kq := LaurentSeries ℚ

abbrev R := AddMonoidAlgebra Kq V

abbrev S := PowerSeries R

def tFactors : S :=
  EPos 1 0 2 two_ne_zero * EPos 1 (2 • xv) 2 two_ne_zero * EPos 1 (-2 • xv) 2 two_ne_zero

def hyperFactors : S := (hyperVecs.map fun v => EInv 1 v 1 one_ne_zero).prod

def kernelM : S := tFactors * hyperFactors

def thK (i : ℤ) : Kq := if Even i then sgn (i / 2) * (qC (i / 2) + qC (i / 2 + 1)) else 0

abbrev ctTheta : R →ₗ[Kq] R := ctFun thK

def Ikernel : S := smap ctTheta kernelM

def ev1 : R →ₗ[Kq] Kq := Finsupp.linearCombination Kq fun _ => (1 : Kq)

def d4SchurQ (N : ℕ) : ℚ :=
  ∑ m ∈ Finset.range (N + 1), ((2 : Kq)⁻¹ * ev1 (coeff (2 * m) Ikernel)).coeff (2 * ((N : ℤ) - m))
\end{leancode}
\noindent Files \leanpath{DeligneCharacterD4Schur.lean}, \leanpath{DeligneCharacterD4Final.lean},
\leanpath{DeligneCharacterD4Positivity.lean}.
\begin{leancode}
structure D4DemazurePositivity : Prop where
  energy : ∀ m : ℕ, ∀ t ∈ demazureChar cartanD4 thetaD4 wordD4 m, 0 ≤ t.1.2 + (m : ℤ) * m
  coeff : ∀ (m : ℕ) (N : ℕ), 0 ≤ hilbInt m N

def D4DemazureEqSchur : Prop := ∀ N, (demazureSumCoeff d4DemazureHilb N : ℚ) = d4SchurQ N

theorem d4DemazureEqSchur_of_positivity (hpos : D4DemazurePositivity) : D4DemazureEqSchur :=

theorem d4DemazurePositivity_of_U5Z
    (hU5 : DeligneSector.FusionModuleBoundG DeligneSector.D4.D4sector d4FusionFZ) :
    D4K.D4DemazurePositivity :=
\end{leancode}
\begin{dictionary}
\textbf{Decoding the series.} \lean{Kq} is $\Bbbk=\Q((r))$ with $r=q^{1/2}$ (\lean{rr}; \lean{qq} $=r^2$), so
half-integer $q$-shifts are ring maps. \lean{R} $=\Bbbk[x^{\pm1},z_1^{\pm1},\dots,z_4^{\pm1}]$ (exponent lattice
\lean{V := Fin 5 → ℤ}, coordinate $0$ is $x$) and \lean{S} $=R[[s]]$. \lean{EPos c v k} and \lean{EInv c v k} are
$(c\,z^vs^k;q)_\infty$ and its inverse; so \lean{tFactors} $=(t;q)(tx^2;q)(tx^{-2};q)$ with $t=s^2$, \lean{hyperFactors}
$=1/\prod(sx^{\pm1}z_i^{\pm1};q)$ over the sixteen weights \lean{hyperVecs}, and \lean{kernelM} is $M$.
\lean{thK i} is the coefficient $t_i$ of $x^i$ in $\Theta$ (\lean{qC j} $=q^{\binom j2}$, \lean{sgn} a sign; odd
$i$ give $0$); \lean{ctFun θf} sends $c\,x^az^v$ to $c\,t_{-a}z^v$, i.e.\ it is $F\mapsto\mathrm{CT}_x\Theta F$.
\lean{smap} applies it coefficientwise in $s$, so \lean{Ikernel} is $I$ (``$2I$'' in the file's docstring
convention). \lean{ev1} sums the coefficients of a Laurent polynomial, i.e.\ sets $x=z_i=1$. \lean{coeff (2 * m)}
takes $[s^{2m}]$, and \lean{.coeff (2 * (N - m))} reads off $r^{2(N-m)}=q^{N-m}$: the $s^{2m}$ term is weighted by
$q^m$, i.e.\ the series is evaluated at $s^2=t=q$. So \lean{D4K.d4SchurQ N} is exactly the displayed
$\mathtt{d4SchurQ}_N$. \textbf{Modelling step, not Lean.} That this series is the Schur index of the $SU(2)$,
$N_f=4$ gauge theory (the contour integral of the vector-multiplet and Haar-measure factor against four
hypermultiplets) is a hand transcription made when the file was written; no Lean statement relates
\lean{d4SchurQ} to a gauge-theory integral. The theorem is about the series as defined.

\textbf{Positivity.} \lean{hilbInt m N} is the unclamped integer Demazure coefficient. \lean{D4DemazurePositivity}:
energies $\ge-m^2$ and \lean{hilbInt}$\ge0$. The \lean{energy} field is proved (\lean{D4Wp.energy_holds}, from
integrality of the $D_4$ identity). The \lean{coeff} field is \emph{not} proved: the only route in the library is
\lean{d4DemazurePositivity_of_U5Z}, which takes the unclamped fusion bound \lean{hU5} (input U5 of
Chapter~\ref{ch:L26}, open) as an explicit hypothesis. So the $D_4$ Demazure-sum identity is conditional on
\lean{D4DemazurePositivity}, whose remaining field is currently available only from the open U5. (On paper, positivity
follows from the Demazure character formula, i.e.\ from U7; no Lean statement backs that reading.)
\end{dictionary}
\begin{proof}[Proof on paper]
\emph{Schur $=\ltwentyfiveEu E_4$.} Prove the classical $\wp'(z)=-\sigma(2z)/\sigma(z)^4$ in theta form from a two-variable identity
of lattice series (files \leanpath{DeligneCharacterD4Wp}--\leanpath{Wp3}). After the half-period substitution the Schur
integrand becomes a perfect square $-\eta^{-10}\mathcal W_1^2$, with $\mathcal W_1$ an explicit Lambert series
(\lean{integrand_sq}), and $\mathrm{CT}_x\mathcal W_1^2=-2\sum(N+1)\sigma_3(N+1)r^{2N}$. Since
$\sum_N(N+1)\sigma_3(N+1)q^{N}=q^{-1}\cdot\frac1{240}\ltwentyfiveEu E_4$, this is the closed form up to normalisation, and no
modular forms are used. Then $d4E4Target=c_N$ by Theorem~\ref{l25:thm:closed}. \emph{Demazure $=$ Schur.} The kernel $M$ satisfies a
$q$-difference recursion $I-I(qt)=tq\,B\,I(q^2t)$ given the finite polynomial identity \lean{D4FiniteIdentity}
(\lean{Ikernel_recursion}); its even coefficients are $\frac{2}{(q)_m}C_m$ with $C_m$ the Demazure characters
(\lean{Ikernel_coeff_even}). The finite identity is proved by kernel checkpoints (Section~\ref{l25:sec:data}).
Positivity is needed to compare with the clamped \lean{d4DemazureHilb}.
\end{proof}
\begin{physics}
$D_4$ at $k=-2$ is the chiral algebra of $SU(2)$ with four flavours, a Lagrangian theory: its Schur index is a contour
integral over the gauge fugacity $x$ of vector-multiplet and hypermultiplet contributions. With $s=q^{1/2}$ the
factor $M$ is the four hypermultiplets ($sx^{\pm1}z_i^{\pm1}$, $z_i$ the $SO(8)$ flavour fugacities) and the
$W$-boson part of the vector multiplet ($(s^2x^{\pm2};q)$), while $\Theta$ supplies the rest of the vector multiplet
and the Haar measure. This reading is physics modelling; the Lean theorem is about the series \lean{d4SchurQ}. Granting
it, the theorem evaluates the index in closed form and identifies it with the MLDE solution and with a fusion-product
sum.
\end{physics}
\begin{status}
Standard axioms for \lean{d4SchurEqE4} and \lean{d4SchurQ_eq_mldeCoeff_holds}. \lean{d4SchurQ}
$=c_N$: \textbf{unconditional} as a $q$-series statement (its reading as a gauge-theory Schur index is modelling, not
Lean). Demazure $=$ \lean{d4SchurQ}: conditional on \lean{D4DemazurePositivity}; the \lean{energy} field is proved, the
\lean{coeff} field only from the open U5 (\lean{d4DemazurePositivity_of_U5Z}).
Equivalently in one-sided form: \lean{d4StarBound_iff_leSchur}.
\end{status}

\subsection{\texorpdfstring{$G_2$}{G2}, \texorpdfstring{$F_4$}{F4}, \texorpdfstring{$E_6$}{E6}, \texorpdfstring{$E_7$}{E7}, \texorpdfstring{$E_8$}{E8}: named, in one-sided form}

\begin{definition}[The open star bounds]\label{l25:def:starbounds}
\lean{XStarBound} is $D_N\le c_N$ for all $N$, with the type's Demazure table and $h^\vee$.
\end{definition}
\noindent Files \leanpath{VOA/\allowbreak DeligneAudit/\allowbreak DeligneCharacterStarBound.lean},
\leanpath{G2StarReduction.lean}, \leanpath{G2StarUnconditional.lean}, \leanpath{G2StarFermionic.lean}.
\begin{leancode}
abbrev G2StarBound : Prop := StarBoundQ g2DemazureHilb (mldeCoeff 4)
abbrev F4StarBound : Prop := StarBoundQ f4DemazureHilb (mldeCoeff 9)
abbrev E6StarBound : Prop := StarBoundQ e6DemazureHilb (mldeCoeff 12)
abbrev E7StarBound : Prop := StarBoundQ e7DemazureHilb (mldeCoeff 18)
abbrev E8StarBound : Prop := StarBoundQ e8DemazureHilb (mldeCoeff 30)
\end{leancode}
\begin{leancode}
structure G2DemazureEqKW : Prop where
  eq : g2DemazureSeries = g2KW

theorem g2StarIdentity_of_demazure (hD : G2DemazureEqKW) : G2StarIdentity :=

structure G2DemazureFermionic : Prop where
  eq : ∀ n N, g2DemazureHilb n N = g2FermHilb n N

structure G2FermionicKW : Prop where
  eq : ∀ N, ((demazureSumCoeff g2FermHilb N : ℕ) : ℤ) = coeff N g2KW
\end{leancode}
\begin{dictionary}
For $G_2$ the target side is proved (\lean{g2KWEqMLDE}), so \lean{G2StarIdentity} rests only on
\lean{G2DemazureEqKW}, which splits further into the representation-theoretic \lean{G2DemazureFermionic} (Demazure
$=$ HKOTY fermionic formula, \lean{g2FermHilb}) and one $q$-series identity \lean{G2FermionicKW}
(\lean{g2DemazureEqKW_of_fermionic}). For $F_4$, $E_{6,7,8}$ nothing is split.
\end{dictionary}
\begin{physics}
Each is a Nahm-type ``fermionic $=$ bosonic'' identity of rank equal to the rank of $\ltwentyfiveg$; for $E_{6,7,8}$
these are the conjectural quasi-particle forms of the Minahan--Nemeschansky Schur indices.
\end{physics}
\begin{status}
\textbf{Named inputs; conjectural} (Conjecture~\ref{l27:conj:star}). Numerical evidence (Python, exact integers,
the same Demazure recursion truncated in degree; not machine-checked; \leanpath{Notes/deligne/STAR\_IDENTITY\_SCOPING.md},
2026-10-28): $D_N=c_N$ for $G_2$ through $q^{60}$, $F_4$ through $q^{25}$, $E_6$ through $q^{20}$, $E_7$ through
$q^{16}$, $E_8$ through $q^{14}$; no negative Demazure coefficient occurred. For $G_2$,
\lean{G2DemazureFermionic} follows from published results (Fourier--Littelmann, Chari--Moura, Ardonne--Kedem,
Di Francesco--Kedem; conventions per Lin, IMRN 2021), not formalized; \lean{G2FermionicKW} is open (no known proof)
and was checked through $q^{60}$ (\leanpath{Notes/deligne/G2\_STAR\_FERMIONIC.md}). Kernel-checked guards: $G_2$ $N\le2$, $F_4$ $N\le1$.
\end{status}

\subsection{Termwise domination cannot prove the star bound}\label{l25:sec:rigidity}

The one-sided star bound $D_N\le c_N$ is all that is needed downstream. One might hope it is easier than the identity.
The next proposition shows that, for the most natural kind of proof, it is not. It says nothing about other kinds of
proof.

\begin{proposition}[Rigidity of termwise domination; on paper]\label{l25:prop:rigidity}
Write $D=\sum_nq^nH_n/(q)_n$ with $H_n\ge0$ coefficientwise. Suppose $U_n\ge H_n$ coefficientwise for every $n$ and
$\sum_nq^nU_n/(q)_n\le c$ coefficientwise. If in fact $D=c$, then $U_n=H_n$ for every $n$.
\end{proposition}
\begin{proof}
The series $\sum_nq^n(U_n-H_n)/(q)_n$ has nonnegative coefficients and is $\le c-D=0$, so it is $0$. Each summand has
nonnegative coefficients, because $1/(q)_n$ does, so each summand is $0$. Since $1/(q)_n$ is a unit, $U_n=H_n$.
\end{proof}
\begin{physics}
Downstream there is no slack either: $c_N=\dim L_N$ (Caveat~\ref{l25:cav:equality}), so $D_N\le\dim L_N$ \emph{is}
classical freeness in counting form. A crystal-injection proof would have to be a bijection between Demazure-crystal
data and a combinatorial model of $c_N$, and no such model is known.
\end{physics}

\begin{warning}[title={Dead ends for the star bound (do not retry)}]
The status of each entry is stated: ``proved'' means on paper here; ``observed'' means a computer-algebra computation
in exact arithmetic, outside Lean.
\begin{itemize}
\item \textbf{Termwise domination proofs of $\le$} are forced to be exact by Proposition~\ref{l25:prop:rigidity} (proved),
  since numerically $D=c$ (observed).
\item \textbf{Positive MLDE expansions.} The AK kernel $10(h^2-1)\sigma_3-(24(N-r)+2(h+1))\sigma_1$ changes sign (proved: it
  is negative for large $N-r$). The product (plethystic) exponents $a_n$ of $\sum_Nc_Nq^N=\prod_n(1-q^n)^{-a_n}$
  (glossary) alternate in sign for every type, e.g.\ $G_2$: $14,-13,78,-337,\dots$ (observed). The only positive form
  found is $\mathrm{ch}=B\cdot(q)_\infty^{-(2h-2)}$ with $B\ge0$, observed through $q^{150}$ for $E_{6,7,8}$. For $E_6$,
  $B_n=((n+2)\sigma_9(n+2)-\tau(n+2))/1050$, which is positive by Deligne's bound $|\tau(n)|\le d(n)n^{11/2}$ (proved on
  paper). $B$ has no known combinatorial meaning.
\item \textbf{$G_2$ via $D_4$'s kernel.} $G_2$ at $k=-5/3$ has no known 4d Lagrangian theory, so there is no Schur-index
  contour integral to start from.
\item \textbf{Non-product structure.} Introduce a variable $z$ counting $n$. The series $\sum_nz^nH_n/(q)_n$ is not an
  infinite product in $(z,q)$, because its plethystic exponents grow and alternate (observed). The branching functions
  of the $G_2$ level-one character into $\ltwentyfiveg$-isotypic pieces are not products either (observed).
\end{itemize}
\end{warning}

\begin{warning}[title={The stopped Route A pilot for $G_2$}]
\emph{Route A} is the paper's proof: Macdonald/Cherednik constant terms at $t=0$. A pilot for $G_2$ was run and
stopped. Its first layer R1 is in the library, \leanpath{VOA/\allowbreak DeligneAudit/\allowbreak DeligneCharacterG2RouteA.lean}
($K=\mathrm{Frac}\,\Q[X^{\pm},Y^{\pm},q^{\pm}]$, monomial substitutions):
\begin{leancode}
def mono (v : E3) : K := algebraMap L K (monoL v)

def substL (M : E3 →+ E3) : L →+* L := mapDomainRingHom ℚ M
\end{leancode}
Numerically (exact rationals) $HB=q^{-2}BH$ holds as an identity of difference operators with 19 vanishing shift
coefficients, while $B=\text{(algorithm step)}$ and $H\circ\text{step}=q^{-2}\text{step}\circ H$ hold only on
Weyl-group invariants. \textbf{Blocker.} Proving R3 ($HB=q^{-2}BH$) by certificate is infeasible: the shift identities
carry 17--41 distinct binomial denominators, beyond what computer algebra could clear in reasonable time, and no
small vanishing subgroup was found among the 28 terms of the single-root shifts. The only viable proofs are
structural (DAHA/Cherednik $Y$-operators for $G_2$) or a citation of the literature (van Diejen--Emsiz plus Chari--Ion)
as a named input. The pilot was therefore stopped, and \lean{G2DemazureEqKW} stays named. R1 is kept as infrastructure;
it proves nothing about characters.
\end{warning}

%======================================================================
\section{\texorpdfstring{$q=1$}{q=1} counts and energies}\label{l25:sec:counts}
%======================================================================

\noindent\emph{Side results} for this chapter's chain. They supply the count and positivity inputs of the U5 route in
Chapter~\ref{ch:L26} and the \lean{energy} field of \lean{D4DemazurePositivity}.

\begin{theorem}[Tensor factorisation at $q=1$]\label{l25:thm:total}
Let $\theta$ have $r$ entries and let every operator of the word, and the affine one, have nonzero root $\alpha$. If
$\chi=C_1|_{q=1}$ is invariant under every reflection of the word and under $s_\theta$, then
$C_m|_{q=1}=\chi^m$; in particular the coefficient sum of $C_m$ is $(\dim\ltwentyfiveg+1)^m$. This holds (and is checked)
for $D_4$ ($29^m$), $G_2$ ($15^m$), $F_4$ ($53^m$), $E_6$ ($79^m$), $E_7$ ($134^m$), $E_8$ ($249^m$).
\end{theorem}
\noindent Files \leanpath{VOA/\allowbreak DeligneAudit/\allowbreak DeligneCharacterDemazureCount.lean},
\leanpath{DeligneCharacterE8Count.lean}.
\begin{leancode}
theorem demazureChar_total (A : List (List ℤ)) (θ : List ℤ) (hθ : θ.length = r)
    (word : List ℕ)
    (hα : ∀ ρ ∈ ρZero (r := r) A θ :: word.map (ρFin A), ρ.α ≠ 0)
    (hinv : ∀ ρ ∈ ρZero (r := r) A θ :: word.map (ρFin A),
      ρ.sLin (toR1 r (demazureChar A θ word 1)) = toR1 r (demazureChar A θ word 1)) (m : ℕ) :
    ((demazureChar A θ word m).map (·.2)).sum = (((demazureChar A θ word 1).map (·.2)).sum) ^ m

theorem e8_demazure_total (m : ℕ) :
    ((demazureChar cartanE8 thetaE8 wordE8 m).map (·.2)).sum = 249 ^ m := by
\end{leancode}
\begin{dictionary}
Namespace \lean{VOA.DeligneAudit.DeligneCharacter.DemCount}. \lean{toR1 r} maps a term list into the group ring
$\Z[\Z^r]$ at $q=1$; \lean{ρ.sLin} is the linear reflection of one operator. \lean{((…).map (·.2)).sum} is the sum of all
coefficients. The hypotheses \lean{hθ} (length of $\theta$) and \lean{hα} (every operator's root $\alpha\ne0$; this is
where the domain argument of the proof needs non-degeneracy) are listed in the statement above. Instances \lean{d4_demazure_total}, \lean{g2_demazure_total}, \lean{e6_demazure_total},
\lean{e7_demazure_total} have the same shape.
\end{dictionary}
\begin{proof}[Proof on paper]
At $q=1$ each operator satisfies $(1-e^{-\alpha})Df=f-e^{-\alpha}s^{\rm aff}(f)$ with $s^{\rm aff}(fg)=s(f)s^{\rm aff}(g)$; since
$\Z[\Z^r]$ is a domain, $D(fg)=fD(g)$ whenever $sf=f$. Iterate. The invariance of $\chi$ is a finite check.
\end{proof}
\begin{physics}
At $q=1$ the grading is forgotten, and the statement is the Fourier--Littelmann factorisation
$D(1,m\theta)\cong D(1,\theta)^{\otimes m}$ as $\ltwentyfiveg$-modules, at the level of dimensions. In the quasi-particle
picture (heuristic) it says that $m$ quasi-particles, each with internal states $\ltwentyfiveg\oplus\C$, are
independent once energies are ignored: the $q$-grading records only their interaction energies.
\end{physics}
\begin{status}
Standard axioms; unconditional; no literature input.
\end{status}

\begin{theorem}[Energies by a convexity invariant]\label{l25:thm:energy}
Let $\theta$ have $r$ entries and every letter of the word be $<r$. Let $F$ be a finite list of functionals stable under
the word's reflections and $s_\theta$, containing $-\theta^\vee$, with
$\varphi(\theta)\le2$ for all $\varphi\in F$. Then every term $z^\mu q^e$ of $C_m$ has $e\ge-m^2$.
\end{theorem}
\noindent File \leanpath{VOA/\allowbreak DeligneAudit/\allowbreak DeligneCharacterDemazureEnergy.lean}.
\begin{leancode}
structure OrbitData (r : ℕ) (A : List (List ℤ)) (θ : List ℤ) (word : List ℕ) (F : List (List ℤ)) :
    Prop where
  hθ : θ.length = r
  hw : ∀ i ∈ word, i < r
  hFin : ∀ i ∈ word, ∀ φ ∈ F, reflFunFin r A i φ ∈ F
  hFθ : ∀ φ ∈ F, reflFunTheta r A θ φ ∈ F
  hneg : negTheta r A θ ∈ F
  hθ2 : ∀ φ ∈ F, (List.zipWith (· * ·) φ θ).sum ≤ 2

theorem energy_of_orbit (hD : OrbitData r A θ word F) (m : ℕ) :
    ∀ t ∈ demazureChar A θ word m, 0 ≤ t.1.2 + (m : ℤ) * m
\end{leancode}
\begin{dictionary}
\lean{t.1.2} is the energy $e$; the conclusion is $e+m^2\ge0$. \lean{reflFunFin}, \lean{reflFunTheta} are $\varphi\circ s_i$,
$\varphi\circ s_\theta$; \lean{negTheta} is $-\theta^\vee$. All hypotheses are decidable and checked per type with $F$ the
$W$-orbit of $\theta^\vee$ (\lean{g2DemazureEnergy_holds}, \lean{e6DemazureEnergy}, \lean{e7DemazureEnergy},
\lean{e8DemazureEnergy}, \lean{f4DemazureEnergy}).
\end{dictionary}
\begin{proof}[Proof on paper]
Invariant: every term has $\varphi(\mu)\le2m$ for $\varphi\in F$. Finite $D_i$ move $\mu$ along the segment to $s_i\mu$
(linear interpolation preserves the bound) and keep $e$. $D_0$ sends $\mu\mapsto\mu+k\theta$ with
$0\le k\le1-\langle\mu,\theta^\vee\rangle\le2m-1$ (using $-\theta^\vee\in F$) and lowers $e$ by $k$; so
$-(m-1)^2-(2m-1)=-m^2$.
\end{proof}
\begin{physics}
In the character $C_m$ the variable $q$ records $L_0$-energy inside the level-one WZW module. The bound $e\ge-m^2$,
together with the prefactor $q^{m^2}$ in $C_m$, says that the Demazure grading of $D(1,m\theta)$ starts at degree $0$. So $\mathrm{Hilb}_qD(1,m\theta)$ is a genuine polynomial in $q$ with no negative powers.
This is what makes the clamped and unclamped tables comparable, and the $D_4$ positivity input partly provable.
\end{physics}
\begin{status}
Standard axioms; unconditional. Replaces the $D_4$-only integrality trick.
\end{status}

\begin{remark}[Unclamped fusion series; Gaussian transform]
With energies and totals, the unclamped series $\mathtt{hilbSeriesG}$ and $\mathtt{fusionFZ}=\mathtt{fusionOf}(\cdot)$
(\leanpath{DeligneCharacterFusionZ.lean} and per-type \leanpath{DeligneCharacter*FusionZ.lean},
\leanpath{F4FusionCountZ.lean}) have total $(\dim\ltwentyfiveg)^m$ (\lean{fusionFZ_count}), removing every count and
positivity input from the U5 route of Chapter~\ref{ch:L26}. \leanpath{GaussianTransform.lean} proves Ravinder's mixed
fusion recursion in closed form (\lean{VOA.DeligneGaussian.mixed_eq}, \lean{demazure_eq_gaussian}); both are used in
Chapter~\ref{ch:L26}.
\end{remark}

%======================================================================
\section{The large kernel-checked data, described collectively}\label{l25:sec:data}
%======================================================================

About half of this chapter's files are certificates. Each consists of literal data and theorems proved by
\lean{decide +kernel} (the kernel evaluates the check; no \lean{native_decide}, no \lean{implemented_by}). A referee
should check \emph{what} is asserted and that the bridge lemmas connect it to the definitions above; the literals
themselves need no reading.
\begin{center}\small
\begin{tabular}{p{0.32\textwidth}p{0.6\textwidth}}
\toprule
files & what is checked\\
\midrule
\leanpath{E8VacuumSteps}, \leanpath{E8VacuumStepCheck00--30}, \leanpath{E8VacuumStepAll} & the 3874 steps
$\lambda\,\mathrm{ad}_bR=\mu R'$ ($\mu\ne0$) of the breadth-first propagation search through the 3875 quadrics of
$I_2^{\rm ns}$ for $E_8$, described in the proof of Theorem~\ref{l25:thm:carrierdata} (giving \lean{sym_rel_all}). These
certify relation symbols, not characters\\
\leanpath{E6VacuumChecks}, \leanpath{E7LieZhu} & Casimir match and divisibility checks for $E_6$, $E_7$; the $E_7$ chain of 1539
relation symbols\\
\leanpath{DeligneCharacterE8Chain0--57}, \leanpath{E8Inv0--9}, \leanpath{E8Light} & the intermediate Demazure characters
$P_0,\dots,P_{57}$ of $C_1$ for $E_8$ pinned to literals, one operator per file, in a light-import copy of the algorithm
bridged by \lean{rfl}/induction; nodup and reflection-membership checks for invariance\\
\leanpath{E7Chain0--2}, \leanpath{E7Inv0--4}, \leanpath{E6Chain0}, \leanpath{E6Inv0--1}, \leanpath{F4FusionCountChain0--3},
\leanpath{F4FusionCountInv} & the same for $E_7$, $E_6$, $F_4$\\
\leanpath{D4FiniteCore}, \leanpath{D4FiniteCompute}, \leanpath{D4FiniteData}, \leanpath{D4FiniteCheckQ},
\leanpath{D4FiniteCheckSteps}, \leanpath{D4FiniteCheckSteps2}, \leanpath{D4FiniteIdentity}, \leanpath{D4FiniteStatementProof}
& the $D_4$ finite identity (P) in $\Z[q^\pm,s^\pm,x^\pm,z_1^\pm..z_4^\pm]$ (\lean{d4FiniteIdentity_holds}), via a
$(q,u,t,z)$ world with 4$\times$ fewer terms and one-step literal checkpoints. The elaborator \lean{d4_precompute}\texttt{\%} only
\emph{proposes} each literal; the kernel re-verifies it\\
\leanpath{DeligneCharacterD4Wp--Wp12} (and the other \leanpath{DeligneCharacterD4*} layers) & the theta-product chain of
Theorem~\ref{l25:thm:d4}; theory, not data, but long ($\wp'$ identity, half-period substitution, CT count)\\
\bottomrule
\end{tabular}
\end{center}

\subsection{The files of this chapter}\label{l25:sec:files}

The chapter covers the $239$ files assigned to it in the book's assignment table. About $120$ are theory (roughly
$26\,600$ lines), and the other $119$ are kernel-checked certificate and data files (roughly $3\,300$ lines).
\begin{center}\small
\begin{tabular}{p{0.45\textwidth}p{0.45\textwidth}}
\toprule
files & content\\
\midrule
\leanpath{VOA/\allowbreak DeligneAudit/\allowbreak DeligneCharacter*.lean} & interface, AK coefficients, Demazure algorithm, closed forms, $D_4$ evaluation, counts\\
\leanpath{DeligneMLDE*.lean} & the recurrence, the trace identity, Ramanujan, the heat equation, level two and three\\
\leanpath{DeligneVacuumRelation*.lean}, \leanpath{E6Vacuum*.lean}, \leanpath{E7LieVacuum.lean}, \leanpath{E7LieZhu.lean},
\leanpath{F4LieVacuum.lean}, \leanpath{G2LieVacuum.lean}, \leanpath{E8VacuumStep*.lean} & the vacuum quartic relation per type,
with propagation certificates\\
\leanpath{A2Character*.lean} & the $A_2$ character side and star identity\\
\leanpath{G2Star*.lean} & the $G_2$ target side and the named star reduction\\
\leanpath{D4Finite*.lean}, \leanpath{D4CharacterTop.lean} & the $D_4$ finite identity\\
\leanpath{F4FusionCount*.lean}, \leanpath{GaussianTransform.lean} & $F_4$ counts; Ravinder's mixed recursion in closed form\\
\leanpath{VOA/\allowbreak Commutative/\allowbreak JacobiTriple*.lean}, \leanpath{FourTriangular*.lean} & formal $q$-series identities\\
\leanpath{VOA/\allowbreak Modularity/\allowbreak ZhuCoordinate*.lean} & the square-bracket current calculus\\
\bottomrule
\end{tabular}
\end{center}
All paths without a directory are in \leanpath{VOA/\allowbreak DeligneAudit/}, except the \leanpath{FourTriangular*}
files, which are in \leanpath{VOA/\allowbreak Commutative/}.

\begin{warning}[title={Kernel-cost dead ends (recorded)}]
One \lean{decide} over the whole $D_4$ identity: 650 s and 41 GB (killed; a memory incident). \lean{![…]} table
lookups at high index: $O(\text{index})$ in the kernel (one 248-entry row 25 s / 8.7 GB vs 4 s / 2.4 GB with a literal
\lean{List} and \lean{getD}). Heavy imports in check files: use light-import copies bridged by \lean{rfl}.
\lean{decide} of \lean{e8DemazureHilb 1}: $>10$ min. The $A_2$ guard of $(\star)$ beyond $N=8$ did not finish.
\end{warning}

%======================================================================
\section{Lean defects and stale documentation}\label{l25:sec:defects}
%======================================================================

This section is the single place where the chapter's documentation defects are listed (Caveat~\ref{l25:cav:docs}).
Each is also entered in the project's defect log. No Lean statement of this chapter was found to be wrong, vacuous or
mis-hypothesised.

\paragraph{Missing statements and out-of-library checks (severity S3).}
\begin{itemize}
\item \textbf{Missing equality corollaries} (warning after Theorem~\ref{l25:thm:alltypes}). $\dim L_k(\ltwentyfiveg)_N=c_N$
  is a Lean theorem only for $A_2$ (\lean{A2.finrank_eq_mldeCoeff'}). For $G_2$, $D_4$, $F_4$, $E_6$, $E_7$, $E_8$ only
  \lean{TargetLowerBoundQ} ($\ge$) is exported. The equality is derivable by composing the per-type vacuum quartic
  relation (for $E_6$, $E_7$: \lean{vacuumQuarticRelation_of_carrierData} with the \lean{*_hrel} data),
  \lean{isAKSolution_of_vacuumQuarticRelation} and \lean{finrank_eq_mldeCoeff_of_isAKSolution}, but no such corollary is
  written. Composition checked in a scratch file for $G_2$, $F_4$, $D_4$ only.
\item \textbf{Missing $A_2$ character formula} (Remark~\ref{l25:rem:a2cons}). $\dim V_N=[q^N]\prod(1+q^r)^8$ for
  $V=L_{-3/2}(\mathfrak{sl}_3)$ is not stated. It follows in one line from \lean{A2.finrank_eq_mldeCoeff'} and
  \lean{A2.productCoeff_eq_mldeCoeff'}; only the inequality \lean{KacWakimotoLowerBound} is exported.
\item \textbf{Checks outside the library.} \lean{e7DemazureHilb_one} ($\mathrm{Hilb}\,D(1,\theta)=133+q$ for $E_7$, about
  50\,s) and \lean{A2_demazure_eq_GT_three} (agreement of the $A_2$ Gelfand--Tsetlin table with the affine algorithm, for
  $m=3$ and $N<10$ only) live in a separate heavy-check file in the project's notes, outside \leanpath{VOA/}. They are
  therefore not built by \lean{lake build}; they were run once by hand. For smaller $m$, the $A_2$ comparison was
  carried out outside Lean (the docstring of \leanpath{A2CharacterInterface.lean} cites an audit for $n\le7$).
\item \textbf{Untraced modelling.} No Lean statement connects \lean{D4K.d4SchurQ} to the $SU(2)$, $N_f=4$ Schur index,
  although the docstring calls it ``the Schur index''. The identification is a hand transcription.
\end{itemize}

\paragraph{Stale documentation (severity S4).}
\begin{itemize}
\item \leanpath{DeligneCharacterExceptional.lean}: the docstring describes \lean{XTargetLowerBound W} as a literature input.
  It is proved for the concrete $W$ (Theorem~\ref{l25:thm:alltypes}).
\item \leanpath{A2CharacterInterface.lean} and \leanpath{A2Character.lean}: the docstrings call the star identity and
  Kac--Wakimoto ``not formalized'' or ``literature input''. Both now follow (\lean{A2.starIdentity},
  \lean{A2.kacWakimotoLowerBound_of_quartic}).
\item \lean{A2ProductMLDE} (``checked numerically through $q^{40}$''), \lean{D4SchurEqE4} (``Named input''),
  \lean{LorenzA2} (``named''), \lean{RamanujanFormal} and \lean{RamanujanIdentities} (``named input''): all are proved.
\item \leanpath{D4CharacterTop.lean} calls \lean{D4DemazurePositivity} a literature input. Its \lean{energy} field is
  proved, and \lean{coeff} is available only from the open U5 (\lean{d4DemazurePositivity_of_U5Z}). The docstring of
  \leanpath{DeligneCharacterD4Schur.lean} calls positivity ``a consequence of the Demazure character formula'', a reading
  that no Lean statement backs.
\item \leanpath{DeligneCharacterDemazure.lean}: the title says ``simply-laced'', but the algorithm is used, correctly, for
  $G_2$ and $F_4$ (dictionary of Definition~\ref{l25:def:demazure}).
\item The name \lean{mldeCoeff} is used for two unrelated declarations (warning in Section~\ref{l25:sec:interface}).
\end{itemize}

%======================================================================
\section{Referee's checklist}\label{l25:sec:checklist}
%======================================================================

\begingroup\small
\begin{longtable}{p{0.37\textwidth}p{0.23\textwidth}p{0.19\textwidth}p{0.10\textwidth}}
\toprule
Lean declaration & Intended meaning & Hypotheses / status & Check\\
\midrule\endhead
\lean{CarrierCharacterLowerBound}, \lean{ordinaryCarrierCharacter_of_le} & lower half suffices & packages; uncond.\ impl. & Thm.~\ref{l25:thm:ofle}\\
\lean{demazureSumCoeff}, \lean{StarIdentityFor}, \lean{DemazureBoundFor} & Demazure sum and split & definitions & Def.~\ref{l25:def:dsum}\\
\lean{carrierCharacterLowerBound_of_demazureSum} & glue & \lean{hcar}, \lean{hchr} & Thm.~\ref{l25:thm:glue}\\
\lean{DeligneCharacter.mldeCoeff} & AK coefficients & definition & Def.~\ref{l25:def:mlde}\\
\lean{eq_mldeCoeff_of_isAKSolution} & uniqueness & unconditional & Thm.~\ref{l25:thm:unique}\\
\lean{torusOnePoint_Lsq_sq_vacuum} & $S(L[-2]^2\ltwentyfivevac)$ formula & unconditional & Thm.~\ref{l25:thm:tower}\\
\lean{VacuumQuarticRelation}, \lean{FormalMLDE} & trace identity / MLDE & definitions & Def.~\ref{l25:def:vqr}\\
\lean{kappa_eq_of_formalMLDE}, \lean{isAKSolution_of_vacuumQuarticRelation} & $\kappa$ forced; recurrence & $c$, $\ltwentyfivevac\ne0$, relation & Thm.~\ref{l25:thm:kappa}\\
\lean{sqCurrent_comm}, \lean{Lsq_sqCurrent_comm}, \lean{vSqZ_L_neg_one} & square current algebra & unconditional & Thm.~\ref{l25:thm:sqcomm}\\
\lean{vacuumQuarticRelation_of_symbols} & symbols $\Rightarrow$ relation & symbol vanishing & Thm.~\ref{l25:thm:symbols}\\
\lean{targetLowerBound_of_carrierData} & carrier data $\Rightarrow$ bound & \lean{hc}, \lean{hcas}, \lean{hrel} & Thm.~\ref{l25:thm:carrierdata}\\
\lean{A2.vacuumQuarticRelation}, \lean{A2.finrank_eq_mldeCoeff'} & $\dim L_{-3/2}(\mathfrak{sl}_3)_N=c_N$ & \textbf{unconditional} & Thm.~\ref{l25:thm:alltypes}\\
(none) & $\dim L_{-3/2}(\mathfrak{sl}_3)_N=[q^N]\prod(1+q^r)^8$ & derivable, \textbf{not stated} & Rem.~\ref{l25:rem:a2cons}\\
\lean{g2_targetLowerBound} & $\dim L_{-5/3}(\mathfrak g_2)_N\ge c_N$, $h=4$ & no hypotheses ($\ge$ only) & Thm.~\ref{l25:thm:alltypes}\\
\lean{targetLowerBoundQ4} & $\dim L_{-2}(\mathfrak{so}_8)_N\ge c_N$, $h=6$ & no hypotheses ($\ge$ only) & Thm.~\ref{l25:thm:alltypes}\\
\lean{f4_targetLowerBound} & $\dim L_{-5/2}(\mathfrak f_4)_N\ge c_N$, $h=9$ & no hypotheses ($\ge$ only) & Thm.~\ref{l25:thm:alltypes}\\
\lean{e6_targetLowerBound} & $\dim L_{-3}(\mathfrak e_6)_N\ge c_N$, $h=12$ & no hypotheses ($\ge$ only) & Thm.~\ref{l25:thm:alltypes}\\
\lean{e7_targetLowerBound} & $\dim L_{-4}(\mathfrak e_7)_N\ge c_N$, $h=18$ & no hypotheses ($\ge$ only) & Thm.~\ref{l25:thm:alltypes}\\
\lean{targetLowerBoundQ8} & $\dim L_{-6}(\mathfrak e_8)_N\ge c_N$, $h=30$ & no hypotheses ($\ge$ only) & Thm.~\ref{l25:thm:alltypes}\\
(none) & $\dim L_k(\ltwentyfiveg)_N=c_N$ for $G_2,D_4,F_4,E_6,E_7,E_8$ & derivable, \textbf{not stated} & Caveat~\ref{l25:cav:equality}\\
\lean{vacuumQuarticRelation_of_carrierData}, \lean{g2_/f4_vacuumQuarticRelation}, \lean{vacuumQuarticRelation4/8} & vacuum quartic relation & unconditional; none for $E_6$, $E_7$ & Thm.~\ref{l25:thm:alltypes}\\
\lean{jacobiTriple}, \lean{euler_eq_pentagonal}, \lean{quintuple_sum}, \lean{lambertS_eq} & classical $q$-identities & unconditional & Thms.~\ref{l25:thm:jtp}--\ref{l25:thm:legendre}\\
\lean{ramanujanIdentities} & Ramanujan & unconditional (Mathlib Sturm) & Thm.~\ref{l25:thm:ramanujan}\\
\lean{A2.a2ProductMLDE}, \lean{g2KWEqMLDE}, \lean{lorenzA2}, \lean{g2Level3Core} & level 2, 3 targets & unconditional & Thms.~\ref{l25:thm:a2prod}, \ref{l25:thm:g2target}\\
\lean{eq_mldeCoeff_of_HOde}, \lean{e6/e7/e8Target_eq_mldeCoeff} & closed forms & unconditional & Thms.~\ref{l25:thm:HOde}, \ref{l25:thm:closed}\\
\lean{demazureChar}, \lean{levelOneDemazureHilb} & Demazure characters & algorithm; U7 to interpret & Def.~\ref{l25:def:demazure}\\
\lean{A2.starIdentity} & $A_2$ star identity & unconditional ($q$-series) & Thm.~\ref{l25:thm:a2star}\\
\lean{d4SchurEqE4}, \lean{d4DemazureEqSchur_of_positivity}, \lean{d4DemazurePositivity_of_U5Z} & series \lean{d4SchurQ} $=$ MLDE; Demazure $=$ \lean{d4SchurQ} & uncond.; \lean{D4DemazurePositivity} (\lean{coeff} only from open U5) & Thm.~\ref{l25:thm:d4}\\
\lean{G2StarBound}, \lean{G2DemazureEqKW} & $D_N\le c_N$ for $G_2$; its Demazure half & \textbf{open, named}; conjectural; numerics through $q^{60}$ & Def.~\ref{l25:def:starbounds}\\
\lean{F4StarBound} & $D_N\le c_N$ for $F_4$ & \textbf{open, named}; conjectural; numerics through $q^{25}$ & Def.~\ref{l25:def:starbounds}\\
\lean{E6StarBound}, \lean{E7StarBound}, \lean{E8StarBound} & $D_N\le c_N$ for $E_6,E_7,E_8$ & \textbf{open, named}; conjectural; numerics through $q^{20},q^{16},q^{14}$ & Def.~\ref{l25:def:starbounds}\\
\lean{D4DemazurePositivity} (field \lean{coeff}) & $D_4$ Demazure coefficients $\ge0$ & \textbf{open} except via U5 & Thm.~\ref{l25:thm:d4}\\
\lean{demazureChar_total}, \lean{energy_of_orbit} & $q=1$ counts; energies & unconditional & Thms.~\ref{l25:thm:total}, \ref{l25:thm:energy}\\
\bottomrule
\end{longtable}
\endgroup

%======================================================================
\section{Notes and references}\label{l25:sec:notes}
%======================================================================

The second-order MLDE for the Deligne series is due to T.~Arakawa and K.~Kawasetsu, ``Quasi-lisse vertex algebras and
modular linear differential equations'', in \emph{Lie Groups, Geometry, and Representation Theory} (Kostant memorial
volume), Progr.\ Math.\ 326 (2018), arXiv:1610.05865, \S7; the physics prediction is C.~Beem and L.~Rastelli, ``Vertex
operator algebras, Higgs branches and modular differential equations'', JHEP 08 (2018) 114, arXiv:1707.07679. The
square-bracket modes are from Y.~Zhu, ``Modular invariance of characters of vertex operator algebras'', J.\ AMS 9
(1996). Second-order MLDEs: S.~Mathur, S.~Mukhi and A.~Sen, Phys.\ Lett.\ B 213 (1988). The Deligne series:
P.~Deligne, ``La s\'erie exceptionnelle de groupes de Lie'', C.\ R.\ Acad.\ Sci.\ 322 (1996). Admissible characters:
V.~Kac and M.~Wakimoto, arXiv:1612.07423. Demazure modules and fusion: B.~Ravinder, arXiv:1409.0274; E.~Feigin,
SIGMA 4 (2008), arXiv:0806.4851; G.~Fourier and P.~Littelmann, Adv.\ Math.\ 211 (2007); V.~Chari and B.~Ion,
Compos.\ Math.\ (2015); Y.~Sanderson, J.\ Algebraic Combin.\ 11 (2000); S.~Kumar, Invent.\ Math.\ 89 (1987);
O.~Mathieu, Ast\'erisque 159--160 (1988). Fermionic formulas: Hatayama--Kuniba--Okado--Takagi--Yamada, ``Remarks on
fermionic formula'', Contemp.\ Math.\ 248 (1999). Macdonald $t=0$ constant terms: J.~F.~van Diejen and E.~Emsiz.
The $q$-series identities are classical (Jacobi, Euler, Gauss, Legendre, Ramanujan; the quintuple product: Watson;
Lorenz's formula and the cubic identity: J.~M.~and P.~B.~Borwein and F.~Garvan, Trans.\ AMS 343 (1994)). Rogers'
Mehler formula for the Rogers--Szeg\H{o} polynomials and $q$-Vandermonde: G.~E.~Andrews, \emph{The Theory of Partitions}
(Addison--Wesley, 1976), Chapter~3 and its examples; Ramanujan's ``theory of signature 3'': B.~C.~Berndt,
S.~Bhargava and F.~G.~Garvan, Trans.\ AMS 347 (1995). Rogers--Szeg\H{o} polynomials and $q$-Whittaker functions at $t=0$:
I.~G.~Macdonald, \emph{Symmetric Functions and Hall Polynomials}, 2nd ed.\ (Oxford, 1995), Chapter~VI. Backward
pointers: Chapters~\ref{ch:L14}, \ref{ch:L15}, \ref{ch:L16}, \ref{ch:L23}, \ref{ch:L24}; forward: Chapter~\ref{ch:L26}
(fusion side, U5--U7) and Chapter~\ref{ch:L27} (the seven final theorems).

\endgroup
